\documentclass[journal,10pt,letterpaper]{IEEEtran}
\usepackage[utf8]{inputenc}
\usepackage[T1]{fontenc}
\usepackage[spanish,es-tabla,es-nodecimaldot,es-noshorthands]{babel}
\usepackage{amsmath,amssymb,amsthm,mathtools,bm}
\usepackage{graphicx,booktabs,array,tabularx,multirow}
\usepackage{enumitem,microtype,stfloats,placeins}
\usepackage[dvipsnames]{xcolor}
\usepackage{tikz}
\usetikzlibrary{arrows.meta,positioning,fit}
\usepackage{cite}
\usepackage{fancyhdr}
\usepackage[hidelinks]{hyperref}
\hypersetup{pdftitle={Metodología para la localización y evaluación de atractores ocultos en sistemas enteros y fraccionarios},pdfauthor={María Fernanda Moreno-López},pdfsubject={Balance armónico, puntos perpetuos, integración de Caputo y cuencas de atracción}}
\graphicspath{{figuras/}}
\setlength{\headheight}{13pt}
\setlength{\headsep}{12pt}
\setlength{\columnsep}{0.24in}
\pagestyle{fancy}
\fancyhf{}
\fancyhead[C]{\footnotesize Metodología de localización y evaluación de atractores ocultos}
\fancyfoot[C]{\footnotesize\thepage}
\renewcommand{\headrulewidth}{0.3pt}
\fancypagestyle{IEEEtitlepagestyle}{\fancyhf{}\fancyfoot[C]{\footnotesize\thepage}\renewcommand{\headrulewidth}{0pt}}
\raggedbottom
\emergencystretch=0.5em
\newcommand{\R}{\mathbb R}
\newcommand{\C}{\mathbb C}
\newcommand{\dd}{\,\mathrm d}
\newcommand{\ii}{\mathrm i}
\newcommand{\T}{\mathsf T}
\newcommand{\norm}[1]{\left\lVert#1\right\rVert}
\newcommand{\abs}[1]{\left|#1\right|}
\newcommand{\Caputo}[1]{{}^{\mathrm C}D_{t_0}^{#1}}
\newcolumntype{L}[1]{>{\raggedright\arraybackslash}p{#1}}
\newcolumntype{Y}{>{\raggedright\arraybackslash}X}
\newtheorem{proposition}{Proposición}[section]
\newtheorem{definition}[proposition]{Definición}
\newtheorem{remark}[proposition]{Observación}
\renewcommand{\proofname}{Demostración}
\renewcommand{\IEEEkeywordsname}{Palabras clave}
\setlist[enumerate]{leftmargin=*,itemsep=2pt,topsep=3pt}
\setlist[itemize]{leftmargin=*,itemsep=2pt,topsep=3pt}
\title{Metodología para la localización y evaluación de atractores ocultos en sistemas enteros y fraccionarios}
\author{María Fernanda Moreno-López}
\date{}
\begin{document}
\maketitle
\begin{abstract}
La localización de un atractor oculto requiere construir condiciones iniciales
fuera de las regiones accesibles mediante perturbaciones de los equilibrios
y distinguir la atracción de la recurrencia observada a tiempo finito.
Se desarrolla un procedimiento que combina análisis algebraico del campo,
balance armónico, puntos perpetuos, continuación e integración numérica.
Las ecuaciones de Caputo se formulan mediante integrales de Volterra, con
terminal inferior y memoria explícitos. Se justifican las condiciones de
estabilidad, los procedimientos de generación de semillas, la discretización
y los criterios de comparación dinámica. El análisis geométrico y topológico
precisa el papel de las separatrices, los retornos y los conjuntos invariantes.
Los resultados comprenden sistemas tipo Chua y no Chua de orden entero y
fraccionario, con trayectorias, indicadores numéricos y sondeos de cuenca.
Las demostraciones de existencia o acotación se distinguen de los resultados
de localización y de las clasificaciones sustentadas en muestras finitas.
\end{abstract}
\begin{IEEEkeywords}
Atractores ocultos, derivada de Caputo, función descriptiva, puntos perpetuos,
continuación numérica, cuencas de atracción.
\end{IEEEkeywords}
\section{Problema de localización y secuencia metodológica}
\label{bre:fund:problema}

Sea $F:\Omega\subset\R^n\to\R^n$ un campo definido en un abierto.
Se consideran el problema entero $\dot X=F(X)$ y el problema conmensurable
\begin{equation}
 \Caputo q X(t)=F(X(t)),\qquad X(t_0)=X_0,
 \quad 0<q<1.
 \label{bre:fund:pvi}
\end{equation}
La formulación conmensurable asigna el mismo orden a todas las componentes.
El caso $q=1$ se interpreta como una ecuación diferencial ordinaria (EDO).
La estructura tipo Chua y el orden de la derivada son características
independientes: una realización de Lur'e puede describir un sistema no Chua.

En orden entero, sea $\varphi_t$ el flujo hacia adelante, sobre los datos
cuyas soluciones existen. Para un compacto invariante y atractivo $A$, su cuenca es
\begin{equation}
 \mathcal B(A)=\{X_0:\operatorname{dist}(\varphi_t(X_0),A)\to0\}.
 \label{bre:fund:cuenca}
\end{equation}
Aquí se exige existencia global de las trayectorias incluidas en la cuenca.
La clasificación usual denomina oculto a un atractor cuya cuenca no
interseca vecindades suficientemente pequeñas de ningún equilibrio
\cite{LeonovKuznetsov2013,DudkowskiEtAl2016}. Con
$\mathcal E=\{e:F(e)=0\}$, la condición precisa es
\begin{equation}
 \forall e\in\mathcal E\ \exists r_e>0:
 B(e,r_e)\cap\mathcal B(A)=\varnothing.
 \label{bre:fund:hidden}
\end{equation}
El orden de los cuantificadores permite radios distintos para equilibrios
distintos. Un contacto a gran distancia de $e$ no contradice la existencia
de un radio menor excluido de la cuenca. Si $\mathcal E$ es vacío, la
condición es vacua; todavía debe establecerse la existencia de un atractor.

La metodología comprende seis operaciones encadenadas:
\begin{enumerate}
 \item Fijar campo, parámetros, dominio, orden, terminal inferior y
 convención de solución; resolver los equilibrios y su estabilidad local.
 \item Construir semillas mediante balance armónico, puntos perpetuos,
 continuación desde un problema resoluble o búsqueda acotada.
 \item Transportar las semillas al campo objetivo conservando la historia
 en los problemas fraccionarios.
 \item Integrar el campo autónomo bajo un mismo convenio de inicialización,
 separar el transitorio y comparar paso, horizonte y memoria.
 \item Identificar destinos mediante trayectorias, secciones, espectros,
 recurrencia e indicadores variacionales compatibles con el modelo.
 \item Estudiar la atracción y la separación respecto de los equilibrios;
 aplicar retornos, regiones invariantes o métodos topológicos cuando sus
 hipótesis estén verificadas.
\end{enumerate}
La resolución algebraica determina candidatos; la integración determina
su destino a tiempo finito. La demostración de atracción y la exclusión de
vecindades completas son problemas adicionales, con hipótesis propias.

\section{Fundamento analítico de la formulación fraccionaria}
\label{bre:fund:caputo}

\subsection{Integral fraccionaria y dominio del operador}

Para $a>0$ y $g\in L^\infty([t_0,T],\R^n)$, se define
\begin{equation}
 (I_{t_0}^a g)(t)=\frac1{\Gamma(a)}
 \int_{t_0}^{t}(t-s)^{a-1}g(s)\dd s.
 \label{bre:fund:integral}
\end{equation}
El núcleo es integrable en $s=t$ porque $a-1>-1$. En particular,
\begin{equation}
 \norm{I_{t_0}^a g(t)}\le
 \frac{(t-t_0)^a}{\Gamma(a+1)}\norm g_\infty,
 \label{bre:fund:integralbound}
\end{equation}
al integrar la potencia y usar $\Gamma(a+1)=a\Gamma(a)$.
Esta cota prueba la continuidad en el terminal inferior y la anulación
del incremento allí.

La composición $I^aI^b=I^{a+b}$ requiere el mismo terminal inferior.
Para justificarla, escriba la integral doble sobre
$t_0\le s\le u\le t$. Su integral absoluta es finita por la acotación de
$g$ y la integrabilidad de ambos núcleos; Fubini permite cambiar el orden.
La integral interior se calcula con $u=s+(t-s)v$:
\begin{align}
 &\int_s^t(t-u)^{a-1}(u-s)^{b-1}\dd u \nonumber\\
 &\quad=(t-s)^{a+b-1}\int_0^1(1-v)^{a-1}v^{b-1}\dd v\nonumber\\
 &\quad=(t-s)^{a+b-1}\frac{\Gamma(a)\Gamma(b)}{\Gamma(a+b)}.
 \label{bre:fund:beta}
\end{align}
Sustituir esta expresión en la integral exterior da la composición.
La identidad se extiende a $g\in L^1$: Tonelli aplicado al valor absoluto
proporciona, con $H=T-t_0$,
\[
 \norm{I^a g}_{L^1}\le\frac{H^a}{\Gamma(a+1)}\norm g_{L^1}.
\]
La integral doble es entonces absolutamente integrable en el intervalo y
Fubini demuestra la composición casi en todas partes. Esta extensión
permite aplicarla a la derivada ordinaria de una función absolutamente continua.
La misma sustitución demuestra, para $\mu>-1$,
\begin{equation}
 I_{t_0}^a(t-t_0)^\mu
 =\frac{\Gamma(\mu+1)}{\Gamma(\mu+a+1)}(t-t_0)^{\mu+a}.
 \label{bre:fund:power}
\end{equation}

Si $X$ es absolutamente continua, la derivada clásica de Caputo es
\begin{equation}
 \Caputo q X=I_{t_0}^{1-q}X'
 =\frac1{\Gamma(1-q)}\int_{t_0}^{t}(t-s)^{-q}X'(s)\dd s,
 \label{bre:fund:classical}
\end{equation}
donde la integral se entiende en los puntos en los cuales existe, al menos
casi en todas partes. Para no imponer una regularidad temporal que no se
haya demostrado, se utiliza también la definición
\begin{equation}
 \Caputo q X=\frac{\dd}{\dd t}I_{t_0}^{1-q}(X-X_0),
 \label{bre:fund:generalized}
\end{equation}
cuando $X$ es continua, $X(t_0)=X_0$ y el término derivado es
absolutamente continuo. Ambas definiciones coinciden sobre su dominio
común: $X-X_0=I^1X'$ y la composición de integrales da
$I^{1-q}(X-X_0)=I^1I^{1-q}X'$, cuya derivada es $I^{1-q}X'$.
Las soluciones se buscan en la clase $X_0+I_{t_0}^q(L^\infty)$ sobre
cada intervalo compacto de existencia \cite{Gomoyunov2018Convex}.

\subsection{Elección de Caputo y significado del dato inicial}
\label{bre:fund:initialchoice}

La elección del operador forma parte del modelo. Para $0<q<1$, Caputo
permite prescribir el valor ordinario $X(t_0)=X_0$ y anula las constantes.
La derivada de Riemann--Liouville, definida por
$D_{\rm RL}^qX=(\dd/\dd t)I^{1-q}X$, satisface sobre el dominio anterior
\begin{equation}
 D_{\rm RL}^qX=\Caputo qX+
 \frac{X_0}{\Gamma(1-q)}(t-t_0)^{-q},\quad t>t_0.
 \label{bre:fund:RLbridge}
\end{equation}
En efecto, se separa $X=(X-X_0)+X_0$ y se deriva
$I^{1-q}X_0=X_0(t-t_0)^{1-q}/\Gamma(2-q)$.
Por tanto, el mismo PVI puede escribirse con Riemann--Liouville si se
incorpora ese término singular de inicialización. Sustituir el operador
y conservar sin cambios el segundo miembro produce, en general, otro
problema. La conveniencia de Caputo aquí procede de los datos iniciales
puntuales y de la formulación causal elegida; no establece una superioridad
universal sobre otros operadores.

El valor inicial de $X$ no debe confundirse con el valor de su derivada
fraccionaria en el extremo. Para el campo constante $F(X)=v$, Volterra da
\begin{equation}
 X(t)=X_0+\frac{(t-t_0)^q}{\Gamma(q+1)}v.
 \label{bre:fund:startupconstant}
\end{equation}
Su derivada ordinaria es $v(t-t_0)^{q-1}/\Gamma(q)$, integrable pero
no acotada cuando $v\ne0$. Sustituirla en Caputo y aplicar la integral beta
da $\Caputo qX(t)=v$ para todo $t>t_0$. Así, el límite por la derecha
puede ser distinto de cero. Asignar cero a una integral de intervalo vacío
en $t=t_0$ es una convención de extremo; no autoriza a imponer
$F(X_0)=0$. La ecuación se exige para $t>t_0$ en el sentido de solución
declarado, y el dato inicial se obtiene por continuidad de $X$.

El operador de Liouville--Weyl, con historia desde $-\infty$, permite
analizar armónicos estacionarios. Su relación con el operador causal de
Caputo justifica el predictor armónico mediante un término explícito de
historia; no identifica ambos problemas iniciales. Las diferencias de
Grünwald--Letnikov pueden aproximar operadores compatibles con estas
formulaciones cuando se fijan la inicialización y la memoria apropiadas.
El nombre de la discretización, por sí solo, no determina el PVI que resuelve.

\subsection{Equivalencia con Volterra y existencia local}

Si $g(t)=F(X(t))$ es acotada, la ecuación integral correspondiente es
\begin{equation}
 X(t)=X_0+I_{t_0}^q[F(X)](t).
 \label{bre:fund:volterra}
\end{equation}
Aplicar $I^{1-q}$ a $X-X_0$ y usar \eqref{bre:fund:beta} produce
$I^1g$. Su derivada es $g$ casi en todas partes, de modo que
\eqref{bre:fund:generalized} satisface \eqref{bre:fund:pvi}.
Recíprocamente, integrar \eqref{bre:fund:generalized} da
$I^{1-q}(X-X_0)=I^1g$: la constante es cero por la continuidad en $t_0$.
Al restar $I^{1-q}I^qg=I^1g$, queda $I^{1-q}h=0$, con
$h=X-X_0-I^qg$. Aplicar $I^q$ da $I^1h=0$; al derivar se obtiene
$h=0$ casi en todas partes y, por continuidad, en todo el intervalo.
Así se demuestra la equivalencia, sin introducir $X'$ donde no esté
justificado \cite{Podlubny1999}.

Supóngase ahora $F$ localmente Lipschitz en $X$. Elija una bola cerrada
$\overline B(X_0,R)\subset\Omega$ y constantes $M,L$ tales que
$\norm{F(X)}\le M$ y
$\norm{F(X)-F(Y)}\le L\norm{X-Y}$ en ella. Sobre
\[
 \mathcal X=\{X\in C([t_0,t_0+\tau],\R^n):
               \norm{X-X_0}_\infty\le R\}
\]
se define $\mathcal TX=X_0+I^q[F(X)]$. La estimación anterior da
\begin{align}
 \norm{\mathcal TX-X_0}_\infty&\le\frac{M\tau^q}{\Gamma(q+1)},\nonumber\\
 \norm{\mathcal TX-\mathcal TY}_\infty
 &\le\frac{L\tau^q}{\Gamma(q+1)}\norm{X-Y}_\infty.
 \label{bre:fund:contraction}
\end{align}
Por tanto, escoger $M\tau^q/\Gamma(q+1)\le R$ y
$L\tau^q/\Gamma(q+1)<1$ hace de $\mathcal T$ una contracción de un
subconjunto cerrado de un espacio de Banach en sí mismo. Su único punto
fijo es la solución local. Para campos dependientes del tiempo se requieren
las mismas cotas uniformes, además de medibilidad temporal y continuidad
espacial; la formulación integral cubre campos por tramos durante la
continuación. Un campo PWL continuo es localmente Lipschitz, aunque su
jacobiano no exista en las superficies de conmutación.

La extensión a tiempos mayores conserva la integral acumulada. Si la
solución permanece en un compacto de $\Omega$ hasta un supuesto tiempo
máximo finito, $F(X)$ permanece acotada. La cota de continuidad de $I^q$
proporciona un límite al aproximarse a ese tiempo y la misma construcción
contractiva, con la historia fijada, extiende la solución. En consecuencia,
la pérdida de existencia global exige abandonar todos los compactos del
dominio. Esta alternativa se usa al demostrar acotación en Lorenz.

\subsection{Transformada de Laplace y respuesta lineal}

Trasladando $t_0$ a cero, supóngase $X-X_0=I^qg$ y
$e^{-\sigma t}g\in L^1([0,\infty))$ para algún $\sigma>0$.
Fubini con peso exponencial garantiza las transformadas de $X-X_0$,
$H=I^{1-q}(X-X_0)=I^1g$ y $H'=g$ en un semiplano derecho.
La transformada del núcleo de \eqref{bre:fund:integral}
es $s^{-a}$. Se obtiene de la integral gamma mediante $v=st$ para
$s>0$ real; la identidad se extiende analíticamente al semiplano de
convergencia con la rama principal. Fubini, ahora con peso exponencial,
convierte la convolución en producto. Como
$I^{1-q}(X-X_0)$ se anula inicialmente,
\begin{align}
 \mathcal L\{\Caputo qX\}(s)
 &=s\,s^{q-1}\left(\widehat X(s)-\frac{X_0}{s}\right)\nonumber\\
 &=s^q\widehat X(s)-s^{q-1}X_0.
 \label{bre:fund:laplace}
\end{align}
El término inicial no desaparece salvo cuando $X_0=0$.
Para $\Caputo qX=AX+Bu$, $y=CX$, se sigue
\begin{align}
 \widehat X(s)&=(s^qI-A)^{-1}s^{q-1}X_0\nonumber\\
 &\quad+(s^qI-A)^{-1}B\widehat u(s),\nonumber\\
 W_q(s)&=C(s^qI-A)^{-1}B.
 \label{bre:fund:transfer}
\end{align}
La transferencia describe la respuesta a la entrada con condición inicial
nula. En el balance armónico se usa
$(\ii k\omega)^q=(k\omega)^q e^{\ii q\pi/2}$ para $k>0$ y el conjugado
para $k<0$. El modo constante tiene multiplicador cero. Se fija
$s^q=\exp(q\operatorname{Log}s)$, con
$\operatorname{Arg}s\in(-\pi,\pi)$.

La expansión del resolvente para $\norm A<|s|^q$ produce
\[
 (s^qI-A)^{-1}s^{q-1}
 =\sum_{j=0}^\infty A^j s^{-qj-1}.
\]
La transformada inversa de cada término y la convergencia de la serie
de Mittag--Leffler dan
\begin{equation}
 X(t)=E_q(At^q)X_0,\quad
 E_{q,\beta}(Z)=\sum_{j=0}^\infty\frac{Z^j}{\Gamma(qj+\beta)},
 \label{bre:fund:ML}
\end{equation}
con $E_q=E_{q,1}$. Para una fuerza $g$, el mismo resolvente da
\begin{equation}
 X(t)=E_q(At^q)X_0+
 \int_0^t K_A(t-s)g(s)\dd s,
 \label{bre:fund:variation}
\end{equation}
donde $K_A(t)=t^{q-1}E_{q,q}(At^q)$. La convolución se interpreta sobre
intervalos finitos aun cuando no se haya demostrado una cota exponencial
global para el problema no lineal.
Para justificar la fórmula sin una inversión formal, supóngase
$g\in L^\infty([0,T])$ e itérese
$X=X_0+I^q(AX+g)$. La composición de integrales produce las series
$\sum_{j\ge0}A^jt^{qj}X_0/\Gamma(qj+1)$ y
$\sum_{j\ge0}A^jI^{q(j+1)}g$. Sus normas están dominadas,
respectivamente, por las series escalares con términos
$\norm A^jT^{qj}\norm{X_0}/\Gamma(qj+1)$ y
$\norm A^jT^{q(j+1)}\norm g_\infty/\Gamma(q(j+1)+1)$.
Ambas convergen; la convergencia uniforme permite integrar la suma,
verificar Volterra y obtener \eqref{bre:fund:variation} por unicidad.

\subsection{Equilibrios, linealización y estabilidad}

Un equilibrio satisface $F(e)=0$, porque Caputo anula las constantes.
Si $F\in C^1$ cerca de $e$, escriba
\begin{equation}
 F(e+\xi)=J_e\xi+R(\xi),\qquad J_e=DF(e).
 \label{bre:fund:linearization}
\end{equation}
El teorema fundamental del cálculo en una bola convexa proporciona
\begin{align}
 R(\xi)-R(\eta)
 &=\int_0^1[DF(e+\eta+s(\xi-\eta))-J_e]\nonumber\\
 &\hspace{24mm}\cdot(\xi-\eta)\dd s,\nonumber\\
 \norm{R(\xi)-R(\eta)}&\le\ell_r\norm{\xi-\eta},\nonumber\\
 \ell_r&=\sup_{\norm u\le r}\norm{DF(e+u)-J_e}\longrightarrow0.
 \label{bre:fund:smallremainder}
\end{align}
En particular, $R(\xi)=o(\norm\xi)$; una cota cuadrática requiere
$DF$ localmente Lipschitz, por ejemplo $F\in C^2$. En un campo PWL,
este argumento clásico se aplica a equilibrios interiores a una región.

El criterio de Matignon para $0<q<1$ establece estabilidad asintótica del
sistema lineal si y sólo si todos los autovalores son no nulos y
\begin{equation}
 |\operatorname{Arg}\lambda_j(J_e)|>q\pi/2.
 \label{bre:fund:matignon}
\end{equation}
La condición caracteriza el decaimiento de \eqref{bre:fund:ML}, incluidos
los bloques de Jordan \cite{Matignon1996}. Su aplicación no lineal requiere
el teorema de linealización y la pequeñez Lipschitz del resto
\cite{CongDoanSiegmundTuan2016}. El mecanismo puede verse en
\eqref{bre:fund:variation}: bajo el sector estricto existen
\begin{align*}
 M_A&=\sup_{t\ge0}\norm{E_q(J_et^q)}<\infty,\\
 C_A&=\int_0^\infty\norm{K_{J_e}(s)}\dd s<\infty.
\end{align*}
Estas estimaciones son consecuencias de las cotas sectoriales de
Mittag--Leffler. Cerca de cero el núcleo es del orden $t^{q-1}$,
integrable; en el infinito su decaimiento es del orden $t^{-q-1}$.
Se elige $r$ con $C_A\ell_r<1$ y
$M_A\norm{\xi_0}<(1-C_A\ell_r)r$.
La fórmula de variación impide que la solución alcance por primera vez
la frontera de la bola y da estabilidad. Para la atracción, se separa
la convolución en una parte de tiempo inicial fijo, que tiende a cero,
y una cola. Si $L=\limsup_{t\to\infty}\norm{\xi(t)}$, resulta
$L\le C_A\ell_rL$, de donde $L=0$.
Así se explicita la función de cada hipótesis del resultado no lineal.

El margen $m_q(e)=\min_j(|\operatorname{Arg}\lambda_j|-q\pi/2)$
cuantifica la distancia angular a la frontera. $m_q>0$ verifica el criterio
estricto; $m_q<0$ señala un modo lineal creciente. Si $m_q=0$ o existe un
autovalor nulo, el argumento anterior no decide la estabilidad. Para $q=1$
se recupera el criterio de Hurwitz $\operatorname{Re}\lambda_j<0$.
Un margen pequeño requiere controlar el error de los autovalores antes
de asignar su signo.

\subsection{Memoria y espacio de estados}

La evolución puntual de Caputo no satisface, en general, la ley de
semigrupo. El ejemplo escalar $\Caputo q x=\lambda x$ lo muestra:
\begin{align}
 &E_q\bigl(\lambda(2h)^q\bigr)-E_q(\lambda h^q)^2\nonumber\\
 &\quad=\frac{\lambda(2^q-2)}{\Gamma(1+q)}h^q+O(h^{2q}).
 \label{bre:fund:nosemigroup}
\end{align}
Para $\lambda\ne0$ y $0<q<1$ el coeficiente principal no es cero.
La composición de integrales de \eqref{bre:fund:beta} concierne a los
órdenes de los operadores; no autoriza a reiniciar la evolución temporal.

Si $t_0=0$, dividir Volterra en $[0,t]$ y $[t,t+\theta]$ produce
\begin{align}
 m_t(\theta)&=X_0+\frac1{\Gamma(q)}
  \int_0^t(t+\theta-s)^{q-1}F(X(s))\dd s,\nonumber\\
 X(t+\theta)&=m_t(\theta)+\frac1{\Gamma(q)}
  \int_0^\theta(\theta-u)^{q-1}F(X(t+u))\dd u.
 \label{bre:fund:profile}
\end{align}
El perfil $m_t$ determina la continuación; $m_t(0)=X(t)$ es sólo una
evaluación. La construcción de un semiflujo sobre perfiles exige
bien planteamiento global y continuidad en el espacio elegido
\cite{DoanKloeden2021Semidynamical}. No basta adoptar sin restricciones
todo $C([0,\infty),\R^n)$ con la topología de convergencia uniforme en
compactos. En efecto, para cualquier $c$, el término libre
\begin{equation}
 m_c(\theta)=c-\frac{\theta^q}{\Gamma(q+1)}F(c)
 \label{bre:fund:artificial}
\end{equation}
produce la solución constante $X\equiv c$. La continuación conserva
$m_c$, aunque $F(c)\ne0$. En el campo $F(x)=-x$, estos perfiles fijos
se aproximan al perfil cero cuando $c\to0$: el equilibrio estable de
reinicio no atrae una vecindad de todos los perfiles.

Para definir atracción funcional se especifica un espacio admisible
$\mathcal X_{\rm adm}$, positivamente invariante, con semiflujo continuo
$S_t$. Para un compacto estrictamente invariante $A$ y una familia
de reinicios $\iota(B)$, con $\iota(X_0)(\theta)=X_0$, la atracción uniforme
requiere $\iota(B)\subset\mathcal X_{\rm adm}$ y
significa
\begin{equation}
 \sup_{X_0\in B}\operatorname{dist}_{\mathcal X}
       (S_t\iota(X_0),A)\longrightarrow0.
 \label{bre:fund:resetattraction}
\end{equation}
Se debe indicar qué conjuntos $B$ se incluyen; un atractor local no tiene
que atraer todos los reinicios acotados. La ocultedad funcional excluye
vecindades de los perfiles de equilibrio en $\mathcal X_{\rm adm}$.
La ocultedad bajo reinicio excluye bolas de $X_0$ según
\eqref{bre:fund:hidden}, usando esa misma atracción. Los sondeos
numéricos de este estudio exploran reinicios finitos y destinos observables;
no construyen por sí solos un espacio funcional admisible ni verifican
la convergencia uniforme de \eqref{bre:fund:resetattraction}.

\subsection{Acotación sin aplicar una regla de la cadena ordinaria}

Para Caputo no se sustituye $\Caputo q[V(X)]$ por
$\nabla V(X)\cdot F(X)$. Si $V$ es convexa y $C^1$ con gradiente
localmente Lipschitz, y $X-X_0\in I^q(L^\infty)$, se dispone de
\begin{equation}
 \Caputo q[V(X(t))]\le\nabla V(X(t))\cdot\Caputo qX(t).
 \label{bre:fund:convex}
\end{equation}
La desigualdad se cumple casi en todas partes. El resultado también asegura la regularidad fraccionaria de
$V(X)-V(X_0)$ \cite{Gomoyunov2018Convex}. Para ver el signo en el caso
suave, fije $t$ y defina el resto convexo
$\Phi(s)=V(X(s))-V(X(t))-\nabla V(X(t))\cdot(X(s)-X(t))\ge0$.
Entonces $\Phi(t)=0$ y la diferencia entre ambos miembros, multiplicada
por $\Gamma(1-q)$, es
\begin{align}
 \int_0^t(t-s)^{-q}\Phi'(s)\dd s
 &=-t^{-q}\Phi(0)\nonumber\\
 &\quad-q\int_0^t(t-s)^{-q-1}\Phi(s)\dd s\le0.
 \label{bre:fund:convexproof}
\end{align}
La integración por partes se hace primero hasta $t-\varepsilon$.
Con trayectoria Lipschitz y gradiente Lipschitz,
$\Phi(s)=O((t-s)^2)$, por lo cual el término del extremo tiende a cero
y la integral converge. Para la clase $I^q(L^\infty)$, la extensión se
realiza por aproximación y compacidad; no se presupone la existencia de
$X'$ en toda esa clase. En la sección de resultados se aplica esta
desigualdad al Lorenz generalizado y se exhibe la cota obtenida.

Los modelos autónomos de Caputo con terminal inferior fijo y orden no
entero no admiten las órbitas periódicas no constantes propias de una EDO
bajo las hipótesis del teorema de no periodicidad
\cite{TavazoeiHaeri2009,AreaLosadaNieto2014}. Por ello, un ciclo del
problema armónico estacionario, una curva cerrada dibujada por una muestra
y una solución periódica exacta del PVI causal son objetos distintos.
La construcción armónica emplea el primero como predictor; la evaluación
numérica posterior se refiere siempre al PVI declarado.

\subsection{Qué permite trasladar un método de orden entero}
\label{bre:fund:transfermethods}

Las ecuaciones $F(e)=0$ y las realizaciones algebraicas del mismo campo no
dependen de $q$. Su reutilización es exacta. En cambio, sustituir
$\ii\omega$ por $(\ii\omega)^q$ sólo traslada la respuesta del operador
lineal estacionario: la selección de una trayectoria causal requiere
resolver nuevamente Volterra. La misma separación se aplica a los puntos
perpetuos, cuya ecuación algebraica se conserva, pero cambia su
interpretación como aceleración.

Una relación rigurosa con el caso entero puede establecerse a tiempo
finito. Sean $X_q$ y $X_1$ soluciones con el mismo dato inicial, sobre
$[0,T]$, contenidas en un compacto común en el cual $F$ esté acotado por
$M$ y sea Lipschitz con constante $L$. Ponga
\[
 k_q(t)=\frac{t^{q-1}}{\Gamma(q)},\qquad
 \varepsilon_q=\int_0^T|k_q(t)-1|\dd t.
\]
Para $q\in[q_*,1]$, $q_*>0$, el integrando está dominado por un múltiplo
de $t^{q_*-1}+1$, integrable en $[0,T]$. Como $k_q(t)\to1$ para
$t>0$ cuando $q\to1$, la convergencia dominada da
$\varepsilon_q\to0$. Con $e_q(t)=\norm{X_q(t)-X_1(t)}$, restar ambas
ecuaciones integrales produce
\begin{align}
 e_q(t)
 &\le M\varepsilon_q\nonumber\\
 &\quad+L\int_0^t k_q(t-s)e_q(s)\dd s,\nonumber\\
 \sup_{0\le t\le T}e_q(t)
 &\le M\varepsilon_q E_q(LT^q)\longrightarrow0.
 \label{bre:fund:qtoone}
\end{align}
La última desigualdad procede de iterar la integral positiva:
el término de orden $j$ es
$(Lt^q)^j/\Gamma(qj+1)$, y su suma es $E_q(Lt^q)$.
El factor de Mittag--Leffler permanece uniformemente acotado para
$q\in[q_*,1]$: ponga $B=L\max(T,T^{q_*})$. Para $j$ suficientemente
grande, la monotonía de $\Gamma$ en el semieje correspondiente permite
dominar su término por $B^j/\Gamma(q_*j+1)$; la serie dominante converge.
Los términos restantes son un número finito de funciones continuas de $q$.
La hipótesis de compacto común se verifica directamente o se utiliza la
estimación hasta el primer tiempo de salida. Esta prueba justifica
comparar trayectorias enteras y fraccionarias en una ventana controlada.
No proporciona una estimación uniforme para $T\to\infty$; por tanto, no
demuestra persistencia de un atractor, de su cuenca ni de su ocultedad al
variar $q$. Esos problemas requieren controlar por separado la dinámica
asintótica y el espacio de estados.

\begin{table*}[t]
\centering
\caption{Condiciones de adaptación de los métodos enteros al problema causal de Caputo.}
\label{bre:fund:adaptationtable}
\small
\begin{tabularx}{\textwidth}{@{}L{0.19\textwidth}L{0.30\textwidth}Y@{}}
\toprule
Operación & Resultado utilizable & Hipótesis o desarrollo adicional\\
\midrule
Equilibrios y realización & $F(e)=0$ e identidad de la realización del campo. &
Mismo campo, dominio y parámetros; tratar las regiones no suaves antes de derivar.\\
Linealización & Matriz $DF(e)$ y resto de Taylor. &
Regularidad $C^1$, pequeñez Lipschitz del resto y sector estricto de Matignon.\\
Balance armónico & Multiplicador de Fourier $(\ii k\omega)^q$ y coeficientes de la no linealidad. &
Problema auxiliar estacionario; rama compleja fijada, residuo físico y control de armónicos omitidos.\\
Continuación & Familia de campos con extremos algebraicamente comprobados. &
Una sola historia para la familia causal; un reinicio posterior define un nuevo PVI. La persistencia asintótica no se sigue de continuidad a tiempo finito.\\
Puntos perpetuos & Ecuación $DF\,F=0$ para generar semillas. &
Expansión de arranque en $t^q$; no identificarla con $X''=0$. Los eventos secuenciales dependen de la historia.\\
Retornos y cuencas & Cruces observados y destinos de reinicios. &
Operador funcional cuando se conserva memoria; hipótesis de retorno, atracción y exclusión de vecindades completas.\\
Lyapunov y compacidad & Desigualdad convexa y estimaciones integrales. &
Verificar la clase $I^q(L^\infty)$, continuación global y el espacio de perfiles; no usar una regla de la cadena ordinaria.\\
\bottomrule
\end{tabularx}
\end{table*}

Danca estudia Lorenz generalizado, Rabinovich--Fabrikant y Chua PWL con
Caputo, estabilidad sectorial y un predictor--corrector fraccionario
\cite{Danca2017}. Wu et al. construyen condiciones iniciales para Chua
con arctangente mediante una adaptación del procedimiento de
Leonov y emplean ADM en la simulación \cite{Wu2023ArctanChua}.
Ambos trabajos reconocen la dificultad de interpretar la periodicidad en
orden fraccionario; Wu et al. también discuten la pérdida de historia de
su aproximación por intervalos. Estas contribuciones suministran modelos,
parámetros y procedimientos de búsqueda. La validez de una operación
concreta se determina, además, por las hipótesis de la
Tabla~\ref{bre:fund:adaptationtable}. Reproducir una recurrencia publicada
y aproximar el PVI de Caputo con memoria completa son verificaciones
distintas, que se comparan mediante el operador integral correspondiente.

\section{Balance armónico y construcción de semillas}
\label{bre:bal:sec}

La localización comienza al separar una parte lineal resoluble y una
no linealidad cuyos coeficientes de Fourier puedan calcularse. El balance
determina frecuencia, amplitud y, cuando es necesaria, una componente
media. Estos datos se convierten en un estado inicial mediante el
resolvente del sistema; después se integra el campo original o una
continuación que termina exactamente en él. La construcción sigue el
método de función descriptiva y localización por continuación
\cite{GelbVanderVelde1968,LeonovKuznetsov2013,KuznetsovEtAl2017}.

\subsection{Realización escalar y convención de transferencia}
\label{bre:bal:lure}

Considérese un sistema autónomo con orden conmensurable
\(0<q\leq1\), terminal inferior fijo \(t_0\) y realización exacta
\begin{equation}
 \Caputo{q}X=PX+b\psi(\sigma),\qquad
 \sigma=r^{\T}X,
 \label{bre:bal:lure-eq}
\end{equation}
donde \(X\in\R^n\), \(P\in\R^{n\times n}\) y
\(b,r\in\R^n\). Para \(q=1\), el operador es la derivada ordinaria.
La palabra exacta exige comprobar
\(PX+b\psi(r^{\T}X)-F(X)\equiv0\) en el dominio del campo.
En particular, una aproximación local de \(F\) no constituye por sí
sola una realización de Lur'e del sistema completo.

Para Chua, escriba \(h(x)=\ell x+\psi(x)\) en
\begin{equation}
 \begin{aligned}
 \Caputo{q}x&=\alpha[y-x-h(x)],\\
 \Caputo{q}y&=x-y+z,\\
 \Caputo{q}z&=-\beta y-\gamma z.
 \end{aligned}
 \label{bre:bal:chua}
\end{equation}
La primera componente se expande como
\(-\alpha(1+\ell)x+\alpha y-\alpha\psi(x)\). Por ello,
\begin{equation}
 \begin{gathered}
 P=\begin{pmatrix}
 -\alpha(1+\ell)&\alpha&0\\
 1&-1&1\\0&-\beta&-\gamma
 \end{pmatrix},\\
 b=(-\alpha,0,0)^{\T},\qquad r=(1,0,0)^{\T}.
 \end{gathered}
 \label{bre:bal:chua-matrices}
\end{equation}
Las dos características consideradas corresponden a
\begin{equation}
 \begin{array}{c|c|c}
 &\ell&\psi(\sigma)\\\hline
 \text{saturación}&m_1&\Delta\operatorname{sat}(\sigma)\\
 \text{arctangente}&a_1&a_2\arctan(\rho\sigma)
 \end{array}
 \label{bre:bal:nonlinearities}
\end{equation}
con \(\Delta=m_0-m_1\), \(\rho>0\) y
\(\operatorname{sat}(s)=\max(-1,\min(1,s))\). El signo negativo de
\(b\) se conserva: cambiarlo sin modificar \(\psi\) cambia el campo.

Tras trasladar \(t_0\) a cero, la transformada de Laplace satisface
\begin{equation}
 (s^qI-P)\widehat X(s)=s^{q-1}X(t_0)+bU(s),
 \label{bre:bal:laplace}
\end{equation}
donde \(U\) es la transformada de \(\psi(\sigma)\). Se utiliza
exclusivamente la convención
\begin{equation}
 \begin{aligned}
 W_q(s)&=r^{\T}(s^qI-P)^{-1}b,\\
 \widehat\sigma(s)
 &=r^{\T}(s^qI-P)^{-1}s^{q-1}X(t_0)\\
 &\quad+W_q(s)U(s).
 \end{aligned}
 \label{bre:bal:transfer}
\end{equation}
La transferencia relaciona entrada y salida después de separar la respuesta
del dato inicial. Las expresiones con resolvente requieren
\(\det(s^qI-P)\ne0\).

\subsection{Resolvente, polinomio y condición de fase de Chua}
\label{bre:bal:resolvent}

Defina \(\lambda=s^q\), \(a_P=\alpha(1+\ell)\),
\(d_0=\beta+\gamma\) y \(d_1=1+\gamma\). Entonces
\begin{equation}
 \lambda I-P=\begin{pmatrix}
 \lambda+a_P&-\alpha&0\\
 -1&\lambda+1&-1\\
 0&\beta&\lambda+\gamma
 \end{pmatrix}.
 \label{bre:bal:resolvent-matrix}
\end{equation}
El menor superior izquierdo es
\begin{equation}
 Q(\lambda)=(\lambda+1)(\lambda+\gamma)+\beta
 =\lambda^2+d_1\lambda+d_0.
 \label{bre:bal:minor}
\end{equation}
La expansión por la primera fila da
\begin{equation}
 \begin{aligned}
 D_\ell(\lambda)
 &:={\det}(\lambda I-P)\\
 &=(\lambda+a_P)Q(\lambda)-\alpha(\lambda+\gamma)\\
 &=\lambda^3+(d_1+a_P)\lambda^2\\
 &\quad +(d_0+a_Pd_1-\alpha)\lambda+a_Pd_0-\alpha\gamma.
 \end{aligned}
 \label{bre:bal:characteristic}
\end{equation}
Al reemplazar la primera columna por \(b\), el determinante de Cramer
es \(-\alpha Q(\lambda)\). Como \(r^{\T}\) selecciona la primera
coordenada,
\begin{equation}
 W_q(s)=-\frac{\alpha Q(s^q)}{D_\ell(s^q)}.
 \label{bre:bal:chua-transfer}
\end{equation}
Este cálculo fija simultáneamente numerador, denominador y signo.

Para \(q=1\), la ecuación de fase se reduce a un polinomio real.
Ponga \(z=\omega^2\), con \(\omega>0\), y escriba
\begin{equation}
 \begin{aligned}
 Q(\ii\omega)&=d_0-z+\ii d_1\omega,\\
 D_\ell(\ii\omega)&=D_R+\ii D_I,\\
 D_R&=a_Pd_0-\alpha\gamma-(d_1+a_P)z,\\
 D_I&=\omega(d_0+a_Pd_1-\alpha-z).
 \end{aligned}
 \label{bre:bal:phase-parts}
\end{equation}
Multiplicar por el conjugado del denominador proporciona
\begin{equation}
 \Im W_1(\ii\omega)
 =-\alpha\frac{d_1\omega D_R-(d_0-z)D_I}{D_R^2+D_I^2}.
 \label{bre:bal:phase-imag}
\end{equation}
Fuera de los polos, dividir el numerador por \(\omega\) y agrupar
términos equivale a resolver
\begin{equation}
 \begin{aligned}
 0={}&z^2+(\alpha+d_1^2-2d_0)z\\
 &+d_0^2-\alpha d_0+\alpha\gamma d_1.
 \end{aligned}
 \label{bre:bal:phase-polynomial}
\end{equation}
Sólo se conservan raíces \(z>0\) que mantengan
\(D_\ell(\ii\sqrt z)\ne0\). Para \(q<1\) se evalúa
\eqref{bre:bal:chua-transfer} en
\begin{equation}
 \zeta=(\ii\omega)^q
 =\omega^q\left(\cos\frac{\pi q}{2}
       +\ii\sin\frac{\pi q}{2}\right),
 \label{bre:bal:spectral-frequency}
\end{equation}
con la rama principal. La frecuencia temporal es \(\omega\);
\(\zeta\) es la variable espectral que entra en el resolvente.

\subsection{Proyección de Fourier y balance centrado}
\label{bre:bal:df}

Para \(\sigma=A\cos\theta\), \(A>0\), el fasor fundamental de la
salida estática es
\begin{equation}
 Y_1(A)=\frac1\pi\int_0^{2\pi}
 \psi(A\cos\theta)e^{-\ii\theta}\,\dd\theta,
 \qquad N(A)=Y_1(A)/A.
 \label{bre:bal:fourier-df}
\end{equation}
La parte seno se anula porque \(\psi(A\cos\theta)\) es par en
\(\theta\). Si además \(\psi\) es impar, su media es cero y, mediante
un desplazamiento de fase y simetría entre semiperiodos,
\begin{equation}
 N(A)=\frac{2}{\pi A}\int_0^\pi
 \psi(A\sin\theta)\sin\theta\,\dd\theta\in\R.
 \label{bre:bal:df-sine}
\end{equation}
La integral pertenece a la característica estática: no contiene \(q\).
El orden fraccionario entra a través del bloque dinámico \(W_q\).
Si la no linealidad posee histéresis o estado interno, debe calcularse el
fasor de su respuesta periódica; en ese caso \(N\) puede ser complejo
y depender de la frecuencia y de la historia interna.

Al retener sólo el fundamental, la salida del bloque lineal debe
reproducir la entrada \(\sigma\). Así,
\(A=W_q(\ii\omega)Y_1\), y la ecuación de cierre es
\begin{equation}
 1-W_q(\ii\omega)N(A)=0.
 \label{bre:bal:physical-closure}
\end{equation}
Para \(N\) real, una solución finita no nula exige
\begin{equation}
 \Im W_q(\ii\omega)=0,\qquad
 N(A)=k:=\frac1{\Re W_q(\ii\omega)}.
 \label{bre:bal:phase-gain}
\end{equation}
Después de la condición de fase debe comprobarse que \(k\) pertenece al
rango real de \(N\). Una intersección de fase cuya ganancia no puede
alcanzar la no linealidad no produce amplitud admisible.

\subsubsection{Saturación}
Para \(\psi(s)=\Delta\operatorname{sat}(s)\) y \(0<A\leq1\),
\(\psi(A\sin\theta)=\Delta A\sin\theta\); como
\(\int_0^\pi\sin^2\theta\,\dd\theta=\pi/2\), resulta
\(N(A)=\Delta\). Si \(A>1\), sea
\(\theta_0=\arcsin(1/A)\). Los tramos no saturados son
\([0,\theta_0]\) y \([\pi-\theta_0,\pi]\); en el tramo central la
saturación vale uno. Por ello,
\begin{equation}
 \begin{aligned}
 N(A)&=\frac{2\Delta}{\pi A}
 \left[2A\int_0^{\theta_0}\sin^2\theta\,\dd\theta
       +2\cos\theta_0\right]\\
 &=\frac{2\Delta}{\pi A}
 \left[A\theta_0-\frac A2\sin(2\theta_0)+2\cos\theta_0\right]\\
 &=\frac{2\Delta}{\pi}
 \left[\arcsin\frac1A+\frac{\sqrt{A^2-1}}{A^2}\right].
 \end{aligned}
 \label{bre:bal:saturation-df}
\end{equation}
El último paso usa \(A\sin\theta_0=1\). Para \(\Delta>0\),
\begin{equation}
 N'(A)=-\frac{4\Delta\sqrt{A^2-1}}{\pi A^3}<0
 \quad(A>1).
 \label{bre:bal:saturation-derivative}
\end{equation}
Además, \(N(1)=\Delta\) y \(N(A)\to0\) al crecer \(A\).
Existe una única amplitud saturada para \(0<k<\Delta\). Si
\(k=\Delta\), toda amplitud \(0<A\leq1\) satisface el balance lineal
y éste no selecciona una amplitud aislada.

\subsubsection{Arctangente}
Sea \(\psi(s)=a_2\arctan(\rho s)\), \(\xi=\rho A\), y defina
\(I(\xi)=\int_0^\pi\arctan(\xi\sin\theta)\sin\theta\,\dd\theta\).
La diferenciación bajo la integral es válida por continuidad y acotación
del integrando y de su derivada en intervalos compactos de \(\xi\).
Para \(\xi>0\),
\begin{equation}
 \begin{aligned}
 I'(\xi)&=\int_0^\pi
 \frac{\sin^2\theta}{1+\xi^2\sin^2\theta}\,\dd\theta\\
 &=\frac1{\xi^2}\left[\pi-
 \int_0^\pi\frac{\dd\theta}{1+\xi^2\sin^2\theta}\right].
 \end{aligned}
 \label{bre:bal:arctan-integral}
\end{equation}
La integral restante se reduce por simetría a dos veces el intervalo
\([0,\pi/2]\). Con \(u=\tan\theta\), se transforma en
\(2\int_0^\infty[1+(1+\xi^2)u^2]^{-1}\,\dd u
=\pi/\sqrt{1+\xi^2}\). La función
\(\pi(\sqrt{1+\xi^2}-1)/\xi\) tiene la derivada obtenida y límite
cero en \(\xi=0\), igual a \(I(0)\). Se concluye
\begin{equation}
 \begin{aligned}
 N(A)&=\frac{2a_2}{\pi A}I(\rho A)\\
 &=\frac{2a_2}{\rho A^2}
       \left(\sqrt{1+\rho^2A^2}-1\right)\\
 &=\frac{2a_2\rho}{1+\sqrt{1+\rho^2A^2}}.
 \end{aligned}
 \label{bre:bal:arctan-df}
\end{equation}
La última forma evita la resta de números próximos cuando \(A\) es
pequeño. Si \(a_2\ne0\), la ecuación \(N(A)=k\) se despeja como
\begin{equation}
 A^2=\frac{4a_2(a_2\rho-k)}{k^2\rho}.
 \label{bre:bal:arctan-amplitude}
\end{equation}
Antes de elevar al cuadrado se tiene
\(\sqrt{1+\rho^2A^2}=2a_2\rho/k-1\). Por tanto, para \(A>0\),
la condición completa es \(0<k/(a_2\rho)<1\); impone el signo
correcto y evita aceptar raíces introducidas por el cuadrado.

\subsection{Balance de media y fundamental sesgado}
\label{bre:bal:bias}

Para localizar una oscilación alrededor de una media distinta de cero,
tome \(\sigma=c+A\cos\theta\). Defina los coeficientes físicos
\begin{equation}
 \begin{aligned}
 \Psi_0(c,A)&=\frac1{2\pi}\int_0^{2\pi}
       \psi(c+A\cos\theta)\,\dd\theta,\\
 \Psi_1(c,A)&=\frac1\pi\int_0^{2\pi}
       \psi(c+A\cos\theta)\cos\theta\,\dd\theta,\\
 N_1(c,A)&=\Psi_1(c,A)/A.
 \end{aligned}
 \label{bre:bal:bias-fourier}
\end{equation}
La media del bloque dinámico satisface \(P\bar X+b\Psi_0=0\).
Si \(P\) es invertible,
\begin{equation}
 \bar X=-P^{-1}b\Psi_0,\qquad
 c=r^{\T}\bar X=W_q(0)\Psi_0.
 \label{bre:bal:dc}
\end{equation}
Para Chua,
\begin{equation}
 W_q(0)=-\frac{\beta+\gamma}
 {(1+\ell)(\beta+\gamma)-\gamma}.
 \label{bre:bal:chua-dc}
\end{equation}
Si \(P\) es singular, se resuelven directamente las ecuaciones de
media, sin usar \(P^{-1}\) ni una ganancia continua infinita.

La condición armónica sigue siendo
\begin{equation}
 1-W_q(\ii\omega)N_1(c,A)=0.
 \label{bre:bal:bias-closure}
\end{equation}
La ecuación de media y las partes real e imaginaria de
\eqref{bre:bal:bias-closure} son tres ecuaciones reales para
\((c,A,\omega)\). Mantener la media evita tratar el desplazamiento
del ciclo como una coordenada libre sin ecuación.

Para la saturación, las integrales se evalúan por tramos incluso cuando
la entrada queda totalmente saturada. Con
\([u]_{-1}^{1}=\max(-1,\min(1,u))\), sean
\begin{equation}
 \theta_+=\arccos\left[\frac{1-c}{A}\right]_{-1}^{1},\qquad
 \theta_-=\arccos\left[\frac{-1-c}{A}\right]_{-1}^{1}.
 \label{bre:bal:bias-crossings}
\end{equation}
Se cumple \(0\leq\theta_+\leq\theta_-\leq\pi\). En
\([0,\pi]\), la salida normalizada vale sucesivamente
\(1\), \(c+A\cos\theta\) y \(-1\), con cambios en esos
ángulos. Integrar cada tramo da
\begin{equation}
 \begin{aligned}
 \Psi_0=\frac\Delta\pi\bigl[&\theta_++\theta_--\pi
       +c(\theta_--\theta_+)\\
       &+A(\sin\theta_--\sin\theta_+)\bigr].
 \end{aligned}
 \label{bre:bal:bias-sat-dc}
\end{equation}
Para el fundamental se usan las primitivas
\(\int\cos\theta\,\dd\theta=\sin\theta\) y
\(\int\cos^2\theta\,\dd\theta=\theta/2+\sin(2\theta)/4\).
La suma de los tres tramos es
\begin{equation}
 \begin{aligned}
 \Psi_1=\frac{2\Delta}\pi\bigl[&\sin\theta_++\sin\theta_-
       +c(\sin\theta_--\sin\theta_+)\\
       &+\tfrac A2(\theta_--\theta_+)\\
       &+\tfrac A4(\sin2\theta_--\sin2\theta_+)\bigr].
 \end{aligned}
 \label{bre:bal:bias-sat-first}
\end{equation}
Los ángulos recortados hacen desaparecer automáticamente un tramo de
longitud cero. En \(c=0\), la media se anula y
\(\Psi_1/A\) recupera \eqref{bre:bal:saturation-df}.

\subsection{Modo normalizado, semilla y equivalencias}
\label{bre:bal:modal}

Fijada una raíz de balance, sea \(k=N(A)\) o \(k=N_1(c,A)\), y
\(P_k=P+bkr^{\T}\). Si \(\zeta I-P\) es invertible, el lema del
determinante matricial proporciona
\begin{equation}
 \begin{aligned}
 \det(\zeta I-P_k)
 &=\det(\zeta I-P)\\
 &\quad\times\left[1-kr^{\T}(\zeta I-P)^{-1}b\right].
 \end{aligned}
 \label{bre:bal:det-lemma}
\end{equation}
En \(\zeta=(\ii\omega)^q\), el corchete es cero por el cierre.
Un modo normalizado se construye directamente, sin una elección
arbitraria de escala, mediante
\begin{equation}
 v=k(\zeta I-P)^{-1}b.
 \label{bre:bal:mode}
\end{equation}
En efecto, \(r^{\T}v=kW_q(\ii\omega)=1\) y
\((\zeta I-P_k)v=bk-bkr^{\T}v=0\).
Para Chua, las dos últimas filas de esta ecuación dan
\begin{equation}
 v=\left(1,\frac{\zeta+\gamma}{Q(\zeta)},
               -\frac\beta{Q(\zeta)}\right)^{\T}.
 \label{bre:bal:chua-mode}
\end{equation}
La igualdad presupone \(Q(\zeta)\ne0\), condición que se comprueba
junto con la existencia del resolvente.

En la rama sesgada, \(\Psi_1=kA\) y
\(X_1=(\zeta I-P)^{-1}b\Psi_1=Av\). Se obtiene la familia de semillas
\begin{equation}
 X_{\mathrm{sem}}(\varphi)=\bar X+Re(Av e^{\ii\varphi}),
 \qquad 0\leq\varphi<2\pi,
 \label{bre:bal:seed}
\end{equation}
con \(\bar X=0\) en la rama centrada. La salida es exactamente
\(r^{\T}X_{\mathrm{sem}}=c+A\cos\varphi\). El par
\(\zeta,\overline\zeta\) está en la frontera angular
\(|\arg\zeta|=q\pi/2\) del auxiliar. Para usarlo como predictor de
una oscilación dominante se comprueba que los demás modos satisfagan el
criterio angular estricto de estabilidad; el balance por sí solo no
establece esa propiedad \cite{Matignon1996}.

Una similitud real \(X=SY\), con \(S\) invertible, transforma
\(P,b,r^{\T}\) en \(S^{-1}PS,S^{-1}b,r^{\T}S\).
Como
\((s^qI-S^{-1}PS)^{-1}=S^{-1}(s^qI-P)^{-1}S\),
la transferencia no cambia y la semilla se transforma por
\(Y_{\mathrm{sem}}=S^{-1}X_{\mathrm{sem}}\).
Si se añade un desplazamiento \(X=SY+d\), deben conservarse también
el término constante \(S^{-1}Pd\) y el argumento
\(r^{\T}SY+r^{\T}d\); si \(d\) es equilibrio, se cancelan usando
\(Pd+b\psi(r^{\T}d)=0\) y la no linealidad incremental
\(\psi(r^{\T}SY+r^{\T}d)-\psi(r^{\T}d)\).

Otra equivalencia consiste en redistribuir una ganancia \(a>0\):
\(\widetilde b=ab\), \(\widetilde\psi=\psi/a\).
Entonces \(\widetilde W_q=aW_q\),
\(\widetilde N=N/a\), y
\(\widetilde W_q\widetilde N=W_qN\). El campo, el cierre y el modo
\(\widetilde k(\zeta I-P)^{-1}\widetilde b=v\) permanecen iguales.
Esta comprobación permite comparar realizaciones con diferentes escalas
sin confundir una ganancia de representación con un cambio del modelo.

\subsection{Del armónico estacionario al PVI causal}
\label{bre:bal:weyl-caputo}

El balance usa la acción del operador estacionario de Liouville--Weyl
sobre los armónicos:
\(D_{\mathrm W}^{q}e^{\ii\omega t}=(\ii\omega)^qe^{\ii\omega t}\).
Para \(0<q<1\), esta identidad se interpreta por Fourier o mediante
regularización de Abel de la integral sobre toda la historia. En efecto,
para \(\varepsilon>0\),
\begin{equation}
 \int_0^\infty u^{-q}e^{-(\varepsilon+\ii\omega)u}\,\dd u
 =\Gamma(1-q)(\varepsilon+\ii\omega)^{q-1}.
 \label{bre:bal:abel}
\end{equation}
Multiplicar por \(\ii\omega/\Gamma(1-q)\) y tomar
\(\varepsilon\downarrow0\) produce la rama de
\eqref{bre:bal:spectral-frequency}.

En el cálculo siguiente se fija \(t_0=0\). La derivada de Caputo
retiene sólo el intervalo \([0,t]\).
Para el mismo armónico suave,
\begin{equation}
 \Caputo{q}e^{\ii\omega t}
 =\frac{\ii\omega e^{\ii\omega t}}{\Gamma(1-q)}
   \int_0^t u^{-q}e^{-\ii\omega u}\,\dd u.
 \label{bre:bal:caputo-harmonic}
\end{equation}
Restar la integral sobre \([0,\infty)\) revela el defecto inicial:
\begin{equation}
 \begin{aligned}
 \Caputo{q}e^{\ii\omega t}-\zeta e^{\ii\omega t}
 ={}&-\frac{\ii\omega e^{\ii\omega t}}{\Gamma(1-q)}\\
 &\times\int_t^\infty u^{-q}e^{-\ii\omega u}\,\dd u.
 \end{aligned}
 \label{bre:bal:initial-defect}
\end{equation}
La integral impropia converge para \(\omega>0\) por el criterio de
Dirichlet. Cerca de cero,
\(\Caputo{q}e^{\ii\omega t}
=\ii\omega t^{1-q}/\Gamma(2-q)+O(t^{2-q})\), mientras
\(\zeta e^{\ii\omega t}\to\zeta\ne0\).
Por consiguiente, incluso el armónico del auxiliar lineal exige una
respuesta de arranque en el PVI de Caputo. Para \(q=1\), la identidad
ordinaria no presenta este defecto.

En el sistema no lineal se añade un segundo defecto: los armónicos
descartados de \(\psi\). El balance de primer armónico no los hace
cero. La semilla de \eqref{bre:bal:seed} localiza un estado, pero la
trayectoria causal debe obtenerse resolviendo el PVI. Los resultados
de no periodicidad para sistemas autónomos de Caputo con terminal
inferior finito impiden interpretar sin más el cierre estacionario como
una órbita periódica exacta del PVI
\cite{TavazoeiHaeri2009,AreaLosadaNieto2014}.

La distinción permite precisar el traslado de los métodos publicados.
Danca utiliza el PVI de Caputo y ABM, y describe la localización mediante
linealización armónica y continuación, junto con exploraciones numéricas
de las vecindades de equilibrios \cite{Danca2017}. Wu et al. aplican a
Chua con arctangente el algoritmo de localización de Leonov y construyen
un auxiliar mediante transferencias racionales en una variable compleja
\(p\), con un bloque de denominador \(p^2+\omega_0^2\)
\cite{Wu2023ArctanChua}. La reducción algebraica de una matriz por
similitud sigue siendo válida en Caputo; interpretar sus polos como una
frecuencia temporal requiere además la relación \(p=(\ii\omega)^q\).
En particular, conservar autovalores \(\pm\ii\omega_0\) con \(q<1\)
los sitúa estrictamente en el sector estable de Matignon, no en su
frontera. Es posible conservar una semilla obtenida en orden entero y
probarla en Caputo, pero eso constituye un PVI distinto del predictor
que recalcula la fase con \(s^q\).

Wu et al. reconocen expresamente tanto la no periodicidad exacta como
la pérdida de parte de la memoria histórica al emplear su aproximación
por descomposición de Adomian en subintervalos. También distinguen los
contornos numéricamente periódicos de la periodicidad matemática
\cite{Wu2023ArctanChua}. El problema pendiente al transferir una
prueba entera es demostrar persistencia y atracción con la memoria
causal declarada; no se resuelve sólo reproduciendo un contorno o
reutilizando una fórmula de estado inicial. La independencia de \(q\)
en \eqref{bre:bal:df-sine} sí es legítima para una característica
estática: elevar \(N\) a \(q\) no se deduce ni de Fourier ni de
la transformada de Caputo.

\subsection{Continuación que conserva la media y el campo final}
\label{bre:bal:continuation}

La descomposición
\(\psi(\sigma)=\Psi_0+k(\sigma-c)
+[\psi(\sigma)-\Psi_0-k(\sigma-c)]\) define el auxiliar afín
\(P_k=P+bkr^{\T}\), \(d_k=b(\Psi_0-kc)\), y la familia
\begin{equation}
 \begin{aligned}
 F_\varepsilon(X)&=P_kX+d_k+\varepsilon b\delta\psi(r^{\T}X),\\
 \delta\psi(\sigma)&=\psi(\sigma)-\Psi_0-k(\sigma-c),
 \qquad 0\leq\varepsilon\leq1.
 \end{aligned}
 \label{bre:bal:homotopy}
\end{equation}
En \(\varepsilon=0\),
\(F_0(\bar X)=P\bar X+bk c+b(\Psi_0-kc)=0\), y
\(F_0(\bar X+\xi)=P_k\xi\). En \(\varepsilon=1\), los términos
\(bk r^{\T}X\), \(b\Psi_0\) y \(bkc\) se cancelan con los del
corchete y queda exactamente \(F_1(X)=PX+b\psi(r^{\T}X)\).
Una homotopía que omita \(d_k\) pierde esta correspondencia cuando
\(c\ne0\).

Se inicia en \(X_{\mathrm{sem}}\), se hace crecer
\(\varepsilon(t)\) por etapas y se evalúa una única integral causal:
\begin{equation}
 \begin{aligned}
 X(t)=X_{\mathrm{sem}}+\frac1{\Gamma(q)}
 \int_{t_0}^t(t-s)^{q-1}F_{\varepsilon(s)}(X(s))\,\dd s.
 \end{aligned}
 \label{bre:bal:causal-continuation}
\end{equation}
Cada valor pasado del campo conserva el parámetro que tuvo en ese
instante. Al cambiar de etapa se prolonga esa historia. La continuidad
numérica de una rama sólo permite seguir candidatos: no demuestra que
exista un atractor en todos los parámetros intermedios. Para un estudio
de cuencas con reinicios se realiza después un PVI autónomo del campo
final con el mismo terminal inferior y convenio de memoria usados en
los demás lanzamientos.

\subsubsection{Recuperación del método entero}
Las matrices \(P,b,r\), las proyecciones \(\Psi_0,\Psi_1\) y
la identidad algebraica de la homotopía no dependen del orden. Para
una frecuencia positiva fija y fuera de los polos,
\begin{equation}
 \begin{aligned}
 (\ii\omega)^q&\longrightarrow\ii\omega,\\
 W_q(\ii\omega)&\longrightarrow r^{\T}(\ii\omega I-P)^{-1}b
 \qquad(q\to1^-).
 \end{aligned}
 \label{bre:bal:integer-limit}
\end{equation}
La continuidad del inverso matricial se obtiene, por ejemplo, de su
expresión mediante la adjunta y el determinante no nulo. Por tanto,
el cierre converge a \(1-W_1N=0\); para Chua, su ecuación de fase
es precisamente \eqref{bre:bal:phase-polynomial}.

La continuidad de las \emph{raíces} requiere una hipótesis adicional.
Para la rama sesgada, forme el residual real
\begin{equation}
 \mathcal F(c,A,\omega;q)=\begin{pmatrix}
 c-W_q(0)\Psi_0\\
 \Re(1-W_q(\ii\omega)N_1)\\
 \Im(1-W_q(\ii\omega)N_1)
 \end{pmatrix}.
 \label{bre:bal:branch-residual}
\end{equation}
Si \(\mathcal F\) es \(C^1\) cerca de una raíz con \(q=1\),
\(A>0\), \(\omega>0\), y su Jacobiano respecto de
\((c,A,\omega)\) es invertible, el teorema de la función implícita
produce una rama local \((c(q),A(q),\omega(q))\) que converge a
esa raíz. Para el balance centrado se aplica el mismo argumento a
las dos ecuaciones reales en \((A,\omega)\).

En esa rama, \(k(q)=N_1(c(q),A(q))\), \(P_{k(q)}\), \(v(q)\),
\(\bar X(q)\) y la semilla de \eqref{bre:bal:seed} convergen a sus
expresiones enteras porque se construyen mediante funciones continuas
y resolventes no singulares. En \(q=1\),
\(\bar X+\Re(Av e^{\ii\omega t})\) resuelve exactamente el auxiliar
afín: al derivarla, \(\ii\omega v=P_kv\) y \(F_0(\bar X)=0\)
dan la igualdad componente a componente. El campo final sigue
requiriendo la continuación y su validación dinámica. Si falla la
invertibilidad del Jacobiano del balance, pueden fusionarse o terminar
ramas; el límite de las fórmulas no demuestra su continuación global.

\subsection{Ejemplo calculado: Chua con saturación y \texorpdfstring{\(q=0.9998\)}{q=0.9998}}
\label{bre:bal:example}

Tome
\begin{equation}
 \begin{gathered}
 (\alpha,\beta,\gamma)=(8.4562,12.0732,0.0052),\\
 (m_0,m_1)=(-0.1768,-1.1468),\qquad q=0.9998.
 \end{gathered}
 \label{bre:bal:example-parameters}
\end{equation}
La separación da \(\ell=-1.1468\), \(\Delta=0.9700\),
\(d_0=12.0784\), \(d_1=1.0052\). Como control entero, el polinomio
de fase de \eqref{bre:bal:phase-polynomial} es
\(z^2-14.69017296z+43.794581375648\). Para el orden fraccionario
indicado, se resuelve directamente \(\Im W_q(\ii\omega)=0\).
Una de las ramas proporciona
\begin{equation}
 \begin{aligned}
 \omega&\simeq2.04028605108,\\
 W_q(\ii\omega)&\simeq4.76138797078,\\
 k&\simeq0.210022792962,\\
 \zeta&\simeq0.000640883348+2.03999498963\ii.
 \end{aligned}
 \label{bre:bal:example-frequency}
\end{equation}
Como \(0<k<\Delta\), la ecuación centrada tiene una amplitud única
saturada, \(A\simeq5.85176778549\). Para localizar una rama sesgada
se conserva la misma condición de fase y se resuelve
\begin{equation}
 c-6.79206998612\Psi_0(c,A)=0,\qquad
 \Psi_1(c,A)-kA=0.
 \label{bre:bal:example-system}
\end{equation}
Las fórmulas por tramos dan, con cifras redondeadas,
\begin{equation}
 \begin{aligned}
 c&\simeq2.776397,& A&\simeq4.577883,\\
 \theta_+&\simeq1.969299,&\theta_-&\simeq2.540861,\\
 \Psi_0&\simeq0.408770354,&\Psi_1&\simeq0.961459759.
 \end{aligned}
 \label{bre:bal:example-bias}
\end{equation}
Se verifica \(W_q(0)\Psi_0\simeq c\) y
\(\Psi_1/A\simeq k\). Las dos últimas ecuaciones de la media dan
\(\bar y=\gamma c/(\beta+\gamma)\) y
\(\bar z=-\beta c/(\beta+\gamma)\), por lo que
\begin{equation}
 \bar X\simeq(2.776397,\ 0.001195296,\ -2.775202)^{\T}.
 \label{bre:bal:example-mean}
\end{equation}
La sustitución de \(\zeta\) en \eqref{bre:bal:chua-mode} produce
\begin{equation}
 v\simeq\begin{pmatrix}
 1\\0.063298582+0.241242519\ii\\
 -1.428794382+0.370525917\ii
 \end{pmatrix}.
 \label{bre:bal:example-mode}
\end{equation}
Para fase \(\varphi=0\), la suma \(\bar X+A\Re v\) proporciona
\begin{equation}
 X_{\mathrm{sem}}\simeq
 (7.354280,\ 0.290969,\ -9.316055)^{\T}.
 \label{bre:bal:example-seed}
\end{equation}
El tercer valor propio del auxiliar es aproximadamente
\(-1.54110635\): se obtiene restando \(2\Re\zeta\) de
\(\operatorname{tr}P_k\), y está en el sector estable. Así quedan
determinados el modo, la media, la escala y el estado inicial. La
continuación usa \eqref{bre:bal:homotopy}; la naturaleza del destino se
establece después mediante integración y análisis de cuencas.

\subsection{Familia de Machado y residuo del modelo físico}
\label{bre:bal:machado}

Machado propone elevar una función descriptiva a un orden real
\cite{Machado2015FDF}. Para distinguir esta operación del orden
dinámico \(q\), se denota su exponente por \(\mu>0\). Con una escala
positiva \(N_\star\) de las mismas unidades que \(N\), la familia es
\begin{equation}
 N_{\mu,\star}=N_\star
 \exp\left[\mu\operatorname{Log}(N/N_\star)\right].
 \label{bre:bal:machado-definition}
\end{equation}
Sobre \(N>0\) es real y unívoca. Para \(N\) complejo se fija una
hoja del logaritmo y se continúa su fase; para \(N<0\), la potencia
principal es, en general, compleja. En \(N=0\), la extensión
\(0^\mu=0\) no permite utilizar el logaritmo. La normalización evita
elevar directamente una magnitud dimensional a un exponente distinto
de uno. No hay una identidad que obligue a imponer \(\mu=q\).

En la rama positiva, el cierre auxiliar
\(1-W_q(\ii\omega)N_{\mu,\star}=0\) conserva la condición de
fase, pero sustituye la ecuación de amplitud por
\begin{equation}
 N(A)=N_\star\left(\frac{k}{N_\star}\right)^{1/\mu},
 \qquad k>0.
 \label{bre:bal:machado-amplitude}
\end{equation}
Se comprueba nuevamente el rango de \(N\). Al derivar
\(\mu\log[N(A)/N_\star]=\log(k/N_\star)\), con \(k\) fijo,
se obtiene
\begin{equation}
 \frac{\dd A}{\dd\mu}
 =-\frac{N(A)\log[N(A)/N_\star]}{\mu N'(A)}.
 \label{bre:bal:machado-sensitivity}
\end{equation}
Esta fórmula identifica pérdida de condicionamiento cuando
\(N'(A)\) se aproxima a cero. El modo se calcula con la ganancia
auxiliar \(k\), y la nueva amplitud modifica la semilla.

La potencia no es el coeficiente de Fourier de la característica
original. Por ello se calculan separadamente
\begin{equation}
 \begin{aligned}
 \rho_{\mathrm{aux}}&=|1-W_q(\ii\omega)N_{\mu,\star}(A)|,\\
 \rho_{\mathrm{fis}}&=|1-W_q(\ii\omega)N(A)|.
 \end{aligned}
 \label{bre:bal:machado-residuals}
\end{equation}
En el ejemplo anterior, con \(N_\star=1\) y \(\mu=2\),
\(A\simeq2.62843351478\), \(N(A)\simeq0.458282437981\) y
\(N(A)^2\simeq k\). El cierre auxiliar tiene residuo del orden de
\(10^{-13}\), pero el físico vale aproximadamente \(1.18206048743\).
La nueva semilla diversifica la búsqueda; no resuelve el balance físico
de primer armónico.

La escala también interviene en la invariancia de representación. Bajo
\(b\mapsto ab\), \(N\mapsto N/a\), debe hacerse
\(N_\star\mapsto N_\star/a\); entonces
\(N_{\mu,\star}\mapsto N_{\mu,\star}/a\), y el cierre auxiliar no
cambia. Si se mantiene una potencia desnuda \(N^\mu\), el producto
se transforma en \(a^{1-\mu}W_qN^\mu\), que depende de la escala
elegida cuando \(\mu\ne1\).

En una rama sesgada se mantiene la media física \(\Psi_0\) y sólo
se aplica la potencia a \(N_1\). Se resuelven
\(c=W_q(0)\Psi_0\) y
\(1-W_qN_{1,\mu,\star}=0\), registrando además
\(|1-W_qN_1|\). La ganancia auxiliar correspondiente se inserta en
la homotopía afín para conservar \(F_1=F\). Esta construcción escalar
no autoriza potencias elemento a elemento en un problema multicanal.

\subsection{Balance multicanal y control del residuo completo}
\label{bre:bal:mimo}

Para no linealidades acopladas se emplea la identidad exacta
\begin{equation}
 F(X)=AX+B\Phi(CX),\qquad
 W_q(s)=C(s^qI-A)^{-1}B,
 \label{bre:bal:mimo-model}
\end{equation}
con \(A\in\R^{n\times n}\), \(B\in\R^{n\times m}\),
\(C\in\R^{p\times n}\) y \(\Phi:\R^p\to\R^m\).
Si \(K=D\Phi(C\bar X)\), la linealización tiene matriz
\(A+BKC\). Fuera de los polos, la identidad
\(\det(I-UV)=\det(I-VU)\) permite escribir
\begin{equation}
 \begin{aligned}
 \det(\zeta I-A-BKC)
 &=\det(\zeta I-A)\\
 &\quad\times\det[I_p-W_q(\ii\omega)K].
 \end{aligned}
 \label{bre:bal:mimo-linear}
\end{equation}
Aquí se fija $\zeta=(\ii\omega)^q$. Este predictor usa una derivada local.
Para amplitudes finitas se
calculan los coeficientes de Fourier de \(\Phi\) sobre la señal
propuesta, sin sustituir sus productos por ganancias escalares separadas.

Sea
\begin{equation}
 X^{(H)}(\theta)=\sum_{k=-H}^{H}X_ke^{\ii k\theta},\qquad
 X_{-k}=\overline{X_k},\quad X_0\in\R^n.
 \label{bre:bal:fourier-state}
\end{equation}
Aquí \(X_1\) es un coeficiente complejo de Fourier: una componente
\(A\cos\theta\) tiene coeficientes \(A/2\) en \(k=\pm1\).
En \eqref{bre:bal:seed}, en cambio, \(Av\) era el fasor completo.
La reconstrucción que conserva esta normalización es
\begin{equation}
 X_{\mathrm{sem}}(\varphi)=X_0+
       2\Re\sum_{k=1}^{H}X_ke^{\ii k\varphi}.
 \label{bre:bal:mimo-seed}
\end{equation}
Si \(u=\sum U_ke^{\ii k\theta}\) y
\(v=\sum V_ke^{\ii k\theta}\), al multiplicar las series y reunir
los términos con el mismo exponente se obtiene
\begin{equation}
 \widehat{uv}_k=\sum_jU_jV_{k-j},\qquad
 \widehat{uvw}_k=\sum_j\sum_lU_jV_lW_{k-j-l}.
 \label{bre:bal:convolutions}
\end{equation}
Las sumas se extienden sobre enteros, tomando coeficientes cero fuera
del intervalo retenido. Para polinomios cuadráticos o cúbicos esto
calcula exactamente la no linealidad de la señal truncada, incluidos
los armónicos generados hasta \(2H\) o \(3H\).

Por ejemplo, para el Lorenz generalizado,
\begin{equation}
 \begin{gathered}
 A=\begin{pmatrix}-\sigma&\sigma&0\\r&-1&0\\0&0&-1\end{pmatrix},
 \quad B=\operatorname{diag}(-a,-1,1),\quad C=I,\\
 \Phi(X)=(yz,xz,xy)^{\T}.
 \end{gathered}
 \label{bre:bal:lorenz-mimo}
\end{equation}
La primera componente del coeficiente no lineal es
\(\widehat\Phi_{1,k}=\sum_jX_{2,j}X_{3,k-j}\); las otras dos se
obtienen de \(xz\) y \(xy\). Para una señal de un solo armónico,
\(x=c_x+a_x\cos\theta\), \(y=c_y+a_y\cos\theta\), el producto
contiene media \(c_xc_y+a_xa_y/2\), fundamental
\((c_xa_y+c_ya_x)\cos\theta\) y segundo armónico
\((a_xa_y/2)\cos2\theta\). Así se ve por qué deben conservarse
tanto la media como el acoplamiento entre canales.

Defina
\(\widehat\Phi_k=(2\pi)^{-1}\int_0^{2\pi}
\Phi(CX^{(H)}(\theta))e^{-\ii k\theta}\,\dd\theta\).
El balance estacionario es
\begin{equation}
 \begin{aligned}
 AX_0+B\widehat\Phi_0&=0,\\
 [ (\ii k\omega)^qI-A]X_k-B\widehat\Phi_k&=0,
 \quad 1\leq k\leq H.
 \end{aligned}
 \label{bre:bal:mimo-balance}
\end{equation}
Si \(X_1\) tiene alguna componente no nula, se añade una condición de fase,
\(\Im(e_j^{\T}X_1)=0\), con \(\Re(e_j^{\T}X_1)>0\).
Si el primer coeficiente se anula, se fija la fase mediante otro
armónico no nulo.
Son \(n+2nH+1\) incógnitas reales, incluida \(\omega\), y el mismo
número de ecuaciones. La desigualdad excluye la solución estacionaria
de amplitud cero en la coordenada elegida.

Resolver las ecuaciones retenidas no anula el residuo fuera de ellas.
Sea
\(\mathcal R(\theta)=D_{\mathrm W}^qX^{(H)}-AX^{(H)}
-B\Phi(CX^{(H)})\). Sus coeficientes son
\begin{equation}
 \begin{aligned}
 \mathcal R_0&=-AX_0-B\widehat\Phi_0,\\
 \mathcal R_k&=[(\ii k\omega)^qI-A]X_k-B\widehat\Phi_k
       \quad(1\leq k\leq H),\\
 \mathcal R_k&=-B\widehat\Phi_k\quad(k>H).
 \end{aligned}
 \label{bre:bal:mimo-residual}
\end{equation}
Para señales reales, los índices negativos son los conjugados. Si el
residuo pertenece a \(L^2\), Parseval, con la norma euclídea, da
\begin{equation}
 \frac1{2\pi}\int_0^{2\pi}\norm{\mathcal R(\theta)}^2\,\dd\theta
 =\norm{\mathcal R_0}^2+2\sum_{k=1}^{\infty}\norm{\mathcal R_k}^2.
 \label{bre:bal:parseval-residual}
\end{equation}
La cola física se mide mediante \(B\widehat\Phi_k\), no sólo por
\(\widehat\Phi_k\): el campo depende de la combinación de canales
impuesta por \(B\), que puede tener núcleo. Se controlan las ecuaciones
retenidas y esta cola al aumentar \(H\). Si se emplea cuadratura de
Fourier en lugar de convoluciones exactas, se refina también la malla
angular para distinguir error de cuadratura y truncamiento.

Cuando \((\ii k\omega)^qI-A\) es invertible, el tamaño
\(\norm{[(\ii k\omega)^qI-A]^{-1}B\widehat\Phi_k}\) estima la
respuesta lineal al armónico omitido y detecta resonancias. No sustituye
un teorema de error para el problema no lineal, que requeriría controlar
la inversa del operador de balance y sus derivadas.

La secuencia de cálculo queda determinada: verificar la realización;
resolver fase, media y amplitud o el balance multiharmónico; reconstruir
la semilla con su fase; evaluar el residuo físico completo; continuar
hasta el campo original y analizar allí el PVI causal. Una raíz
algebraica aceptable debe superar tanto las condiciones de dominio
del resolvente como las comprobaciones de residuo y condicionamiento.

\section{Puntos perpetuos y localización geométrico--topológica}
\label{bre:geo:section}

La búsqueda algebraica produce condiciones iniciales antes de integrar el
sistema. La geometría reduce el dominio donde se buscan; el retorno y el
grado pueden convertir una aproximación en una prueba de existencia. Estas
tres operaciones requieren justificaciones distintas: un cero algebraico
no implica atracción y un conjunto invariante no es necesariamente oculto.

\subsection{Qué se anula en un punto perpetuo}

Para una EDO autónoma $\dot x=f(x)$, con $f\in C^1$, la regla de la cadena
da $\ddot x=Df(x)f(x)$. Escriba $A(x)=Df(x)f(x)$.
Un punto perpetuo (PP) satisface $A(p)=0$ y $f(p)\ne0$
\cite{Prasad2015Perpetual,DudkowskiPrasadKapitaniak2017}. Por tanto,
\begin{equation}
 \operatorname{PP}\subset\{x:\det Df(x)=0\}.
 \label{bre:geo:singular}
\end{equation}
En efecto, un Jacobiano invertible sólo anula el vector cero. La inclusión
es necesaria, no suficiente: después de resolver la condición singular se
sustituye en todas las componentes de $A$ y se excluyen los equilibrios.
Una raíz PP tampoco obliga a que $A(x(t))$ siga siendo cero para $t>t_0$.

En el PVI conmensurable de Caputo, $\Caputo{q}x=f(x)$, $0<q<1$, la
identidad $\Caputo{q}(f\circ x)=Df(x)f(x)$ no es válida en general. Su
interpretación correcta en el arranque se obtiene de Volterra. Para
$h=t-t_0\ge0$ y $f\in C^2$ cerca de $x_0$, sea $f_0=f(x_0)$ y
$J_0=Df(x_0)$. La identidad de potencias obtenida mediante la integral
beta en \eqref{bre:fund:power} permite integrar cada término de la
expansión; todos sus exponentes son mayores que $-1$.
En un intervalo local donde $\norm{f}\le M$, Volterra da primero
$\norm{x(t_0+h)-x_0}\le Mh^q/\Gamma(q+1)$.
La diferenciabilidad de $f$ implica $f(x(t_0+h))=f_0+O(h^q)$;
integrar de nuevo produce
$x-x_0=f_0h^q/\Gamma(q+1)+O(h^{2q})$. Al sustituir
$f(x)=f_0+J_0(x-x_0)+O(\norm{x-x_0}^2)$ en Volterra se obtiene
\begin{equation}
 x(t_0+h)=x_0+\frac{f_0h^q}{\Gamma(q+1)}
 +\frac{J_0f_0h^{2q}}{\Gamma(2q+1)}+O(h^{3q}).
 \label{bre:geo:picard}
\end{equation}
En este paso el resto $O(h^{2q})$ del integrando se transforma en
$O(h^{3q})$ por \eqref{bre:fund:power}; no se ha compuesto ninguna
derivada fraccionaria. Así, una semilla FPP--A satisface $J_0f_0=0$, $f_0\ne0$: cancela el
segundo coeficiente de esta expansión. Con $f\in C^1$ basta el límite
\begin{equation}
 \Gamma(q+1)\frac{f(x(t_0+h))-f_0}{h^q}\longrightarrow J_0f_0.
 \label{bre:geo:startup}
\end{equation}
Las raíces son independientes de $q$; su destino dinámico no lo es.
En particular, ni $\ddot x(t_0)$, generalmente singular, ni la derivada
secuencial se pueden sustituir por $J_0f_0$.

Si $v_q(t)=f(x(t))$ es absolutamente continua, un evento FPP--H depende
de toda la historia:
\begin{equation}
 \mathcal A_q[x](t)=\frac1{\Gamma(1-q)}
 \int_{t_0}^t(t-s)^{-q}\frac{\dd}{\dd s}f(x(s))\,\dd s.
 \label{bre:geo:history}
\end{equation}
El evento se define en instantes $t>t_0$ donde esta integral es finita;
la continuidad absoluta sólo asegura su existencia casi en todas partes.
Se exige $\mathcal A_q[x](t)=0$ y $f(x(t))\ne0$. Para el campo constante
$f=c\ne0$, la solución es $x=x_0+ch^q/\Gamma(q+1)$:
$\mathcal A_q=0$, aunque $\ddot x=ch^{q-2}/\Gamma(q-1)\ne0$.
El ejemplo separa las nociones sin recurrir a una aproximación numérica.
Tampoco se supone una ley de semigrupo, excluida en general por
\eqref{bre:fund:nosemigroup}; las propiedades parciales de composición
requieren dominios adicionales \cite{Cong2023Semigroup}.
Un mínimo pequeño de una discretización de \eqref{bre:geo:history} sólo
es un candidato; debe persistir al refinar paso y horizonte común.

\subsection{Coordenadas, curvatura y memoria}

Para un difeomorfismo local $y=g(x)$ de clase $C^2$, el campo de una EDO es
$\widetilde f=Dg\,f$ y
\begin{equation}
 D\widetilde f\,\widetilde f
 =Dg\,Df\,f+D^2g[f,f].
 \label{bre:geo:coordinate-law}
\end{equation}
Por ello los PP euclidianos se conservan bajo cambios afines invertibles,
pero no bajo cambios generales. Fijada una métrica y su conexión, la
aceleración covariante $\nabla_f f$ sí representa un vector intrínseco.
Su parte normal es
\begin{equation}
 a^\perp=\nabla_f f-
 \frac{\langle\nabla_f f,f\rangle}{\langle f,f\rangle}f,
 \qquad f\ne0.
 \label{bre:geo:normal}
\end{equation}
Al cambiar el tiempo mediante $\widehat f=\rho f$, con $\rho\in C^1$ y $\rho>0$,
$\nabla_{\widehat f}\widehat f=\rho^2\nabla_f f+\rho f(\rho)f$;
así $\widehat a^\perp=\rho^2a^\perp$. La condición de curvatura normal
nula es orbitalmente invariante; la condición fuerte $\nabla_f f=0$ no.

En Caputo debe transportarse la curva de arranque, no recalcularse una
ecuación fraccionaria usando la regla de la cadena ordinaria. Con
$s=h^q$, \eqref{bre:geo:picard} tiene la forma
$x=x_0+Us+Vs^2/2+o(s^2)$, donde
\begin{equation}
 U=\frac{f_0}{\Gamma(q+1)},\qquad
 V=\frac{2J_0f_0}{\Gamma(2q+1)}.
 \label{bre:geo:jet}
\end{equation}
Taylor aplicado a $g$ da exactamente
\begin{equation}
 U_y=Dg\,U,\qquad V_y=Dg\,V+D^2g[U,U].
 \label{bre:geo:jet-law}
\end{equation}
Aquí $\Gamma[U,U]$ denota las componentes
$\Gamma^i{}_{jk}U^jU^k$ de la conexión, distintas de la función gamma.
La ley de transformación de los símbolos de Christoffel cancela el
término hessiano de \eqref{bre:geo:jet-law}; por ello
$V+\Gamma[U,U]$ es un vector. Se transportan los coeficientes de la
curva, sin afirmar que el nuevo campo satisfaga una ecuación Caputo de
la misma forma. En cambio, recalcular
$2D\widetilde f\,\widetilde f/\Gamma(2q+1)$ introduce el defecto
\begin{equation}
 \left[\frac2{\Gamma(2q+1)}-\frac1{\Gamma(q+1)^2}\right]D^2g[f_0,f_0].
 \label{bre:geo:jet-defect}
\end{equation}
Para $q=1/2$, $f=(1,0)$ y $g(x)=(x_1,x_2+x_1^2)$, su segunda
componente es $4-8/\pi\ne0$. Esta diferencia obliga a declarar las
coordenadas físicas, la métrica y el tipo de historia de cada semilla.

\subsection{Resolución algebraica de las semillas}

\subsubsection{MAVPD}
Considérese
\begin{equation}
 \begin{aligned}
 f_1&=\delta(\gamma u+v-u^3),\\
 f_2&=u-\xi v-w,\qquad f_3=\rho v.
 \end{aligned}
 \label{bre:geo:mavpd}
\end{equation}
Con $\delta\rho\ne0$, las filas tercera y segunda de $Df\,f=0$
imponen $f_2=0$ y $f_1=f_3=\rho v$. La primera queda
$\delta(\gamma-3u^2)f_1=0$. En un PP, $f_1\ne0$, por lo que
$u^2=\gamma/3$. Para $\gamma>0$, $\rho\ne\delta$,
\begin{equation}
 \begin{aligned}
 u_\star&=\pm\sqrt{\gamma/3},\qquad
 v_\star=\frac{2\delta\gamma u_\star}{3(\rho-\delta)},\\
 w_\star&=u_\star-\xi v_\star.
 \end{aligned}
 \label{bre:geo:mavpd-roots}
\end{equation}
La segunda fórmula procede de
$\delta(2\gamma u_\star/3+v_\star)=\rho v_\star$; la tercera, de
$f_2=0$. Estas sustituciones comprueban todas las ecuaciones y agotan las
ramas no estacionarias bajo las hipótesis indicadas. Para
$(\delta,\rho,\gamma,\xi)=(100,200,0.1,3.1)$ hay un par simétrico.
Sus integraciones alcanzan ciclos interiores recurrentes en los horizontes
estudiados; la estabilidad asintótica requiere evidencia adicional.

\subsubsection{Muñoz--Pacheco}
Para $a=0.35$,
\begin{equation}
 f=(yz+x(y-a),\ 1-|x|,\ -xy-z)^{\T}.
 \label{bre:geo:munoz}
\end{equation}
No hay equilibrios: $f_2=0$ exige $|x|=1$, $f_3=0$ da $z=-xy$, y
\[
 f_1=-x(y^2-y+a)
 =-x\bigl[(y-\tfrac12)^2+a-\tfrac14\bigr]\ne0
\]
porque $a=0.35>1/4$ y $x=\pm1$.
En cada región $s=\operatorname{sign}x=\pm1$, la
segunda fila de $Df\,f=0$ da $f_1=0$; la tercera, $f_3=-xf_2$; la
primera, $f_2(z+x-xy)=0$. Si $f_2=0$, todo $f$ sería cero, imposible.
Entonces $z=x(y-1)$. Sustituir en $f_3=-xf_2$ da
$y=1-|x|/2$; sustituir en $f_1=0$ da $x(y^2-a)=0$. Resultan
\begin{equation}
 \begin{aligned}
 u_\varepsilon&=2(1-\varepsilon\sqrt a),\\
 p_{s,\varepsilon}&=(su_\varepsilon,\ \varepsilon\sqrt a,
                   -su_\varepsilon^2/2),
 \end{aligned}
 \label{bre:geo:munoz-roots}
\end{equation}
con $s,\varepsilon=\pm1$. Puesto que $0<a<1$, los cuatro puntos
pertenecen a las regiones usadas en la derivación; son todas sus raíces PP.

En $x=0$ no existe un Jacobiano clásico único. La derivada direccional
de $f$ en dirección $f=(yz,1,-z)$ tiene segunda componente $-|yz|$.
Su anulación exige $yz=0$; si $z\ne0$, entonces $y=0$ y la tercera
componente vale $z\ne0$. Queda la recta $(0,y,0)$, donde ambos
Jacobianos laterales anulan $f$. Estos candidatos direccionales generan
\begin{equation}
 x(t)=z(t)=0,\quad
 y(t)=y_0+\frac{(t-t_0)^q}{\Gamma(q+1)},\quad 0<q\le1,
 \label{bre:geo:munoz-boundary}
\end{equation}
como se comprueba sustituyendo en Volterra. Son soluciones no acotadas.
Las cuatro raíces suaves sí generan trayectorias finitas en los ensayos
Caputo; ese resultado por sí solo no demuestra atracción ni caos.

\subsubsection{Chua arctangente}
Sea $h(x)=a_1x+a_2\arctan(\rho x)$ y
\begin{equation}
 \begin{aligned}
 f_1&=\alpha(y-x-h(x)),\qquad f_2=x-y+z,\\
 f_3&=-\beta y-\gamma z.
 \end{aligned}
 \label{bre:geo:chua}
\end{equation}
Escriba $U=f_1$ y $p=h'(x)$. Las dos primeras filas de $Df\,f=0$
dan $f_2=(1+p)U$ y $f_3=pU$; $U=0$ sería un equilibrio. La tercera
fila impone $[\beta+(\beta+\gamma)p]U=0$. Si
$\alpha\rho(\beta+\gamma)\ne0$, se resuelve primero
\begin{equation}
 p=-\frac\beta{\beta+\gamma},\qquad
 x^2=\rho^{-2}\left(\frac{a_2\rho}{p-a_1}-1\right).
 \label{bre:geo:arctan-x}
\end{equation}
Un valor negativo de $x^2$ excluye raíces reales. Si $p=a_1$, se usa
$h'=a_1+a_2\rho/(1+\rho^2x^2)$ sin dividir: no hay solución finita
cuando $a_2\rho\ne0$; si $a_2\rho=0$ la restricción es degenerada.
Para cada $x$ admisible, $f_1=U$ y $f_2=(1+p)U$ reconstruyen
\begin{equation}
 \begin{aligned}
 y&=x+h(x)+U/\alpha,\\
 z&=h(x)+(1+p)U+U/\alpha.
 \end{aligned}
 \label{bre:geo:arctan-yz}
\end{equation}
La ecuación $f_3=pU$ equivale a
\begin{equation}
 \begin{aligned}
 &\left[\gamma+(1+\gamma)p+\frac{\beta+\gamma}{\alpha}\right]U\\
 &\hspace{2em}=-\beta x-(\beta+\gamma)h(x).
 \end{aligned}
 \label{bre:geo:arctan-U}
\end{equation}
Se resuelve esta ecuación lineal y se exige $U\ne0$. Si su coeficiente
es cero, un término derecho no nulo excluye la raíz; si ambos son cero,
la familia se comprueba en el sistema original. En Wu y c590 los
denominadores son no nulos y se obtiene un par simétrico. Los transitorios
acotados de Wu aún no resuelven un destino asintótico; el par de c590,
muy alejado y con $\norm{f}\approx545$, necesita control de escape.

\subsubsection{PLL lead--lag}
En el cilindro $(x,\theta\bmod2\pi)$, con $T=\tau_1+\tau_2>0$,
\begin{equation}
 \begin{aligned}
 f_x&=-x/T+\tau_1\sin\theta/(2T),\\
 f_\theta&=\omega_\Delta-Lx/T-\tau_2L\sin\theta/(2T).
 \end{aligned}
 \label{bre:geo:pll}
\end{equation}
El determinante es $L\cos\theta/(2T)$. Para $L\ne0$, un PP exige
$\cos\theta=0$. Allí $Df$ tiene columnas $(-1/T,-L/T)^{\T}$ y cero;
por tanto $Df\,f=0$ equivale a $f_x=0$. Así
\begin{equation}
 p_+=(\tau_1/2,\pi/2),\qquad
 p_-=(-\tau_1/2,3\pi/2),
 \label{bre:geo:pll-roots}
\end{equation}
si $\omega_\Delta\mp L/2\ne0$, respectivamente. Cada igualdad nula
convierte esa raíz en equilibrio y obliga a excluirla. La semilla positiva
del caso estudiado llega al foco; la negativa conserva un transitorio
acotado sin destino resuelto. PP localiza semillas, no selecciona por sí
solo el ciclo corredor.

\subsubsection{Exclusiones: Chua PWL, Kalman--Fitts y jerk}
Para Chua PWL, en una región de pendiente $m$,
$\det Df=-\alpha[\beta+(\beta+\gamma)m]$. Si es no nulo en todas las
regiones, \eqref{bre:geo:singular} excluye PP suaves. En una frontera
continua, la derivada direccional usa el Jacobiano de la región hacia la
que apunta $f$; si $f$ es tangente, ambos productos laterales coinciden.
La invertibilidad de ambos excluye también una raíz direccional no
estacionaria. Con $(\alpha,\beta,\gamma)=(8.4562,12.0732,0.0052)$ y
$(m_0,m_1)=(-0.1768,-1.1468)$, los determinantes son
$-84.035507517056$ y $15.037737580544$, respectivamente, ambos no nulos.

Kalman--Fitts tiene $\dot X=AX+e_4\tanh(-X_3/0.01)$, donde $A$ es
compañera de $s^4+0.12s^3+2.0254s^2+0.121308s+0.98191881$.
El término no lineal sólo modifica la tercera entrada de la última fila
del Jacobiano. En una matriz compañera de orden cuatro el determinante
es el coeficiente constante, inalterado: $\det Df=0.98191881\ne0$.
No hay PP, aunque otras rutas sí localizan ciclos.

Para cualquier jerk $f=(y,z,G)^{\T}$,
\begin{equation}
 Df\,f=(z,\ G,\ G_xy+G_yz+G_zG)^{\T}.
 \label{bre:geo:jerk}
\end{equation}
Toda raíz satisface $z=G=0$ y $G_xy=0$. En el jerk cuadrático
$G=-x-y-1.1z-2.59z^2+xy$, esto da $y(y-1)=0$. Para $y=0$ queda
el origen; para $y=1$, $G(x,1,0)=-1$, imposible. En el exponencial
$G=-az-I_c(e^{y/(nV_T)}-1)-x$, con $nV_T\ne0$, $G_x=-1$ fuerza
$y=0$ y después $x=z=0$. Ambos carecen de PP y, por
\eqref{bre:geo:picard}, de semillas FPP--A conmensurables.

\subsection{Eliminación y conteo exhaustivo: Lorenz y RF}

Los sistemas polinomiales requieren distinguir una raíz del eliminante de
una solución del sistema. El resultante $\operatorname{Res}_v(P,Q)$ es
el determinante de Sylvester: anula necesariamente toda raíz común, pero
puede incorporar caídas de grado y denominadores. Se reconstruyen las
variables y se comprueban las ramas excepcionales antes de contar PP.
Para contar raíces reales se usa la sucesión de Sturm
$S_0=P$, $S_1=P'$, $S_{j+1}=-\operatorname{rem}(S_{j-1},S_j)$.
Si $V_P(t)$ cuenta cambios de signo omitiendo ceros,
$V_P(a)-V_P(b)$ cuenta raíces distintas entre extremos que no son raíces.

\subsubsection{Lorenz generalizado}
Para el campo
\begin{equation}
 \begin{aligned}
 f_1&=-17(x-y)/5+yz/2,\\
 f_2&=34x/5-y-xz,\qquad f_3=-z+xy,
 \end{aligned}
 \label{bre:geo:lorenz}
\end{equation}
la rama $x=0$, $y\ne0$ exigiría simultáneamente
$z=-748/135$, $z^2=1206/25$; la rama $y=0$, $x\ne0$ exigiría
$z=748/135$, $z^2=1734/25$. Ambas son imposibles; $x=y=0$ da el
origen. Para $xy\ne0$, tome $u=x^2>0$, $v=y/x\ne0$. Dividir las
dos primeras componentes de $Df\,f$ entre $x$ produce
\begin{equation}
 v^2u/2+c_1=0,\quad -vu+c_2=0,\quad uc_3+z=0,
 \label{bre:geo:lorenz-reduced}
\end{equation}
con
\begin{equation}
 \begin{aligned}
 c_1&=867/25-374v/25-27vz/10-z^2/2,\\
 c_2&=-748/25+603v/25+27z/5-vz^2/2,\\
 c_3&=34/5-27v/5+17v^2/5+(v^2/2-1)z.
 \end{aligned}
 \label{bre:geo:lorenz-c}
\end{equation}
Eliminar $u=c_2/v$ deja
\begin{equation}
 \begin{aligned}
 L&=(3468-2992v+1206v^2)/25-(2+v^2)z^2,\\
 Q&=vz+c_2c_3,
 \end{aligned}
 \label{bre:geo:lorenz-LQ}
\end{equation}
y $L=Q=0$. Defina exactamente el polinomio de grado ocho
\begin{equation}
 P_L(z)=\frac{-5^{12}\operatorname{Res}_v(L,Q)}
 {(25z^2-1206)(25z^2-986)}.
 \label{bre:geo:lorenz-eliminant}
\end{equation}
La división es exacta y da
\begin{equation}
 \begin{aligned}
 P_L(z)={}&781250z^8-94421875z^6\\
 &-112200000z^5+3621441250z^4\\
 &+6045336000z^3-60977294900z^2\\
 &-150632718720z-231894425384.
 \end{aligned}
 \label{bre:geo:lorenz-polynomial}
\end{equation}
El resto de $Q$ entre $L$, respecto de $v$, es
$(R_1v+R_0)/[125(25z^2-1206)]$, donde
\begin{equation}
 \begin{aligned}
 R_0&=18307572+4398240z\\
    &\hspace{1.5em}+2697400z^2-50625z^4,\\
 R_1&=-2864840+647540z-1739100z^2\\
    &\hspace{1.5em}-422500z^3+3125z^5.
 \end{aligned}
 \label{bre:geo:lorenz-remainder}
\end{equation}
Se verifica $\gcd(R_1,(25z^2-986)P_L)=1$; así $v=-R_0/R_1$,
$u=c_2/v$, $x=\pm\sqrt u$, $y=vx$. Para $25z^2-1206=0$,
$L=0$ exigiría $v=66/187$, pero
$\gcd(25z^2-1206,Q(66/187,z))=1$ tras eliminar denominadores.
Esa rama es vacía. Fuera de ella, el coeficiente principal de $L$ no
se anula: el resultante cero equivale a una raíz común de $L$ con
su resto lineal, lo que justifica la reconstrucción. Las variaciones de Sturm,
en el orden $-\infty,-7.443,-7.442,9.321,9.322,+\infty$, son
$(5,5,4,4,3,3)$. Hay exactamente dos raíces de $P_L$; sus
reconstrucciones tienen $(z,u)\approx(-7.4423053,0.5564666)$ y
$(9.3218988,-81.6209147)$. La evaluación racional por intervalos, usando
la forma equivalente
\begin{equation}
 u=603/25-z^2/2+(-748/25+27z/5)/v,
 \label{bre:geo:lorenz-sign}
\end{equation}
da respectivamente $0.29<u<0.83$ y $-82<u<-81$ en los dos intervalos
aislantes; sólo la primera permite $x$ real.
El factor $25z^2-986$ da exactamente $u=29/5+z$:
para $z>0$, $u\approx12.0801274$ reconstruye el par de equilibrios;
para $z<0$, $u<0$. Junto al origen quedan tres
equilibrios y dos PP. La sustitución del par PP da
$\norm{f}\approx23.4221>0$; la eliminación es exhaustiva, pero su
dinámica fraccionaria desde esas semillas sigue pendiente.

\subsubsection{Rabinovich--Fabrikant}
Sea
\begin{equation}
 \begin{aligned}
 f_1&=y(z-1+x^2)+ax,\quad a=1/10,\\
 f_2&=x(3z+1-x^2)+ay,\quad b=719/2500,\\
 f_3&=-2z(b+xy).
 \end{aligned}
 \label{bre:geo:rf}
\end{equation}
Para $x=0$, $A_1=2y[(a-b)z-a]$ y
$A_3=2z[2b^2-y^2(z-1)]$. Si $y=0$ queda $z=0$; si $y\ne0$,
$z=a/(a-b)<0$ haría imposible $y^2(z-1)=2b^2>0$. Sólo queda el
origen. Para $x\ne0$, sean $u=x^2>0$, $v=xy$, $H=u+z-1$,
$K=3z+1-u$. La multiplicación $Df\,f$ da
\begin{equation}
 \begin{aligned}
 E_1=xA_1&=u(a^2+HK)+2v[a(u+H)-bz]\\
         &\quad+2(u-1)v^2,\\
 E_2=xA_2&=Dv+C,\\
 E_3=-uA_3/(2z)&=(z-1-u)v^2+2u(a-2b)v\\
               &\quad+u(uK-2b^2),
 \end{aligned}
 \label{bre:geo:rf-reduced}
\end{equation}
donde $E_3$ se usa sólo para $z\ne0$ y
\begin{equation}
 \begin{aligned}
 D&=3z^2-6uz-2z-3u^2+4u+a^2-1,\\
 C&=2u[a(3z+1-2u)-3bz].
 \end{aligned}
 \label{bre:geo:rf-DC}
\end{equation}
En $z=0$, $A_3=0$ idénticamente y
\begin{equation}
 \begin{aligned}
 \operatorname{Res}_v(E_1,E_2)|_{z=0}
 ={}&-10^{-6}u(10u-11)(10u-9)\\
 &\times[100(u-1)^2+1]\\
 &\times[100(3u-1)^2+1].
 \end{aligned}
 \label{bre:geo:rf-z0}
\end{equation}
Los factores cuadráticos son positivos; $u>0$ deja
$(u,v)=(9/10,4/5)$ y $(11/10,-6/5)$. Cada par reconstruye dos
raíces, $x=\pm\sqrt u$, $y=v/x$. No son equilibrios porque, en
$z=0$, el determinante del sistema $f_1=f_2=0$ es
$a^2+(u-1)^2>0$ y sólo admite $x=y=0$.

Para $z\ne0$, se separa $D=0$ antes de dividir. De $C=0$ y $u>0$,
$u=u_D=(1250-7035z)/2500$; imponer además $D=0$ da
$d_D(z)=65000-1953500z-967947z^2=0$. Tras sustituir $u_D$,
\begin{equation}
 \begin{aligned}
 R_D(z)&=\operatorname{Res}_v(E_1(u_D,v,z),E_3(u_D,v,z)),\\
 \gcd(d_D,R_D)&=1.
 \end{aligned}
 \label{bre:geo:rf-exception}
\end{equation}
Por tanto esa rama es vacía. Con $D\ne0$, $v=-C/D$ y las ecuaciones
restantes son los siguientes polinomios explícitos:
\begin{equation}
 \begin{aligned}
 P&=u(a^2+HK)D^2-2CD[a(u+H)-bz]\\
  &\quad+2(u-1)C^2,\\
 Q&=(z-1-u)C^2-2u(a-2b)CD\\
  &\quad+u(uK-2b^2)D^2.
 \end{aligned}
 \label{bre:geo:rf-PQ}
\end{equation}
Defina $S=32264900u^2-215503200u+82528417$,
$E=53150000u^2-71900000u+516961$ y
\begin{equation}
 P_R=\operatorname{pp}\left(
 \frac{\operatorname{Res}_z(P,Q)}{u^{13}S^4E}\right),\qquad\deg P_R=14.
 \label{bre:geo:rf-eliminant}
\end{equation}
Aquí $\operatorname{pp}$ elimina denominadores y el máximo común divisor
de coeficientes, y fija positivo el coeficiente principal. Para evaluar
directamente el polinomio, escriba $P_R(u)=\sum_{k=0}^{14}m_k10^{e_k}u^k$,
con los enteros de la tabla siguiente:
\begin{center}
\footnotesize
\begin{tabular}{@{}rrr@{}}
\hline
$k$&$m_k$&$e_k$\\\hline
0&172890165175624363345259174569803&0\\
1&14469699682074574386521765557812&2\\
2&$-18396580138862579542741337855542518$&2\\
3&15478807059059016799618863592551164&4\\
4&$-483305962260027081263570393541627373$&4\\
5&56915116746512318338898887806915&9\\
6&$-1629237154908128789498842947039615$&8\\
7&95325977171002410189945625&14\\
8&125564435311034818006205569375&13\\
9&$-433337883019314447578125$&19\\
10&81563269511051918346484375&17\\
11&$-99102930618837890625$&23\\
12&77980731527669677734375&20\\
13&$-358729705810546875$&25\\
14&722605133056640625&24\\\hline
\end{tabular}
\end{center}
Las raíces de $S$
en $u>0$ sólo levantan a $D=C=0$, rama ya excluida; las dos raíces
positivas de $E$ dan los cuatro equilibrios no nulos. En efecto,
$f_3=0$, $z\ne0$ exige $v=-b$; de $xf_1=au+vH=0$ se obtiene
$z=1-u+au/b$. La ecuación $xf_2=uK+av=0$ reduce a
$(4b-3a)u^2-4bu+ab^2=0$, exactamente $E(u)=0$ tras multiplicar
por $62500000$. Sus dos raíces son positivas porque su producto y su
suma son positivos y su discriminante es positivo. La sustitución
reconstruye dos signos de $x$ por raíz y verifica las tres componentes.
Para $P_R$, las variaciones de Sturm en
$0,0.0391,0.0392,0.0474,0.0475,+\infty$ son $(7,7,6,6,5,5)$:
exactamente dos raíces positivas.

La reconstrucción tampoco se infiere del resultante. La base reducida
mónica de $\langle P,Q,P_R\rangle$, en orden lexicográfico $z>u$, es
$\{P_R/\operatorname{lc}(P_R),z-Z(u)\}$, con $\deg Z=13$.
Esto define unívocamente $Z$ mediante los polinomios anteriores. La
división exacta da $P(u,Z)\equiv Q(u,Z)\equiv0\pmod{P_R}$ y
\begin{equation}
 \begin{aligned}
 \gcd(P_R,Z)&=\gcd(P_R,D(u,Z))=1,\\
 \gcd(P_R,bD(u,Z)-C(u,Z))&=1.
 \end{aligned}
 \label{bre:geo:rf-lift}
\end{equation}
Cada raíz positiva $u$ tiene un único $z=Z(u)$ real, no nulo, con
$D\ne0$. Se reconstruyen $v=-C/D$, $x=\pm\sqrt u$, $y=v/x$.
El último máximo común divisor excluye $v=-b$, luego $f_3\ne0$.
Los pares $(u,z)$ son aproximadamente $(0.03910075,0.82214520)$ y
$(0.04748353,1.20238081)$. Producen cuatro PP; sumados a los cuatro
del plano $z=0$, resultan ocho PP y cinco equilibrios. Cada exclusión
es algebraica; la clasificación de sus destinos exige integración.

\subsection{Superficies, retorno y prueba de ocultedad}

En una EDO $C^1$, para $h\in C^2$ defina
$L_fh=Dh\,f$ y $L_f^2h=D(L_fh)f$. Si $A$ es compacto e invariante
para el flujo completo, entonces
\begin{equation}
 A\cap\{L_fh=0\}\ne\varnothing,\qquad
 A\cap\{L_f^2h=0\}\ne\varnothing.
 \label{bre:geo:coverage}
\end{equation}
En efecto, $h|_A$ alcanza un máximo; la órbita completa por ese punto
permanece en $A$, por lo que su derivada es cero. Aplicar el mismo
argumento a $L_fh|_A$ prueba la segunda afirmación. Para $h=x_i$ se
obtienen las superficies $f_i=0$ y $(Df\,f)_i=0$
\cite{GodaraEtAl2018CriticalSurfaces}. Si $A$ es un atractor con cuenca
abierta y la superficie es regular en una intersección, la cuenca
contiene allí un abierto relativo de la superficie. Esto justifica
muestrearla; ninguna malla finita garantiza alcanzar ese abierto.

Otra familia de semillas satisface
$\operatorname{rank}[f\ \ Df\,f]\le1$, $f\ne0$, equivalente a
$Df\,f=\lambda f$. En tres dimensiones basta anular los menores
$f_iA_j-f_jA_i$. Los PP corresponden a $\lambda=0$; las curvas
conectantes generales sólo anulan la curvatura, sin un teorema de
cobertura de todas las cuencas \cite{GilmoreEtAl2010Connecting}.
Las variedades estables de sillas y los ciclos repulsivos pueden separar
destinos; una bisección entre condiciones de cuencas distintas aproxima
esa frontera mientras ambas clasificaciones sean fiables.

Para demostrar existencia de un ciclo entero se construye una sección
transversal y un retorno continuo $P(z)=\varphi_{\tau(z)}(z)$ definido
en $\overline\Omega$. Si $P(z)\ne z$ en $\partial\Omega$ y
$\deg(P-I,\Omega,0)\ne0$, existe un punto fijo interior
\cite{Amster2021MetodosTopologicos}. Una homotopía conserva el grado
mientras mantenga el dominio del retorno y no cruce un cero por la
frontera. La estabilidad requiere controlar $DP$; la ocultedad requiere
excluir su cuenca desde vecindades de todos los equilibrios.

Un ejemplo permite realizar las tres pruebas. Para $0<a<b$, tome
\begin{equation}
 \begin{aligned}
 \dot x&=-Qx-y,\qquad\dot y=x-Qy,\\
 Q&=(x^2+y^2-a^2)(x^2+y^2-b^2).
 \end{aligned}
 \label{bre:geo:cycle}
\end{equation}
El único equilibrio es el origen, pues $Q^2+1>0$. Su Jacobiano tiene
autovalores $-a^2b^2\pm\ii$, luego es un foco estable. Para $r>0$,
$x\dot x+y\dot y=-Qr^2$ y $x\dot y-y\dot x=r^2$, de donde
\begin{equation}
 \dot r=g(r)=-r(r^2-a^2)(r^2-b^2),\qquad\dot\theta=1.
 \label{bre:geo:radial}
\end{equation}
Los radios $a,b$ son ciclos de periodo $2\pi$. El signo de $g$ es
negativo en $(0,a)$, positivo en $(a,b)$ y negativo en $(b,\infty)$.
La unicidad impide cruzar esos radios y da
$r(t)\le\max\{r(0),b\}$, excluyendo explosión finita. Cada solución
radial no constante es monótona y acotada. Su límite $\ell$ debe
satisfacer $g(\ell)=0$: de lo contrario $|\dot r|$ estaría separado
de cero y no habría límite. La dirección de monotonía determina
\begin{equation}
 \mathcal B(0)=\{r<a\},\quad
 \mathcal B(C_b)=\{r>a\},\quad
 \partial\mathcal B(C_b)=C_a.
 \label{bre:geo:basins}
\end{equation}
Así $C_a$ es la separatriz y $C_b$ un atractor oculto: toda bola
centrada en el único equilibrio con radio menor que $a$ evita su cuenca.

El retorno sobre $\theta=0$ tiene tiempo $2\pi$. Elija
$a<r_-<b<r_+$. Entonces $P(r_-)-r_->0$ y $P(r_+)-r_+<0$.
La homotopía $(1-\lambda)(P(r)-r)+\lambda(b-r)$ conserva esos signos
en la frontera, de modo que $\deg(P-I,(r_-,r_+),0)=-1$.
Existe un punto fijo; la monotonía radial lo identifica como $b$.
La ecuación variacional $\dot u=g'(b)u$, $u(0)=1$, produce
\begin{equation}
 P'(b)=e^{-4\pi b^2(b^2-a^2)}\in(0,1),
 \label{bre:geo:multiplier}
\end{equation}
porque $g'(b)=-2b^2(b^2-a^2)$. Grado, multiplicador y cuenca prueban,
respectivamente, existencia, atracción y ocultedad; el ejemplo es periódico,
no caótico. El grado del propio campo sólo detectaría el equilibrio.

Para problemas sin reducción radial, una alternativa divide un dominio
$D$ en cajas y encierra rigurosamente sus imágenes a tiempo $\tau$.
Si un bloque $N_A$ cumple
$\varphi_\tau(N_A)\subset\operatorname{int}N_A$ y contiene un atractor
$A\subset\operatorname{int}N_A$, un grafo exterior sin camino desde
vecindades $U_E$ de todos los equilibrios hasta $N_A$ excluye
$U_E\cap\mathcal B(A)$, siempre que esas órbitas no salgan de $D$.
De pertenecer alguna a la cuenca, acabaría entrando en $N_A$ y sus
imágenes sucesivas formarían el camino prohibido. Un grafo de muestras
puntuales puede omitir imágenes y no permite esta prueba.

Estas construcciones de flujo y retorno no se trasladan a Caputo por
reinicios puntuales. En el operador de Volterra un punto fijo representa
una trayectoria en un intervalo, no un ciclo; con terminal inferior
finito no hay soluciones autónomas no constantes exactamente periódicas
bajo las hipótesis usuales \cite{TavazoeiHaeri2009,KaslikSivasundaram2012}.
Una certificación asintótica exige una dinámica de historias admisibles,
compacidad y atracción en su topología relativa. Las superficies y
semillas de arranque conservan utilidad computacional, pero sus
conclusiones se limitan al PVI y al protocolo de memoria declarados.

\section{Integración causal y validación numérica}
\label{bre:num:seccion}

La integración determina si una semilla obtenida mediante balance armónico o
criterios geométricos alcanza un régimen persistente del sistema estudiado.
Para interpretar esa trayectoria deben fijarse el campo, los parámetros, el
orden, la condición inicial, el terminal inferior y el tratamiento de la
memoria. La comparación entre métodos conserva estos datos; los cambios de
paso, horizonte o aproximación del núcleo se evalúan por separado.

\subsection{Discretización de Volterra mediante ABM--PECE}
\label{bre:num:pece}

Considérese el problema conmensurable
\begin{equation}
 {}^{\mathrm C}D_{t_0}^{q}X(t)=F(t,X(t)),\qquad X(t_0)=X_0,
 \quad 0<q\leq1,
 \label{bre:num:pvi}
\end{equation}
con $F$ continua y localmente Lipschitz respecto del estado. En un intervalo
de existencia, su forma integral es
\begin{equation}
 X(t)=X_0+\frac1{\Gamma(q)}
 \int_{t_0}^{t}(t-s)^{q-1}F(s,X(s))\,\dd s.
 \label{bre:num:volterra}
\end{equation}
El esquema de Adams--Bashforth--Moulton con secuencia PECE integra este núcleo
exactamente y aproxima el campo dentro de cada intervalo
\cite{DiethelmFordFreed2002}. En la malla $t_n=t_0+nh$, sea
$F_j=F(t_j,X_j)$. Sustituir el campo por $F_j$ en $[t_j,t_{j+1}]$ da
\[
 \int_{t_j}^{t_{j+1}}(t_{n+1}-s)^{q-1}\,\dd s
 =\frac{h^q}{q}\bigl[(n+1-j)^q-(n-j)^q\bigr].
\]
Como $q\Gamma(q)=\Gamma(q+1)$, el predictor es
\begin{align}
 X_{n+1}^{\mathrm P}
 &=X_0+\frac{h^q}{\Gamma(q+1)}\sum_{j=0}^{n}b_{j,n+1}F_j,
 \label{bre:num:predictor}\\
 b_{j,n+1}&=(n+1-j)^q-(n-j)^q.\nonumber
\end{align}

Para el corrector se interpola linealmente:
\[
 F(t_j+hu,X(t_j+hu))\approx(1-u)F_j+uF_{j+1},
 \qquad 0\leq u\leq1.
\]
Con $r=n+1-j$, las contribuciones del intervalo son
$h^q[L_q(r)F_j+R_q(r)F_{j+1}]$, antes del factor $1/\Gamma(q)$. Las
integrales de las dos funciones base se calculan mediante
\begin{align}
 I_q(r)&=\int_0^1(r-u)^{q-1}\,\dd u
       =\frac{r^q-(r-1)^q}{q},\nonumber\\
 R_q(r)&=\int_0^1u(r-u)^{q-1}\,\dd u\nonumber\\
       &=rI_q(r)-\frac{r^{q+1}-(r-1)^{q+1}}{q+1},\nonumber\\
 L_q(r)&=I_q(r)-R_q(r).
 \label{bre:num:pesos-integrales}
\end{align}
La expresión de $R_q$ sigue del cambio $v=r-u$, que transforma $u$ en
$r-v$. Al reunir denominadores se obtiene
\begin{align*}
 q(q+1)L_q(r)&=(q+1-r)r^q+(r-1)^{q+1},\\
 q(q+1)R_q(r)&=r^{q+1}-(r+q)(r-1)^q.
\end{align*}
El nodo inicial recibe sólo $L_q(n+1)$; un nodo interior $j$ recibe
$R_q(n+2-j)+L_q(n+1-j)$; el extremo nuevo recibe $R_q(1)$. En un nodo
interior, los términos proporcionales a $r^q$ se combinan como
$[(q+1-r)-(r+1+q)]r^q=-2r^{q+1}$. Por tanto, los pesos normalizados son
\begin{align}
 a_{0,n+1}&=n^{q+1}-(n-q)(n+1)^q,\nonumber\\
 a_{j,n+1}&=(n-j+2)^{q+1}+(n-j)^{q+1}\nonumber\\
 &\hspace{1.1cm}-2(n-j+1)^{q+1},\quad 1\leq j\leq n.
 \label{bre:num:pesos-corrector}
\end{align}
Además, $q(q+1)R_q(1)=1$ y
$q(q+1)\Gamma(q)=\Gamma(q+2)$. Evaluar el campo desconocido del extremo en
el predictor produce
\begin{equation}
 \begin{split}
 X_{n+1}=X_0+\frac{h^q}{\Gamma(q+2)}\Bigl[
 &F(t_{n+1},X_{n+1}^{\mathrm P})\\
 &+\sum_{j=0}^{n}a_{j,n+1}F_j\Bigr].
 \end{split}
 \label{bre:num:corrector}
\end{equation}
La secuencia consiste en predecir, evaluar el campo, corregir y evaluar de
nuevo $F_{n+1}=F(t_{n+1},X_{n+1})$. Este último valor se incorpora a la
historia. Las sumas comienzan siempre en $j=0$: la implementación directa de
$N$ pasos requiere $O(N^2)$ operaciones y $O(N)$ almacenamiento para una
dimensión de estado fija.

\subsection{Límite entero y significado del predictor}
\label{bre:num:limite-entero}

Para $q=1$, $b_{j,n+1}=1$, $a_{0,n+1}=1$ y los pesos interiores del
corrector son $2$. Las fórmulas recuperan las cuadraturas compuestas
rectangular izquierda y trapezoidal de la ecuación integral; el extremo
derecho de la segunda se evalúa en el predictor. En particular,
\[
 X_{n+1}^{\mathrm P}=X_0+h\sum_{j=0}^{n}F_j
 =X_n^{\mathrm P}+hF_n\qquad(n\geq1).
\]
La recurrencia parte del predictor anterior, mientras que el estado corregido
satisface
\begin{align*}
 X_n&=X_0+\frac h2\Bigl[F_0+2\sum_{j=1}^{n-1}F_j
                         +F(t_n,X_n^{\mathrm P})\Bigr],\\
 X_n^{\mathrm P}-X_n&=\frac h2\bigl[F_0-F(t_n,X_n^{\mathrm P})\bigr].
\end{align*}
Por ello $X_{n+1}^{\mathrm P}$ no es, en general, el paso local
$X_n+hF_n$ iniciado en el estado corregido. Para $\dot x=x$, $x_0=1$, el
primer corrector da $x_1=1+h+h^2/2$, mientras que
\[
 x_2^{\mathrm P}-(x_1+hx_1)=-h^2/2.
\]
Este ejemplo distingue las identidades exactas de los pesos de una
identificación incorrecta con un método local de un paso.

\subsection{Continuación con historia y reinicio del sistema objetivo}
\label{bre:num:continuacion}

Sea $F_\eta$ una familia que une un sistema auxiliar, en $\eta=0$, con el
campo objetivo, en $\eta=1$. Se prescribe $\eta(t)=\eta_k$ en
$[t_k,t_{k+1})$ y se resuelve un único problema
${}^{\mathrm C}D_{t_0}^{q}X=F_{\eta(t)}(X)$. Cambiar $\eta$ modifica el
campo a partir de ese instante y conserva las evaluaciones anteriores.

La contribución de la historia previa al nivel $k$ es
\begin{equation}
 H_k(t)=\frac1{\Gamma(q)}\int_{t_0}^{t_k}
             (t-s)^{q-1}F_{\eta(s)}(X(s))\,\dd s.
 \label{bre:num:historia-heredada}
\end{equation}
Separar la integral en $t_k$ y usar $X(t_k)=X_0+H_k(t_k)$ da
\begin{equation}
 \begin{split}
 X(t)={}&X(t_k)+H_k(t)-H_k(t_k)\\
 &+\frac1{\Gamma(q)}\int_{t_k}^{t}
            (t-s)^{q-1}F_{\eta_k}(X(s))\,\dd s.
 \end{split}
 \label{bre:num:continuacion-integral}
\end{equation}
Un reinicio en $t_k$ elimina $H_k(t)-H_k(t_k)$. Para $q=1$, el núcleo es
constante y esa diferencia se anula. Para $0<q<1$, el peso de toda la
historia cambia al avanzar $t$, de modo que el estado terminal no basta para
reconstruir la evolución heredada.

Si el parámetro cambia por saltos, el campo dependiente del tiempo es
continuo por tramos. La formulación integral conserva sentido, pero una
estimación de interpolación basada en derivadas temporales acotadas no puede
atravesar esos saltos sin tratamiento adicional. Los tiempos \(t_k\) se
incluyen como fronteras de la partición de cuadratura. En cada lado se usa la
rama correspondiente; una interpolación lineal del campo debe distinguir
\(F_{\eta_{k-1}}(X(t_k))\) de \(F_{\eta_k}(X(t_k))\). La trayectoria es
continua, aunque los dos valores del campo puedan diferir. Si una
discretización usa un único valor nodal, el defecto causado por esa elección
forma parte del control numérico de la transición.

Al finalizar la continuación puede definirse un experimento autónomo nuevo,
\[
 {}^{\mathrm C}D_0^q Z(t)=F_1(Z(t)),\qquad Z(0)=X(t_{\mathrm{fin}}).
\]
Su objetivo es comparar esa condición inicial con los sondeos bajo un mismo
terminal inferior y una misma inicialización de memoria. La continuación
transporta una historia; el reinicio establece la trayectoria de referencia
del campo objetivo. Cambiar el orden $q$ también exige especificar un nuevo
problema, pues cambia el núcleo de la ecuación integral.

\subsection{Comparación con la recurrencia ADM local}
\label{bre:num:adm-local}

Wu et al. distinguen la recurrencia local \(X_{n+1}=\Phi_h(X_n)\) de un
predictor--corrector que utiliza la historia. Su ecuación (27) conserva
explícitamente el incremento de esa historia, y los autores advierten la
pérdida de memoria histórica del ADM por subintervalos en integraciones
largas \cite{Wu2023ArctanChua}. La diferencia se precisa restando la ecuación
de Volterra en \(t_n\) de la misma ecuación en \(t_{n+1}\):
\begin{multline}
 X(t_{n+1})-X(t_n)=\mathcal M_n\\
 +\frac1{\Gamma(q)}\int_{t_n}^{t_{n+1}}
 (t_{n+1}-s)^{q-1}F(s,X(s))\,\dd s,
 \label{bre:num:incremento-exacto}
\end{multline}
\begin{equation}
 \begin{split}
 \mathcal M_n&=\frac1{\Gamma(q)}\int_{t_0}^{t_n}B_n(s)F(s,X(s))\,\dd s,\\
 B_n(s)&=(t_{n+1}-s)^{q-1}-(t_n-s)^{q-1}.
 \end{split}
 \label{bre:num:incremento-memoria}
\end{equation}
La integral sobre $[t_n,t_{n+1}]$ describe el intervalo nuevo.
$\mathcal M_n$ repondera los intervalos anteriores y equivale a
\(H_n(t_{n+1})-H_n(t_n)\) en
\eqref{bre:num:continuacion-integral}, para un campo fijo. Se anula cuando
\(q=1\), pero en general no lo hace para \(0<q<1\). Una expansión ADM
local, aunque aumente su orden de truncamiento, no recupera ese término si
se reinicia usando sólo \(X_n\).

El campo constante \(F(t,x)=v\ne0\), con \(t_0=0\), proporciona un
contraejemplo explícito. La solución causal exacta y su incremento son
\begin{align*}
 X(t)&=X_0+\frac{t^q}{\Gamma(q+1)}v,\\
 X(t_{n+1})-X(t_n)
 &=\frac{h^q}{\Gamma(q+1)}[(n+1)^q-n^q]v.
\end{align*}
Resolver exactamente cada intervalo reiniciado produciría, en cambio,
\(\widehat X_{n+1}-\widehat X_n=h^qv/\Gamma(q+1)\). La contribución omitida
es
\[
 \mathcal M_n=\frac{h^q}{\Gamma(q+1)}
               [(n+1)^q-n^q-1]v.
\]
En un horizonte fijo \(T=Nh\), el reinicio repetido da
\[
 \widehat X_N=X_0+\frac{Th^{q-1}}{\Gamma(q+1)}v,
 \qquad
 X(T)=X_0+\frac{T^q}{\Gamma(q+1)}v.
\]
Para \(q<1\), la primera expresión ni siquiera converge a la segunda al
reducir \(h\). El error procede de omitir la historia, aun con solución
local exacta. Por ello una reconstrucción del algoritmo ADM publicado y una
integración causal con terminal fijo cumplen funciones de comparación
diferentes. La transferencia de una semilla entre ambas rutas requiere
reintegrarla bajo el operador de Caputo y el convenio de memoria elegidos,
y evaluar de nuevo su destino.

\subsection{Defecto integral, error de memoria y refinamiento}
\label{bre:num:error}

Una curva numérica continua $\widetilde X$ se contrasta con el problema de
Caputo mediante su defecto de Volterra,
\begin{equation}
 \begin{split}
 \rho_h(t)={}&\widetilde X(t)-X_0\\
 &-\frac1{\Gamma(q)}\int_{t_0}^{t}
       (t-s)^{q-1}F(s,\widetilde X(s))\,\dd s.
 \end{split}
 \label{bre:num:defecto}
\end{equation}
Supóngase que la solución y la aproximación permanecen en una región donde
$F$ es Lipschitz respecto del estado, con constante $L$. Restar sus
ecuaciones y tomar normas produce, para $t_0=0$,
\[
 u(t)\leq A+\frac L{\Gamma(q)}\int_0^t(t-s)^{q-1}u(s)\,\dd s,
 \quad u=\norm{X-\widetilde X},
\]
donde $A=\norm{\rho_h}_{\infty,[0,T]}$. La estimación se deriva iterando
la integral fraccionaria $I^q$. La identidad beta
\[
 \int_0^t(t-s)^{a-1}s^{b-1}\,\dd s
 =t^{a+b-1}\frac{\Gamma(a)\Gamma(b)}{\Gamma(a+b)}
\]
para $a,b>0$ implica $I^aI^b=I^{a+b}$ y
$I^{jq}1=t^{jq}/\Gamma(jq+1)$. Después de $N$ sustituciones,
\[
 u(t)\leq A\sum_{j=0}^{N-1}\frac{L^jt^{jq}}{\Gamma(jq+1)}
             +L^N(I^{Nq}u)(t).
\]
El último término está acotado por
$\norm{u}_\infty(LT^q)^N/\Gamma(Nq+1)$ y tiende a cero. Así se obtiene
la desigualdad de Gronwall fraccionaria \cite{Gomoyunov2018Convex},
\begin{equation}
 \sup_{0\leq t\leq T}\norm{X(t)-\widetilde X(t)}
 \leq A E_q(LT^q).
 \label{bre:num:gronwall}
\end{equation}
Para $q=1$, $E_1(z)=\exp z$. La cota es condicional a la permanencia en
la región indicada; una certificación requiere justificar también esa
permanencia y acotar el defecto, incluida su cuadratura.

Una aproximación de memoria debe controlar el núcleo en una norma adecuada.
Para $K_q(t)=t^{q-1}/\Gamma(q)$, una suma finita
$K_m(t)=\sum_{\ell=1}^m w_\ell\exp(-\eta_\ell t)$, con pesos y tasas
positivos, es acotada en cero y no aproxima uniformemente la singularidad si
$q<1$. Se separa $[0,\delta]$ y se exige
\begin{equation}
 \begin{split}
 \varepsilon_K(T)&=\int_0^T|K_q(s)-\widetilde K_m(s)|\,\dd s\\
 &\leq\varepsilon_{\mathrm{loc}}+(T-\delta)\varepsilon_\infty,
 \qquad T\geq\delta.
 \end{split}
 \label{bre:num:error-nucleo}
\end{equation}
Aquí $\widetilde K_m$ incluye el tratamiento local elegido,
$\varepsilon_{\mathrm{loc}}$ es su error integral en $[0,\delta]$ y
$\varepsilon_\infty$ acota el error en $[\delta,T]$. Una corrección local
exacta hace $\varepsilon_{\mathrm{loc}}=0$, pero conserva la necesidad de
almacenar la historia reciente.

Si $\norm{F}\leq M$ en la región y el defecto de cuadratura es a lo sumo
$\varepsilon_{\mathrm{quad}}$, sumar y restar $K_q*F(X_m)$ separa el error
de campo del error de núcleo. La segunda contribución está acotada por
$M\varepsilon_K$, por lo que \eqref{bre:num:gronwall} da
\begin{equation}
 \norm{X-X_m}_{\infty,[0,T]}
 \leq(M\varepsilon_K+\varepsilon_{\mathrm{quad}})E_q(LT^q).
 \label{bre:num:error-soe}
\end{equation}
Esta estimación tiene horizonte finito. Por ejemplo, sin corrección local,
las variables $z_\ell(t)=\int_0^t e^{-\eta_\ell(t-s)}F(s,X_m(s))\,\dd s$
satisfacen
\[
 \dot z_\ell=-\eta_\ell z_\ell+F(t,X_m),\qquad
 X_m=X_0+\sum_\ell w_\ell z_\ell.
\]
Para $F(x)=-ax$, $a>0$, el equilibrio proyectado es
$x_m^*=x_0/(1+a\sum_\ell w_\ell/\eta_\ell)$, pues
$z_\ell^*=-ax_m^*/\eta_\ell$. Puede diferir del límite cero del problema
de Caputo. Por ello la concordancia en $[0,T]$ no establece por sí sola
equivalencia de atractores o cuencas al llevar $T$ a infinito.

\paragraph{Regularidad y error de cuadratura.}
Los pesos exactos no eliminan la singularidad del arranque. Si \(F\) es
\(C^1\) cerca de \((t_0,X_0)\), la acotación del campo en
\eqref{bre:num:volterra} da \(X(t_0+s)-X_0=O(s^q)\). Por Lipschitz,
\(F(t_0+s,X(t_0+s))=F_0+O(s^q)\), donde \(F_0=F(t_0,X_0)\). Una nueva
sustitución en la integral produce
\begin{equation}
 X(t_0+s)=X_0+\frac{s^q}{\Gamma(q+1)}F_0+O(s^{2q}).
 \label{bre:num:arranque}
\end{equation}
Si \(0<q<1\) y \(F_0\ne0\), el cociente
\((X(t_0+s)-X_0)/s\) tiene norma no acotada cuando \(s\to0^+\). Así, la
suavidad del campo no implica una solución con derivadas temporales acotadas
en el terminal inferior.

Para precisar el papel de la regularidad, supóngase provisionalmente que
\(g(t)=F(t,X(t))\) tiene derivadas temporales acotadas
\(\norm{g'}\leq M_1\) y \(\norm{g''}\leq M_2\) en toda la ventana.
La interpolación constante tiene error a lo sumo \(M_1h\). La interpolación
lineal \(\ell_j\) en un intervalo de longitud \(h\) satisface
\[
 \norm{g(t)-\ell_j(t)}
 \leq\frac{M_2}{2}(t-t_j)(t_{j+1}-t)
 \leq\frac{M_2h^2}{8}.
\]
Esta cota se obtiene integrando dos veces el resto, nulo en ambos extremos;
el núcleo de esa representación tiene integral absoluta
\((t-t_j)(t_{j+1}-t)/2\). Multiplicar ambos errores por el núcleo positivo y
usar \(\int_0^T K_q(s)\,\dd s=T^q/\Gamma(q+1)\) da
\[
 \varepsilon_{\mathrm P}\leq\frac{M_1T^q}{\Gamma(q+1)}h,
 \qquad
 \varepsilon_{\mathrm L}\leq\frac{M_2T^q}{8\Gamma(q+1)}h^2.
\]
Estas son cotas de cuadratura sobre valores exactos de la trayectoria.
Sustituir el campo del extremo por su evaluación en el predictor añade a la
segunda cota, a lo sumo,
\[
 \frac{Lh^q}{\Gamma(q+2)}\varepsilon_{\mathrm P}.
\]
La consistencia resultante es \(O(h^2+h^{1+q})\) bajo esas hipótesis. Pasar
a un orden global del método requiere controlar además la propagación de
los errores nodales. Si las derivadas de \(g\) no son acotadas cerca de
\(t_0\), las cotas anteriores no son aplicables con constantes uniformes;
se necesita estimar el tramo inicial, introducir correcciones de arranque o
analizar una malla graduada. En las evaluaciones a paso uniforme se conserva
la evidencia de refinamiento sin atribuirle automáticamente el orden del
régimen suave.

\paragraph{Comparaciones a resolución creciente.}
El refinamiento temporal compara el mismo problema con $h,h/2,h/4$, usando
una malla común para medir el error de trayectoria en una ventana corta.
Después compara observables postransitorios al ampliar $T$ y el tiempo de
descarte. Si se aproxima la memoria, se varían además sus modos o longitud,
con error de núcleo controlado. La persistencia de una etiqueta clasifica
el destino al nivel de resolución evaluado; para estimar un orden de
convergencia se necesitan diferencias cuantitativas de la misma magnitud.
Con la distancia de \eqref{bre:num:distancia}, en una ventana corta fija \(I_c\) se define
\[
 E_h=\max_{t\in I_c}d_S(X_h(t),X_{h/2}(t)).
\]
Si \(E_h\simeq Ch^p\), con \(C>0\), entonces
\(E_h/E_{h/2}\simeq2^p\) y
\(p\simeq\log_2(E_h/E_{h/2})\). Esta estimación requiere diferencias no
nulas y un régimen asintótico de discretización; no se aplica a trayectorias
caóticas arbitrariamente largas. Perturbar $X_0$ a paso fijo evalúa otra propiedad: la robustez frente
a condiciones iniciales.

\subsection{Escalas, crecimiento tangente y comparación de nubes}
\label{bre:num:metricas}

La escala de las variables se fija mediante
$S=\operatorname{diag}(s_1,\ldots,s_n)$, $s_i>0$, y la distancia
\begin{equation}
 d_S(x,y)=\norm{S^{-1}(x-y)}_2.
 \label{bre:num:distancia}
\end{equation}
Para comparar escalas temporales se fija además \(\tau_*>0\). Con
\(t=t_0+\tau_*\theta\), \(Y(\theta)=S^{-1}X(t_0+\tau_*\theta)\), el
cambio de variable en Volterra da
\[
 {}^{\mathrm C}D_0^qY(\theta)
 =\tau_*^q S^{-1}F(t_0+\tau_*\theta,SY(\theta)).
\]
El factor \(\tau_*^q\) proviene del producto
\(\tau_*^{q-1}\tau_*\) entre el núcleo y el diferencial. Por tanto,
\(\widehat h=h/\tau_*\), \(\widehat T=T/\tau_*\) para una duración física \(T\); las escalas se
fijan antes de clasificar destinos. Si se usa una tasa espectral
\(\lambda\) de una linealización fraccionaria, su dimensión es
\(\mathrm{tiempo}^{-q}\), y una escala dimensionalmente coherente es
\(\tau_*=|\lambda|^{-1/q}\), para \(\lambda\ne0\).

Con $S=I$ se recupera la distancia euclidiana. En una variable angular se
usa la diferencia mínima módulo $2\pi$. Si se identifica un grupo finito
$G$ de simetrías, la expresión
$d_{S,G}(x,y)=\min_{g\in G}d_S(x,gy)$ define una métrica entre órbitas
cuando cada $g$ es una isometría. Para $g(x)=A_gx+b_g$, esto exige
$A_g^{\mathsf T}S^{-2}A_g=S^{-2}$. La invariancia permite cambiar
representantes, y la desigualdad triangular sigue de
$d_S(x,ghz)\leq d_S(x,gy)+d_S(gy,ghz)$.

En orden entero, para un campo \(C^1\), la matriz tangente satisface
\(\dot Y=DF(X(t))Y\), \(Y(0)=I\). En un campo continuo y \(C^1\) por
regiones, la integración variacional usa el Jacobiano de cada región y
resuelve los cruces de sus interfaces. En un cruce transversal de un campo
continuo, la matriz tangente es continua; los contactos tangenciales
requieren un análisis específico. Sobre una ventana de longitud $T$, el
crecimiento en la norma escalada se calcula con
$Y_S=S^{-1}YS$. Puesto que
$\norm{Y_Sv}_2^2=v^{\mathsf T}Y_S^{\mathsf T}Y_Sv$, los factores
principales de amplificación son los valores singulares $\sigma_i(Y_S)$.
Los exponentes de tiempo finito son
\begin{equation}
 \Lambda_i(T)=\frac1T\log\sigma_i(Y_S(T)).
 \label{bre:num:ftle}
\end{equation}
La acumulación de factores triangulares mediante reortogonalización QR
proporciona estimadores de los exponentes de Lyapunov
\cite{Benettin1980a}; sus valores finitos dependen de la ventana y de la
base inicial. Un exponente máximo positivo indica expansión durante el
intervalo observado. El horizonte heurístico de predicción
$T_{\mathrm{pred}}\simeq\lambda_1^{-1}
\log(\varepsilon_{\mathrm{adm}}/\varepsilon_0)$ se obtiene al resolver
$\varepsilon_0e^{\lambda_1T}=\varepsilon_{\mathrm{adm}}$, para
$0<\varepsilon_0<\varepsilon_{\mathrm{adm}}$ y $\lambda_1>0$
\cite{VallejoSanjuan2017Predictability}. Después de ese horizonte son
especialmente relevantes las distribuciones de retorno, los tiempos de
escape y los observables estadísticos convergentes.

En Caputo, la ecuación variacional conserva el mismo terminal inferior y
la historia del campo tangente. Para $F\in C^1$ en una región que contenga
las trayectorias próximas durante $[0,T]$, derivar Volterra respecto del
dato inicial produce
\begin{equation}
 Y(t)=I+\frac1{\Gamma(q)}\int_0^t(t-s)^{q-1}J(s)Y(s)\,\dd s.
 \label{bre:num:tangent-memory}
\end{equation}
Aquí $J(s)=DF(X(s))$. La derivación se justifica tomando cocientes
incrementales: el teorema
fundamental del cálculo expresa la diferencia de $F$ mediante la integral
de $DF$ sobre el segmento entre las soluciones. Su convergencia uniforme
en el compacto y la desigualdad integral de Gronwall hacen converger esos
cocientes a la única solución lineal anterior. En regiones no suaves se
necesita justificar separadamente esa diferenciabilidad.

Si en un tiempo $\tau$ se factoriza $Y(\tau)=QR$ con $R$ invertible,
el cambio de base $Z(t)=Y(t)R^{-1}$ satisface
\begin{equation}
 Z(t)=R^{-1}+\frac1{\Gamma(q)}
 \int_0^t(t-s)^{q-1}J(s)Y(s)R^{-1}\,\dd s.
 \label{bre:num:qr-memory}
\end{equation}
Por tanto, se transforman también el dato inicial tangente y todos los
valores históricos de $JY$. Conservar sólo $Z(\tau)=Q$ y reiniciar la
integral en $\tau$ elimina una contribución no constante del pasado.
Para $q=1$ esa contribución sí se reduce al estado tangente en $\tau$.
La identidad exige invertibilidad del factor utilizado; ésta debe
comprobarse, pues una matriz fundamental fraccionaria no hereda sin prueba
todas las propiedades de la de una EDO.
Al comparar algoritmos variacionales fraccionarios, incluidos los de
\cite{DancaKuznetsov2018Lyapunov}, se comprueba esta compatibilidad antes
de interpretar sus exponentes.

Para una nube objetivo $A=\{a_k\}$ y una cola de trayectoria
$B=\{b_j\}_{j=1}^{N}$, un criterio de contacto es
\begin{equation}
 D_{90}(B,A)=Q_{0.90}\bigl(\{\min_k d_S(b_j,a_k)\}_j\bigr)
 \leq\varepsilon_A.
 \label{bre:num:contacto}
\end{equation}
El percentil se calcula después del submuestreo declarado. Si las distancias
ordenadas son $d_{(1)}\leq\cdots\leq d_{(N)}$, se interpola entre las
posiciones $i+1$ e $i+2$, con $i=\lfloor0.90(N-1)\rfloor$ y peso
$0.90(N-1)-i$; para $N=1$ se toma la única distancia. Es un estadístico
dirigido: una sonda puede aproximar parte de la nube sin recorrerla toda.
La referencia, las escalas y el umbral se fijan antes del sondeo. La
concordancia geométrica se contrasta además con el destino terminal y la
estabilidad de los observables al variar paso y horizonte.

\subsection{Sondeos, salidas y censura temporal}
\label{bre:num:sondeos}

Alrededor de cada equilibrio $e$ se distinguen muestras volumétricas y
superficiales. En dimensión $n$, una muestra uniforme en una bola euclidiana
de radio $r$ es $x_0=e+rU^{1/n}d$, donde $U$ es uniforme en $[0,1]$ y
$d$ es independiente y uniforme en la esfera unidad. En efecto,
$\Pr(rU^{1/n}\leq a)=(a/r)^n$, la fracción de volumen contenida en el
radio $a$. En el sondeo superficial se usa $x_0=e+rd$; una distribución
determinista de direcciones proporciona un conteo a una resolución angular
fija. Los radios, direcciones y cantidades de muestras determinan el alcance
de la ausencia de contactos.

Para separar recurrencia y escape, sea $D$ una región cerrada de observación de un
flujo $\varphi_t$. El tiempo de primera salida es
\begin{equation}
 \tau_D(x_0)=\inf\{t\geq0:\varphi_t(x_0)\notin D\},
 \qquad \inf\varnothing=\infty.
 \label{bre:num:salida}
\end{equation}
Una integración hasta $T$ proporciona el dato censurado
\begin{equation}
 \widetilde\tau_D=\min(\tau_D,T),\qquad
 \Delta_D=\mathbf1_{\{\tau_D\leq T\}}.
 \label{bre:num:censura}
\end{equation}
En el cálculo se especifican la función de evento, su dirección de cruce,
la interpolación y la tolerancia. Si localmente $D=\{g\leq0\}$, una salida
transversal cumple $g(X(\tau_D))=0$ y
$\nabla g\cdot F>0$. Para $g_j\leq0<g_{j+1}$, el cero de la interpolación
lineal es
\[
 \widehat\tau_D=t_j-
 \frac{g_j}{g_{j+1}-g_j}(t_{j+1}-t_j).
\]
Si el cruce es único y la derivada temporal del evento es al menos
$\nu>0$, un error uniforme $\varepsilon_g$ de la función de evento implica
un error de raíz a lo sumo $\varepsilon_g/\nu$, más la tolerancia del
buscador. Esta estimación incluye el error de trayectoria y de
interpolación. Sin control entre muestras pueden omitirse salidas y
reentradas. La ausencia de escape detectado se considera censura derecha;
un fallo del integrador constituye un resultado distinto.

Para datos iniciales en un conjunto medible $U\subset D$ y una medida con
$0<\mu(U)<\infty$, la supervivencia se define por
\begin{equation}
 S_D(t)=\frac{\mu\{x_0\in U:\tau_D(x_0)>t\}}{\mu(U)}.
 \label{bre:num:supervivencia}
\end{equation}
Con $N$ semillas de igual peso y un horizonte común, para $t<T$ se estima
mediante $N^{-1}\sum_j\mathbf1_{\{\widetilde\tau_D^{(j)}>t\}}$, porque
$\min(\tau_D,T)>t$ equivale a $\tau_D>t$ en ese intervalo. Los intervalos
de confianza requieren un modelo de muestreo; los conteos de una malla
determinista no lo proporcionan por sí solos. Para una trayectoria de Caputo, \eqref{bre:num:salida} se aplica a su
evolución causal con el terminal inferior fijado, sin presuponer un
semiflujo sobre el estado proyectado. Se especifica la familia de historias
o reinicios: dos estados proyectados iguales pueden
tener escapes diferentes si sus memorias difieren. La persistencia, la
atracción y la separación respecto de los equilibrios requieren controles
propios; una trayectoria larga o un exponente positivo de tiempo finito no
sustituye esos argumentos.

\subsection{Definición de los indicadores espectrales y de la prueba 0--1}
\label{bre:diag:definition}

Para una componente postransitoria $u_j$, muestreada uniformemente con
intervalo $\Delta t$, se sustrae su media $\bar u$ y se aplica una ventana
$w_j$. Con $N$ muestras, la transformada discreta es
\begin{equation}
 S_k=\sum_{j=0}^{N-1}w_j(u_j-\bar u)e^{-2\pi\ii jk/N},
 \qquad f_k=\frac{k}{N\Delta t}.
 \label{bre:diag:dft}
\end{equation}
La ventana de Hann es $w_j=[1-\cos(2\pi j/(N-1))]/2$.
La representación de módulo emplea $|S_k|/\sum_jw_j$; es distinta de
una densidad espectral de potencia. Para una PSD unilateral por Welch,
cada segmento usa el factor $\Delta t/\sum_jw_j^2$ sobre $|S_k|^2$;
se duplican las frecuencias positivas interiores, excepto Nyquist
cuando está presente, y se promedian los segmentos. Esta normalización
conserva las unidades de potencia por unidad de frecuencia.

Para cuantificar concentración espectral se excluye el modo cero y se
normalizan las potencias $P_k=|S_k|^2$ en los $m$ modos positivos:
\begin{equation}
 p_k=\frac{P_k}{\sum_{\ell=1}^mP_\ell},\qquad
 H=-\frac{\sum_{k=1}^mp_k\log p_k}{\log m}.
 \label{bre:diag:entropy}
\end{equation}
Se requiere $m>1$ y potencia total positiva; una señal constante se trata
por separado. Con $0\log0=0$, se obtiene $0\le H\le1$:
la cota inferior sigue de $p_k\le1$ y la superior de la concavidad del
logaritmo, que maximiza la entropía en $p_k=1/m$.
El cociente de potencia dominante es $\max_kp_k$, y la frecuencia
dominante corresponde al modo que lo alcanza. Para comparar ventanas se
conservan paso de muestreo, longitud y ventana; de otro modo cambian la
resolución $1/(N\Delta t)$ y el número de modos usados en $H$.
La deriva relativa usada en el clasificador compara los máximos de
ventanas parciales mediante $(f_{\max}-f_{\min})/f_{\max}$, cuando
$f_{\max}>0$ y todas las frecuencias parciales son positivas. Si alguna
es nula o no puede estimarse, el clasificador asigna deriva infinita y
no acepta el filtro de periodicidad. Una señal transitoria o una resolución insuficiente pueden
producir banda ancha sin demostrar caos.

La prueba 0--1 utiliza al menos $N=100$ muestras finitas de una señal
preparada $\phi_j$, después de
declarar el descarte, submuestreo, eliminación de tendencia y
normalización \cite{gottwald2009}. Para $c\in(0,\pi)$, define
\begin{align}
 p_c(n)&=\sum_{j=1}^n\phi_j\cos(jc),\nonumber\\
 q_c(n)&=\sum_{j=1}^n\phi_j\sin(jc).
 \label{bre:diag:translations}
\end{align}
El desplazamiento cuadrático medio a un retardo $n<N$ es
\begin{equation}
 M_c(n)=\frac1{N-n}\sum_{j=1}^{N-n}
 \bigl[\Delta_np_c(j)^2+\Delta_nq_c(j)^2\bigr],
 \label{bre:diag:msd}
\end{equation}
donde $\Delta_np_c(j)=p_c(j+n)-p_c(j)$, y análogamente para $q_c$.
Se elimina el término oscilatorio de la media mediante
\begin{equation}
 D_c(n)=M_c(n)-\bar\phi^{\,2}
                 \frac{1-\cos(nc)}{1-\cos c}.
 \label{bre:diag:modified-msd}
\end{equation}
La corrección procede de sumar la progresión geométrica
$\sum_{j=1}^n e^{\ii jc}$ y calcular su módulo al cuadrado:
$|\sum_{j=1}^n e^{\ii jc}|^2=(1-\cos nc)/(1-\cos c)$.
Por tanto, resta exactamente la contribución de una señal constante.

Para el vector de retardos $n=1,\ldots,n_{\max}$ se calcula
$K_c=\operatorname{corr}(n,D_c(n))$. El valor comunicado es la mediana
de $K_c$ sobre los valores de $c$ elegidos; se conserva esa selección al
comparar ensayos. En el procedimiento de correlación utilizado,
$n_{\max}=\min(\max(20,\lfloor N/10\rfloor),500)$.
Una varianza nula del desplazamiento no admite correlación de Pearson;
se trata como caso degenerado regular. Los umbrales operativos
$K<0.2$, $K>0.8$ y el intervalo intermedio describen, respectivamente,
compatibilidad regular, compatibilidad caótica e indeterminación de ese
diagnóstico finito. Un valor ligeramente negativo es posible por la
correlación muestral. Ruido, transitorios y sobremuestreo afectan el
resultado; la mediana y la comparación de varios intervalos de muestreo
reducen sensibilidades concretas, sin demostrar caos ni ocultedad.

\section{Aplicaciones y resultados numéricos}
\label{bre:res:section}

Las aplicaciones distinguen la obtención de una raíz algebraica, la
construcción de una condición inicial y la clasificación de la trayectoria
resultante. El cierre armónico se evalúa mediante los residuos de la
sección~\ref{bre:bal:sec}; las integraciones y los contactos se interpretan
con los criterios de la sección~\ref{bre:num:seccion}. En las comparaciones
fraccionarias, cada reinicio tiene terminal inferior \(t_0=0\) y memoria
inicial nueva, salvo la continuación expresamente descrita con historia
heredada. Las frecuencias espectrales se expresan en ciclos por unidad de
tiempo; las frecuencias \(\omega\) del balance son angulares.
Las raíces y magnitudes calculadas se muestran redondeadas; las identidades
exactas corresponden a sus expresiones algebraicas. Los parámetros
decimales declarados fijan el problema de cada ensayo.

\subsection{Chua no suave de orden entero}
\label{bre:res:chua-integer}

Se considera el campo de Chua de \eqref{bre:bal:chua}, con
\[
 \begin{aligned}
 (\alpha,\beta,\gamma)&=(8.4562,12.0732,0.0052),\\
 (m_0,m_1)&=(-0.1768,-1.1468),\qquad q=1,
 \end{aligned}
\]
y característica \(h(x)=m_1x+(m_0-m_1)\operatorname{sat}(x)\)
\cite{LeonovKuznetsov2013}. Para \(\beta+\gamma\ne0\), las ecuaciones de
equilibrio segunda y tercera se resuelven como
\begin{equation}
 y=\frac{\gamma x}{\beta+\gamma},\qquad
 z=-\frac{\beta x}{\beta+\gamma},\qquad
 h(x)+\frac{\beta x}{\beta+\gamma}=0.
 \label{bre:res:chua-equilibrium-reduction}
\end{equation}
La última identidad se obtiene sustituyendo las primeras en
\(y-x-h(x)=0\). Sean
\(\kappa=m_1+\beta/(\beta+\gamma)\) y \(d=m_0-m_1\).
En \(|x|\leq1\), queda \((\kappa+d)x=0\); como
\(\kappa+d>0\), sólo se obtiene el origen. En \(|x|>1\),
\(\kappa x+d\operatorname{sign}(x)=0\), por lo que
\[
 x=-\frac d\kappa\operatorname{sign}(x).
\]
Aquí \(\kappa<0\) y \(-d/\kappa>1\): hay exactamente dos raíces externas,
una por signo, cuyas otras coordenadas da
\eqref{bre:res:chua-equilibrium-reduction}. Ninguna está en un quiebre.
Para una característica diferenciable en el equilibrio, con pendiente
\(m=h'(x_e)\), el Jacobiano y su polinomio son
\begin{equation}
 \begin{gathered}
 J_e=\begin{pmatrix}
 -\alpha(1+m)&\alpha&0\\1&-1&1\\0&-\beta&-\gamma
 \end{pmatrix},\\
 p_e(s)=[s+\alpha(1+m)]D(s)-\alpha(s+\gamma),\\
 D(s)=s^2+(1+\gamma)s+\beta+\gamma.
 \end{gathered}
 \label{bre:res:chua-jacobian}
\end{equation}
La expansión del determinante sigue \eqref{bre:bal:characteristic}; se usa
\(m_0\) en el origen y \(m_1\) en los equilibrios externos. La misma
matriz permite aplicar Matignon al cambiar el orden, sin cambiar los
equilibrios del campo.

La raíz de fase, la ganancia y la amplitud
calculadas por \eqref{bre:bal:phase-polynomial} y
\eqref{bre:bal:saturation-df} son
\[
 \begin{aligned}
 \omega_0&=2.039186939959,\\
 k&=0.209867354515,\qquad A_0=5.856145086252.
 \end{aligned}
\]
La reconstrucción modal de \eqref{bre:bal:seed} produce
\begin{equation}
 X_0=\begin{pmatrix}
 5.856145086252\\0.369331578246\\-8.366536168324
 \end{pmatrix}.
 \label{bre:res:chua-integer-seed}
\end{equation}
Los residuos del cierre complejo y de la función descriptiva fueron,
respectivamente, \(3.26\times10^{-15}\) y \(1.98\times10^{-13}\). Estas
cifras verifican las ecuaciones usadas para generar la semilla; la
continuación y los diagnósticos posteriores determinan su destino.
Una búsqueda independiente por recurrencias, con 512 condiciones iniciales,
identificó dos conjuntos recurrentes y 22 contactos. La fracción
\(22/512=0.04296875\) se refiere al dominio y a la muestra elegidos.

La comparación de vecindades empleó las mismas 504 coordenadas en dos
integraciones independientes: tres equilibrios, siete radios entre
\(10^{-5}\) y \(10^{-2}\), y 24 puntos uniformes en cada bola. El horizonte
fue \(T=500\), con descarte inicial de 120; las tolerancias de clasificación
fueron \(10^{-3}\) para equilibrio y \(0.08\) para la distancia normalizada
a la nube. Coincidieron las 504 clases: 152 trayectorias hacia el origen,
198 acotadas no triviales y 154 divergentes según el detector. Ninguna
alcanzó la nube candidata. Los espectros de Lyapunov finitos fueron
\[
 \begin{aligned}
 \widehat\lambda^{(1)}&=(0.14018,0.01319,-1.20304),\\
 \widehat\lambda^{(2)}&=(0.14591,-0.00052,-1.17115).
 \end{aligned}
\]
El exponente máximo positivo y la geometría de la
Fig.~\ref{bre:res:fig-chua-integer} respaldan un candidato caótico bajo esas
ventanas. La ausencia de contactos sólo establece compatibilidad con
ocultedad en el sondeo finito.

El control EFORK--3 de referencia utilizó un protocolo diferente:
\(T=320\), descarte de 80, otras 504 coordenadas y un clasificador por huellas
de sección. Produjo 369 convergencias a equilibrio, 69 divergencias y 66
trayectorias hacia otro destino, también sin contactos. Estos conteos no se
suman a los anteriores ni se interpretan como repeticiones del mismo ensayo.

\begin{figure*}[t]
 \centering
 \includegraphics[width=0.72\textwidth]{chua_entero.png}
 \caption{Trayectoria candidata del Chua no suave entero, localizada mediante
 semilla armónica y continuación. La representación corresponde a una ventana
 postransitoria; la clasificación dinámica utiliza además exponentes y sondeos.}
 \label{bre:res:fig-chua-integer}
\end{figure*}

\subsection{Chua de Caputo no suave: continuación de una semilla sesgada}
\label{bre:res:chua-pwl}

Se mantienen los cinco parámetros anteriores y se fija \(q=0.9998\).
Como control, iniciar directamente el PVI de Caputo desde
\eqref{bre:res:chua-integer-seed}, con ABM de memoria completa,
\(h=0.01\), \(T=500\) y descarte de 250, produjo una respuesta acotada
con frecuencia dominante \(0.3679852806\). Las medianas de la prueba 0--1
fueron
\[
 \begin{aligned}
 K_x&=0.0736852160,\qquad K_y=0.0214448125,\\
 K_z&=-0.0012988924.
 \end{aligned}
\]
Es una reconstrucción numérica relacionada con el caso de
Danca~\cite{Danca2017}; la condición inicial se obtuvo por balance y no se
atribuye a la fuente como dato publicado.

La localización sesgada resolvió el cierre de
\eqref{bre:bal:bias-closure}. Se ensayaron 2430 inicios y se seleccionó,
en el par nominal de pendientes,
\[
 \begin{aligned}
 A&=4.577882721023,&c&=2.776397003274,\\
 \omega&=2.040286051079,&N_1&=0.210022792962,\\
 \Psi_0&=0.408770376181.
 \end{aligned}
\]
La integración trapezoidal periódica de los coeficientes utilizó 8192 nodos
y dio un residuo de \(4.55\times10^{-16}\). Al reevaluar las integrales
continuas con cuadratura adaptativa y separar los quiebres de la saturación,
el residuo fue \(1.9\times10^{-7}\), menor que el umbral \(10^{-4}\). La
segunda cifra controla la aproximación de las integrales, no sólo la
resolución del sistema discretizado. La semilla y el extremo obtenido son
\begin{equation}
 \begin{aligned}
 X_{\rm sem}&=\begin{pmatrix}
 7.354279724297\\0.290968778784\\-9.316054818666
 \end{pmatrix},\\
 X_{\rm end}&=\begin{pmatrix}
 0.822310709850\\1.465272703995\\1.292308250248
 \end{pmatrix}.
 \end{aligned}
 \label{bre:res:pwl-seeds}
\end{equation}
La homotopía \eqref{bre:bal:homotopy} se integró en 21 niveles,
\(\varepsilon_k=0.05k\), con \(h=0.01\) y 60 unidades por nivel. Cada nivel
conservó 30 unidades después de su tramo inicial de ajuste. Las sumas de
Caputo mantuvieron la historia desde el primer nivel, como exige
\eqref{bre:num:continuacion-integral}; el tiempo total fue \(21\cdot60=1260\).
La Fig.~\ref{bre:res:fig-pwl-continuation} muestra ese transporte.

Desde \(X_{\rm end}\) se inició un PVI autónomo nuevo, con \(h=0.01\),
\(T=300\) y descarte de 100. Su frecuencia dominante fue
\(0.3699815009\), y sus entropías espectrales
\((0.11526785,0.10497088,0.10339964)\). La respuesta tiene apariencia
periódica en la ventana retenida, como muestra la
Fig.~\ref{bre:res:fig-pwl-final}; no se afirma una solución periódica exacta
del PVI de Caputo con terminal finito. Las 36 perturbaciones de
\(X_{\rm end}\), doce para cada radio \(10^{-6},10^{-4},10^{-2}\),
alcanzaron la nube objetivo con \(h=0.01\) y memoria completa. Este resultado
mide sensibilidad al dato inicial a paso fijo. El refinamiento
\(h,h/2,h/4\) de ese banco sigue pendiente.

\begin{figure*}[t]
 \centering
 \includegraphics[width=0.96\textwidth]{continuacion_caputo.pdf}
 \caption{Continuación causal del Chua no suave, \(q=0.9998\).
 Los 21 niveles de la homotopía comparten una sola historia. El estado final
 se utiliza después para definir un PVI autónomo con memoria reiniciada.}
 \label{bre:res:fig-pwl-continuation}
\end{figure*}
\begin{figure*}[t]
 \centering
 \includegraphics[width=0.96\textwidth]{chua_caputo_pwl.pdf}
 \caption{Respuesta postransitoria desde \(X_{\rm end}\),
 \(q=0.9998\), \(h=0.01\). La regularidad observada y los contactos desde
 perturbaciones son resultados finitos; no establecen periodicidad exacta.}
 \label{bre:res:fig-pwl-final}
\end{figure*}

El sondeo de vecindades utilizó bolas uniformes, con la construcción de la
sección~\ref{bre:num:sondeos}. Cada sonda se integró con \(h=0.01\),
\(T=200\), descarte de 50 y umbral de contacto \(D_{90}\leq0.5\).
El detector de equilibrio requería, desde \(t=25\), distancia menor que
\(10^{-3}\), norma del campo menor que \(10^{-4}\) y 50 pasos consecutivos;
al final del horizonte se admitía distancia a equilibrio no mayor que 0.5.
El detector de divergencia usó cinco normas consecutivas mayores que 80 o
cinco razones consecutivas de crecimiento mayores que 1.25, con corte duro
120. Son criterios numéricos finitos, distintos de probar convergencia o
explosión asintótica.

Se ejecutaron 17\,400 trayectorias: 7200 en radios \(r\leq0.01\) y
10\,200 en \(r=0.03,0.1,0.3\). No hubo contactos en las 7200 sondas locales
ni en las 6600 correspondientes a \(r=0.03,0.1\). En \(r=0.3\), entre
3600 sondas, se observaron 37 contactos: 19 desde el equilibrio positivo y
18 desde el negativo. Sus distancias iniciales reales al equilibrio
estuvieron entre \(0.1941805298\) y \(0.2997254735\). Por tanto, el primer
radio de bola muestreado con acceso fue 0.3; esta observación no localiza
una frontera continua de cuenca. Todos los conteos corresponden al mismo
convenio de reinicio.

\subsection{Chua con arctangente: control de Wu y candidato c590}
\label{bre:res:arctan}

La característica se fija como
\begin{equation}
 h(x)=a_1x+a_2\arctan(\rho x),\qquad
 \Caputo qx=\alpha[y-x-h(x)].
 \label{bre:res:arctan-convention}
\end{equation}
Así, el coeficiente lineal de \(x\) dentro del corchete es \(-(1+a_1)\).
En el control de Wu, \(a_1=0.4\); el número 1.4 corresponde a \(1+a_1\).
Los restantes parámetros son
\[
 \begin{aligned}
 q&=0.99,&(\alpha,\beta,\gamma)&=(8.4562,12.0732,0.0052),\\
 a_2&=-1.5585,&\rho&=1.
 \end{aligned}
\]
La recurrencia ADM local de orden cuatro, \(h=0.01\), \(T=100\) y descarte
de 50, desde \(X_{0,\pm}=\pm(13.8,0.7093,-19.8768)^{\T}\), produjo dos
respuestas simétricas acotadas, con frecuencia dominante \(0.5398920216\)
y entropías \((0.11607348,0.11106582,0.11095067)\)
\cite{Wu2023ArctanChua}. Este control reconstruye la recurrencia local;
la sección~\ref{bre:num:adm-local} demuestra por qué no se identifica con
la integración causal de memoria completa.

El candidato c590 utiliza la misma familia de campos, con
\begin{equation}
 \begin{aligned}
 q&=0.9999,&\alpha&=21.8493569066,\\
 \beta&=19.0818408409,&\gamma&=0.0073780120,\\
 a_1&=0.0422897934,&a_2&=-3.3367815123,\\
 \rho&=1.7984259333.
 \end{aligned}
 \label{bre:res:c590-parameters}
\end{equation}
La exploración inicial en orden entero evaluó 2400 casos globales y 1000
casos locales; el espectro variacional de la selección entera fue
\((0.4699043683,0.0010414594,-19.0491693716)\). Su semilla
\((7.6733768929,0.5079327671,-14.4903935267)^{\T}\) se utilizó para evaluar
por separado los órdenes \(0.9995,0.9998,0.9999,0.99995\), todos mediante
ABM de memoria completa, \(h=0.005\), \(T=300\) y descarte de 150.
No se transportó una historia entre órdenes distintos.

De la trayectoria con \(q=0.9999\) se extrajeron 16 estados entre
\(t\simeq150\) y 300. Cada estado definió un reinicio independiente. La
criba corta usó \(h=0.0025\), \(T=180\) y descarte de 90; el control largo
comparó \(h=0.0025,0.005,0.01\), con \(T=300\) y descarte de 150.
La selección, extraída cerca de \(t=240\), fue
\begin{equation}
 X_0^{\rm frac}=\begin{pmatrix}
 5.8642449791\\1.5847111486\\3.2155806478
 \end{pmatrix}.
 \label{bre:res:c590-seed}
\end{equation}
Permaneció acotada en los tres pasos; las medianas remuestreadas de la
prueba 0--1 fueron
\[
 (K_{0.0025},K_{0.005},K_{0.01})=(0.865843,0.885089,0.740058).
\]
Los dos pasos finos superaron el umbral de selección 0.8. El valor del paso
mayor muestra sensibilidad numérica y evita atribuir a ese control una
convergencia monótona no observada.

Las ecuaciones estacionarias de las componentes segunda y tercera dan
\(y=\gamma x/(\beta+\gamma)\) y
\(z=-\beta x/(\beta+\gamma)\). La primera se reduce a
\[
 \left(a_1+\frac{\beta}{\beta+\gamma}\right)x
       +a_2\arctan(\rho x)=0.
\]
Esta reducción también permite contar las raíces. Escriba
\(g(x)=\kappa_a x+a_2\arctan(\rho x)\), con
\(\kappa_a=a_1+\beta/(\beta+\gamma)>0\). La función es impar, y
\[
 \begin{aligned}
 g'(0)&=\kappa_a+a_2\rho<0,\\
 g''(x)&=\frac{-2a_2\rho^3x}{(1+\rho^2x^2)^2}>0\quad(x>0).
 \end{aligned}
\]
Además, \(g'(x)\to\kappa_a>0\) y \(g(x)\to+\infty\) cuando
\(x\to+\infty\). La derivada estrictamente creciente tiene un único cero
positivo; la función decrece desde \(g(0)=0\), alcanza un mínimo negativo
y luego crece hasta infinito. Existe exactamente una raíz positiva; por
imparidad, exactamente una negativa, además del origen. Sus valores producen
\[
 E_0=0,\qquad E_\pm=\pm\begin{pmatrix}
 4.6494055135\\0.0017970023\\-4.6476085112
 \end{pmatrix}.
\]
El origen es inestable. El margen angular de Matignon de los equilibrios
externos es \(7.42265\times10^{-4}>0\), por lo que son localmente estables
para el orden fijado. Su coexistencia con una trayectoria candidata exige
estudiar cuencas; la estabilidad local no decide el destino de
\eqref{bre:res:c590-seed}.

Sobre 30\,000 estados postransitorios, ninguna componente cumplió el filtro
de apariencia periódica. Las entropías espectrales fueron
\((0.394751,0.261922,0.315229)\); las fracciones de potencia dominante,
\((0.122163,0.290492,0.241937)\); y las derivas relativas de frecuencia,
\((0.941176,0.029412,0.029412)\). Las
Figs.~\ref{bre:res:fig-c590} y~\ref{bre:res:fig-c590-diagnostics} acompañan
estos indicadores. La conclusión es un candidato aperiódico, con indicios
numéricos de caos; los exponentes enteros no se transfieren automáticamente
al sistema de Caputo.

\begin{figure*}[t]
 \centering
 \includegraphics[width=0.96\textwidth]{chua_caputo_arctan.pdf}
 \caption{Trayectoria postransitoria de c590, \(q=0.9999\), \(h=0.005\),
 desde \eqref{bre:res:c590-seed}. El intervalo de observación es finito.}
 \label{bre:res:fig-c590}
\end{figure*}
\begin{figure*}[t]
 \centering
 \includegraphics[width=0.96\textwidth]{arctan_diagnosticos.pdf}
 \caption{Trayectoria postransitoria de c590 y reconstrucción de su
 primer armónico, con frecuencia \(f_1=0.67451\). La diferencia geométrica
 muestra la parte de la respuesta que no describe ese armónico. La
 clasificación dinámica utiliza además el control de paso y la prueba 0--1.}
 \label{bre:res:fig-c590-diagnostics}
\end{figure*}

El sondeo de cuenca utilizó superficies esféricas, con seis semiejes y
direcciones de Fibonacci, en lugar de bolas uniformes. Cada PVI usó
\(h=0.005\), \(T=180\) y descarte de 90. La referencia incluyó la nube y
su inversión central. La calibración dio
\(d_{90}^{\rm cal}=0.2950508245\) y umbral
\(\delta=0.5901016489\). La parada por equilibrio se activó desde
\(t=45\), con distancia \(<10^{-6}\), norma del campo \(<10^{-4}\) y
50 pasos consecutivos; el corte de divergencia fue \(\norm{X}_2=120\).
Las 8400 sondas con \(r\leq0.3\) dieron cero contactos: 8396 completaron el
horizonte y cuatro terminaron anticipadamente por el criterio de equilibrio.
En \(r=1\) hubo 22 contactos de 2400, todos desde \(E_0\); en \(r=2\),
588 de 2700, distribuidos como 289, 150 y 149 desde \(E_0,E_+,E_-\).
Los contactos macroscópicos describen acceso a la nube bajo el muestreo
superficial. La ausencia local de contactos no excluye todas las condiciones
de una vecindad abierta.

\subsection{MAVPD entero: ciclos, candidato caótico y puntos perpetuos}
\label{bre:res:mavpd}

El campo MAVPD se definió en \eqref{bre:geo:mavpd}. Se fijan
\(\delta=100\), \(\rho=200\), mientras \(\xi\) y \(\gamma\) seleccionan el
régimen \cite{matouk2023}. La parte lineal de Lur'e tiene
\(b=-\delta e_1\), \(c=e_1\) y \(\psi(u)=u^3\). La transferencia en la
convención \(W=c^{\T}(sI-P)^{-1}b\) es
\begin{equation}
 \begin{aligned}
 W(s)&=-\frac{\delta D(s)}{Q(s)},\qquad D(s)=s^2+\xi s+\rho,\\
 Q(s)&=(s-\delta\gamma)D(s)-\delta s.
 \end{aligned}
 \label{bre:res:mavpd-transfer}
\end{equation}
Como \(\cos^3\theta=(3\cos\theta+\cos3\theta)/4\), el primer armónico de
\((A\cos\theta)^3\) tiene amplitud \(3A^3/4\), y
\(N(A)=3A^2/4\). La condición de fase, después de eliminar el denominador
no nulo y poner \(z=\omega^2\), queda
\begin{equation}
 z^2+(\delta-2\rho+\xi^2)z+\rho(\rho-\delta)=0.
 \label{bre:res:mavpd-frequency}
\end{equation}
Para precisar la ganancia, el polinomio de \(P+kbc^{\T}\) se obtiene
reemplazando \(\gamma\) por \(\gamma-k\). Su parte real en
\(s=\ii\omega\) impone
\(\delta(k-\gamma)(\rho-z)=\xi z\); por tanto,
\[
 k=\gamma+\frac{\xi z}{\delta(\rho-z)},\qquad
 A=\sqrt{4k/3},
\]
para una raíz con \(z>0\), \(z\ne\rho\), \(k>0\). El modo normalizado es
\[
 v=\begin{pmatrix}
 1\\ k-\gamma+\ii\omega/\delta\\
 \rho/\delta-\ii\rho(k-\gamma)/\omega
 \end{pmatrix},\qquad Y_0(\theta)=A\operatorname{Re}(v e^{\ii\theta}).
\]
Así se fija la semilla sin escoger sus componentes por separado.

\subsubsection{Régimen periódico y comparación con PP}
Para \((\xi,\gamma)=(3.1,0.1)\), las ramas son
\[
 \begin{array}{c|rr}
 &\text{rama 0}&\text{rama 1}\\\hline
 \omega&10.5975230311&13.3447557342\\
 k&0.1397015948&0.3518790503\\
 A&0.4315886851&0.6849613618
 \end{array}
\]
y las semillas de fase cero son
\[
 Y_{0,0}=\begin{pmatrix}0.4315887\\0.0171348\\0.8631774\end{pmatrix},
 \qquad Y_{0,1}=\begin{pmatrix}0.6849614\\0.1725274\\1.3699227\end{pmatrix}.
\]
Se integraron ambas ramas con fases
\(0,\pi\) y 21 niveles de continuación. Cada nivel utilizó RK4 con
\(h=0.002\), transitorio de 20 y cola de 10. La primera rama alcanzó dos
ciclos internos simétricos; la segunda, el conjunto externo de la
Fig.~\ref{bre:res:fig-mavpd-periodic}. Un refinamiento DOP853 empleó
tolerancias \((2\times10^{-11},2\times10^{-13})\), paso máximo 0.005,
transitorio de 600 y ventana de 1200; registró 2433 retornos positivos a
\(y_2=0\).

En otro control de sección, DOP853 con eventos densos registró 405
cruces; el rango de los últimos retornos fue del orden de \(10^{-13}\).
El periodo fue \(0.493177986749\), con desviación estándar
\(1.25\times10^{-14}\). Un cálculo independiente dio
\(0.493178150753\), y la distancia normalizada entre nubes externas fue
\(0.00061096936\). El espectro finito
\((-0.0017614,-0.3887396,-51.1886634)\) y la mediana 0--1
\(K=0.000440085726\) respaldan un ciclo estable de periodo uno. La
clasificación cuasiperiódica de la fuente no se reproduce para estas
ecuaciones y estos controles.

El sondeo común, con \(h=0.005\), transitorio de 180 y cola de 80, empleó
108 condiciones: tres equilibrios, radios \(10^{-5},10^{-3},10^{-2}\) y
12 direcciones por radio. Hubo 54 destinos en cada ciclo interno y cero
contactos con el externo, coincidentes en ambos cálculos. Para
\(\xi=3.5\), conservando \(\gamma=0.1\), el mismo banco dio 24 contactos
con cada uno de los dos ciclos internos. Este control demuestra la
capacidad del sondeo para detectar acceso; no convierte el radio mínimo
ensayado en una prueba sobre radios arbitrariamente pequeños.

Los PP exactos de \eqref{bre:geo:mavpd-roots}, para \(\xi=3.1\), son
\[
 p_\pm=\pm\begin{pmatrix}
 0.1825741858350554\\0.01217161238900369\\0.1448421874291439
 \end{pmatrix}.
\]
Satisfacen \(J_f(p_\pm)f(p_\pm)=0\) y
\(\norm{f(p_\pm)}=3.44265186329548>0\). El experimento DOP853 desde esos
puntos, \(T=120\), tolerancias \(10^{-10}/10^{-12}\), alcanzó los ciclos
simétricos. En \(t\geq80\), los rangos de las componentes fueron
\((0.37943,0.05722,0.97545)\). El detector registró 74 y 75 máximos,
periodos medios \(0.53760,0.53730\) y coeficientes de variación
\(0.0361,0.0223\). Estos indicadores pertenecen al piloto desde PP y no al
periodo del ciclo externo. La Fig.~\ref{bre:res:fig-mavpd-pp} muestra que
la cancelación inicial de \(J_ff\) no se conserva a lo largo de la órbita.

\begin{figure*}[t]
 \centering
 \includegraphics[width=0.96\textwidth]{mavpd_periodico.png}
 \caption{MAVPD periódico, \(\xi=3.1\). Arriba: proyecciones del ciclo
 externo en los planos \(y_1y_2\) y \(y_1y_3\). Abajo: clasificación de las
 108 sondas, con 54 destinos en cada ciclo interno y ninguno en el externo,
 y continuación de la rama 1. Los paneles describen experimentos distintos
 dentro del mismo caso.}
 \label{bre:res:fig-mavpd-periodic}
\end{figure*}
\begin{figure*}[t]
 \centering
 \includegraphics[width=0.96\textwidth]{mavpd_pp.png}
 \caption{MAVPD entero desde los dos PP exactos. Las cruces indican las
 semillas; el diagnóstico de aceleración distingue su anulación inicial
 de la evolución posterior.}
 \label{bre:res:fig-mavpd-pp}
\end{figure*}

\subsubsection{Continuación hacia un candidato caótico}
Las dos ramas en \((\xi,\gamma)=(3.1,0.1)\) no superaron el cribado
finito \(\widehat\lambda_1>0.02\). Se transportó entonces la rama de menor
frecuencia desde \(\xi=3.10\) hasta 2.85 en pasos de \(-0.01\), y luego
se varió \(\gamma\) respecto de una frontera espectral calculada.
En \(E_\pm=(\pm\sqrt\gamma,0,\pm\sqrt\gamma)\), el polinomio del
Jacobiano es \(\lambda^3+b_1\lambda^2+b_2\lambda+b_3\), con
\[
 b_1=\xi+2\delta\gamma,\quad
 b_2=\rho-\delta+2\delta\gamma\xi,\quad
 b_3=2\delta\gamma\rho.
\]
Para una raíz \(\lambda=\ii\Omega\), \(\Omega>0\), sus partes real e
imaginaria imponen \(b_3-b_1\Omega^2=0\) y
\(\Omega(b_2-\Omega^2)=0\). Por tanto,
\(\Omega^2=b_2>0\) y \(b_1b_2=b_3\). Con \(g=2\delta\gamma\),
\begin{equation}
 \begin{split}
 b_1b_2-b_3
 &=(\xi+g)(\rho-\delta+\xi g)-\rho g\\
 &=\xi g^2+(\xi^2-\delta)g+\xi(\rho-\delta).
 \end{split}
 \label{bre:res:mavpd-hopf}
\end{equation}
En \(\xi=2.85\), la raíz alta da
\(\gamma_H=0.14380379839949112\). Es una frontera candidata de Hopf;
no se ha demostrado aquí transversalidad ni no degeneración de la
bifurcación. El cribado de desplazamientos seleccionó
\[
 \gamma=\gamma_H+0.010=0.15380379839949113.
\]
Se seleccionó el mayor desplazamiento ensayado con exponente válido
\(\widehat\lambda_1>0.5\), cero contactos preliminares, cero ambigüedades y
cero fallos. El valor \(\xi=2.85\) fue fijado como extremo del transporte,
no obtenido por una optimización conjunta.

La trayectoria estricta usó DOP853, \(T=900\), descarte de 300, muestreo
0.02, tolerancias \((2\times10^{-11},2\times10^{-13})\) y paso máximo 0.01.
El cálculo variacional separado usó descarte de 300, acumulación de 1200,
QR cada 0.5 y tolerancias \((2\times10^{-12},2\times10^{-14})\). Se obtuvo
\[
 \begin{aligned}
 \widehat\lambda_1&=0.7133503985,\quad
 \widehat\lambda_2=-0.0009058833,\\
 \widehat\lambda_3&=-20.9236249563,\quad D_{KY}=2.0340497651.
 \end{aligned}
\]
La estimación de Kaplan--Yorke usa
\(D_{KY}=2+(\widehat\lambda_1+\widehat\lambda_2)/|\widehat\lambda_3|\),
pues las dos primeras suman positivamente y la suma total es negativa.
Un control que perturbó el estado y modificó simultáneamente tolerancias,
paso máximo y ventanas dio \(\widehat\lambda_1=0.6830338121\).
La integración independiente produjo
\((0.6938035753,0.0017936287,-20.9005938240)\).
Los exponentes máximos positivos persisten en estos controles finitos.

La sección \(y_2=0\), con orientación positiva, retuvo 1174 cruces mediante
interpolación lineal entre muestras. Las medianas 0--1 de dos coordenadas
de retorno fueron \(0.9985162431\) y \(0.9984323397\). Sobre el flujo,
el resultado fue sensible al submuestreo: tres de diez pasos apoyaron la
clase caótica, uno la regular y seis fueron inconclusos. Por eso la
interpretación conjunta se apoya en retornos y exponentes, con su error
finito, y no en una sola prueba 0--1. La
Fig.~\ref{bre:res:fig-mavpd-chaos} reúne geometría y evolución variacional.
Un control EFORK con \(h=0.002\) alcanzó la misma nube calibrada.

Los 108 sondeos principales y cuatro dirigidos sobre ambos signos del
autovector inestable del origen, en radios \(10^{-7}\) y \(10^{-4}\),
sumaron 112 trayectorias. Todas terminaron en un equilibrio según el
clasificador; no hubo contactos con la nube candidata, ambigüedades ni
fallos numéricos. Los sondeos principales usaron radios
\(10^{-5},10^{-3},10^{-2}\), doce direcciones por equilibrio y radio,
transitorio de 400 y cola de 100. El resultado constituye evidencia
numérica de un candidato caótico compatible con ocultedad en ese banco;
la separación de vecindades continuas permanece sin certificar.

\begin{figure*}[t]
 \centering
 \includegraphics[width=0.96\textwidth]{mavpd_fases.pdf}\\[3pt]
 \includegraphics[width=0.76\textwidth]{mavpd_lyapunov.pdf}
 \caption{MAVPD, \(\xi=2.85\), \(\gamma=0.15380379839949113\).
 Arriba: proyecciones de la trayectoria postransitoria. Abajo: evolución de
 los exponentes variacionales finitos. Los controles de paso, integrador,
 retornos y condiciones iniciales delimitan la clasificación numérica.}
 \label{bre:res:fig-mavpd-chaos}
\end{figure*}

\subsection{Kalman--Fitts: obstrucción de ganancia y mapa de conmutación}
\label{bre:res:kalman}

El sistema de cuarto orden es
\begin{equation}
 \begin{gathered}
 \dot X=AX+b\tanh(c^{\T}X/\varepsilon),\\
 b=e_4,\quad c=-e_3,\quad\varepsilon=0.01,
 \end{gathered}
 \label{bre:res:kalman-field}
\end{equation}
con \(A\) compañera, última fila \((-a_0,-a_1,-a_2,-a_3)\), y polinomio
\[
 \begin{aligned}
 p_A(s)&=[(s+\beta)^2+m_1^2][(s+\beta)^2+m_2^2],\\
 (m_1,m_2,\beta)&=(0.9,1.1,0.03),\\
 (a_0,a_1,a_2,a_3)&=(0.98191881,0.121308,2.0254,0.12).
 \end{aligned}
\]
Las tres primeras filas de \(A\) desplazan consecutivamente las
coordenadas. La condición de equilibrio obliga a \(x_2=x_3=x_4=0\), y
la última ecuación da \(x_1=0\) porque \(a_0>0\)
\cite{kalman2019}. El Jacobiano en el origen tiene polinomio
\(s^4+a_3s^3+(a_2+1/\varepsilon)s^2+a_1s+a_0\). Sus coeficientes son
positivos, y las otras dos desigualdades de Routh--Hurwitz dan
\[
 \begin{aligned}
 a_3(a_2+1/\varepsilon)-a_1&=12.12174>0,\\
 a_3(a_2+1/\varepsilon)a_1-a_1^2-a_3^2a_0
 &=1.456324405056>0.
 \end{aligned}
\]
El único equilibrio es localmente asintóticamente estable.

Al resolver \((sI-A)v=b\), las tres primeras filas imponen
\(v_j=s^{j-1}v_1\), y la cuarta da \(p_A(s)v_1=1\). De aquí
\(W(s)=-s^2/p_A(s)\). La condición de fase da
\(\omega_0^2=a_1/a_3\), y la ganancia requerida es
\[
 \begin{aligned}
 \omega_0&=1.005435229142,\\
 k&=\frac{\omega_0^4-a_2\omega_0^2+a_0}{\omega_0^2}
   =-0.043168701157.
 \end{aligned}
\]
Para \(u=A\cos\theta/\varepsilon\), la función descriptiva es
\[
 N(A)=\frac1{\pi\varepsilon}\int_0^{2\pi}
         \frac{\tanh u}{u}\cos^2\theta\,\dd\theta.
\]
La extensión continua en \(u=0\) vale uno, y
\(0<\tanh u/u\leq1\); por ello \(N(A)>0\) para \(A>0\).
No existe amplitud que satisfaga la ganancia negativa. Esta obstrucción
justifica cambiar de localizador.

El precursor reemplaza \(\tanh\) por \(\operatorname{sign}\).
En un semiespacio con signo fijo \(\sigma\in\{-1,1\}\), la solución exacta
es
\begin{equation}
 \Phi_\sigma(t,X)=e^{At}X+A^{-1}(e^{At}-I)b\sigma.
 \label{bre:res:kalman-flow}
\end{equation}
Se obtiene integrando la fuerza constante:
\(\int_0^t e^{A(t-s)}\,\dd s=A^{-1}(e^{At}-I)\), con \(A\) invertible.
En \(\Sigma=\{c^{\T}X=0\}\), se elige el signo de salida y se busca el
primer cero transversal de \(c^{\T}\Phi_\sigma(t,X)\) después de abandonar
la sección. Si \(t_+(X)\) es ese tiempo, el mapa es
\(\Pi(X)=\Phi_\sigma(t_+(X),X)\). La transversalidad permite distinguir
el retorno de un contacto tangencial. El precursor convergió en 135
iteraciones a un retorno de periodo dos, \(\Pi^2(X)=X\), con semilla
\[
 X_{\rm sw}=\begin{pmatrix}
 0.212347571803\\-1.671123727247\\0\\1.478087383043
 \end{pmatrix}.
\]
Se continuó la no linealidad
\((1-\lambda)\operatorname{sign}\sigma+
\lambda\tanh(\sigma/\varepsilon)\) en 51 nodos
\(\lambda=0,0.02,\ldots,1\). Cada nodo utilizó RK4, \(h=0.01\),
transitorio de 100 y ventana retenida de 20. La corrida final descartó
300 unidades y conservó otras 300.

El ciclo resultante tuvo periodo \(11.9169960234\), frente a
\(11.9169379538\) en una integración independiente. Su espectro finito
fue
\[
 \widehat\lambda=\begin{pmatrix}
 -0.00283635\\-0.02423392\\-0.03056980\\-0.03988642
 \end{pmatrix},
\]
compatible con un exponente tangencial próximo a cero y tres negativos.
Los controles combinaron DOP853 con dos ajustes de tolerancias y RK4 con
\(h=0.01,0.005\), en horizontes 300 y 600. En las ocho configuraciones,
los errores relativos de periodo estuvieron entre \(4.87\times10^{-6}\)
y \(6.11\times10^{-6}\), y las distancias normalizadas entre nubes,
entre \(2.03\times10^{-4}\) y \(9.24\times10^{-4}\).

De las 48 condiciones alrededor del origen, doce en cada radio
\(10^{-5},10^{-4},10^{-3},10^{-2}\), cuarenta se clasificaron en el
equilibrio y ocho como acotadas no triviales; ninguna contactó el ciclo.
Ambos cálculos coincidieron punto por punto. La evidencia favorece un
ciclo estable compatible con ocultedad bajo ese sondeo. Una demostración
requiere además un dominio de retorno con atracción certificada y una
vecindad explícita del origen contenida en su cuenca local.

\subsection{PLL: retorno sobre el cilindro desde una solución exacta}
\label{bre:res:pll}

El PLL de \eqref{bre:geo:pll} tiene
\(\tau_1=0.0448\), \(\tau_2=0.0185\), \(L=500\) y
\(\omega_\Delta=178.9\) \cite{pll2015}. Para separar la constante del
filtro de los horizontes de integración, se escribe
\(T_f=\tau_1+\tau_2=0.0633\). Sus ecuaciones son
\begin{equation}
 \begin{aligned}
 \dot x&=-\frac{x}{T_f}+\frac{\tau_1}{2T_f}\sin\theta,\\
 \dot\theta&=\omega_\Delta-\frac L{T_f}x
                    -\frac{\tau_2L}{2T_f}\sin\theta.
 \end{aligned}
 \label{bre:res:pll-field}
\end{equation}
El estado pertenece a \(\R\times S^1\). Las condiciones estacionarias
imponen \(x_e=(\tau_1/2)\sin\theta_e\) y
\(\sin\theta_e=2\omega_\Delta/L=0.7156\). En una vuelta se obtienen
\[
 \begin{aligned}
 x_e&=0.01602944,\\
 \theta_f&=0.797482699881,\qquad
 \theta_s=2.344109953708.
 \end{aligned}
\]
El Jacobiano tiene
\[
 \operatorname{tr}J_e=-\frac{1+\tau_2L\cos\theta_e/2}{T_f},\qquad
 \det J_e=\frac{L\cos\theta_e}{2T_f}.
\]
El primer punto es un foco estable, con espectro
\(-33.41714171\pm40.52189661\ii\); el segundo es una silla, con espectro
\((-37.78074660,73.01945339)\).

Después de trasladar el equilibrio y conservar la no linealidad sinusoidal,
la transferencia es
\[
 \begin{aligned}
 W(s)&=-\frac L2\frac{1+\tau_2s}{s(1+T_fs)},\\
 \operatorname{Im}W(\ii\omega)
 &=\frac{L(1+\tau_2T_f\omega^2)}{2\omega(1+T_f^2\omega^2)}>0.
 \end{aligned}
\]
No hay cruce real para \(\omega>0\), de modo que el cierre directo con
una función descriptiva estática real no proporciona semilla.

Para \(L>0\), introducir \(y=\dot\theta\) permite despejar
\[
 x=\frac{T_f}{L}(\omega_\Delta-y)-\frac{\tau_2}{2}\sin\theta.
\]
Derivar \(y\), sustituir \(\dot x\) y agrupar
\(\tau_1+\tau_2=T_f\) da
\[
 T_f\dot y=\omega_\Delta-\frac L2\sin\theta
              -\left(1+\frac{\tau_2L}{2}\cos\theta\right)y.
\]
En una rotación positiva, \(y>0\) permite usar \(\theta\) como variable
independiente y dividir entre \(\dot\theta=y\):
\begin{equation}
 \frac{\dd y}{\dd\theta}=H_L(\theta,y)
 =\frac{\omega_\Delta-(L/2)\sin\theta}{T_fy}
    -\frac{1+(\tau_2L/2)\cos\theta}{T_f}.
 \label{bre:res:pll-angle}
\end{equation}
El mapa \(P_L(y_0)=y(2\pi;y_0)\) sólo se define en los datos que conservan
\(y>0\) durante toda la vuelta. Un punto fijo da una órbita corredora, con
periodo físico \(\mathcal T=\int_0^{2\pi}y(\theta)^{-1}\,\dd\theta\).
Si \(z=\partial y/\partial y_0\), la diferenciación del PVI produce
\[
 \begin{gathered}
 z'=-\frac{\omega_\Delta-(L/2)\sin\theta}{T_f y^2}z,\\
 z(0)=1,\qquad P_L'(y_0)=z(2\pi).
 \end{gathered}
\]
Por ello \(|P_L'|<1\) implica atracción transversal local del punto fijo
y \(|P_L'|>1\), repulsión. El criterio se aplica al retorno de la ODE
sobre el cilindro, no a un mapa de estados instantáneos de Caputo.

En \(L=0\), \eqref{bre:res:pll-angle} posee
\(y(\theta)\equiv\omega_\Delta\), con
\[
 \mathcal T_0=\frac{2\pi}{\omega_\Delta},\qquad
 P_0'(\omega_\Delta)=e^{-2\pi/(T_f\omega_\Delta)}.
\]
La reconstrucción que divide entre \(L\) no se usa en ese extremo.
En las variables originales, \(\theta(t)=\omega_\Delta t\), y proponer
\(x_p=A\sin\theta+B\cos\theta\) en la ecuación del filtro conduce a
\(A-\omega_\Delta T_fB=\tau_1/2\) y
\(\omega_\Delta T_fA+B=0\). Resolver estas dos ecuaciones da
\begin{equation}
 x_p(\theta)=\frac{\tau_1}{2[1+(\omega_\Delta T_f)^2]}
              (\sin\theta-\omega_\Delta T_f\cos\theta).
 \label{bre:res:pll-zero-gain}
\end{equation}
Las demás soluciones difieren en \(Ce^{-t/T_f}\). Su contracción durante
una vuelta coincide con \(P_0'\), lo que justifica utilizar esta órbita
exacta para iniciar la continuación en 101 valores hasta \(L=500\).
El retorno usó tolerancias \((2\times10^{-11},2\times10^{-13})\) y paso
angular máximo 0.02.

En la sección \(\theta=0\pmod{2\pi}\), se encontraron los resultados de
la Tabla~\ref{bre:res:tab-pll}. Los errores de retorno cilíndrico fueron
\(2.68\times10^{-15}\) para el ciclo estable y \(6.18\times10^{-14}\)
para la órbita inestable. La reproducción independiente dio periodos
\(0.0826766624969991\) y \(0.108541779641824\), con multiplicadores
\(0.817204143145181\) y \(1.44110733084579\).
Las diferencias angulares se calcularon módulo \(2\pi\).

\begin{table}[t]
 \centering
 \caption{Retornos del PLL en \(L=500\), con fase inicial cero.}
 \label{bre:res:tab-pll}
 \begin{tabular}{@{}lrr@{}}
 \toprule
 Magnitud & Ciclo estable & Órbita inestable\\\midrule
 \(y_0\) & 140.125084913 & 135.172047420\\
 \(x_0\) & 0.004908904250 & 0.005535958797\\
 \(\mathcal T\) & 0.0826766624974 & 0.1085417796382\\
 \(P_L'(y_0)\) & 0.8172041431535 & 1.4411073306967\\
 \bottomrule
 \end{tabular}
\end{table}
\begin{figure*}[t]
 \centering
 \includegraphics[width=0.90\textwidth]{pll_cilindro.pdf}
 \caption{PLL sobre el cilindro de fase. El retorno distingue una órbita
 corredora estable y una órbita periódica inestable que organiza la
 separación entre destinos. Las distancias respetan la periodicidad angular.}
 \label{bre:res:fig-pll}
\end{figure*}

Los sondeos utilizaron los dos equilibrios, cuatro radios
\(10^{-5},10^{-4},10^{-3},10^{-2}\) y doce direcciones por radio.
Las 96 condiciones terminaron en el foco y ninguna contactó el ciclo;
las clasificaciones coincidieron entre los dos cálculos. La órbita
inestable de la Fig.~\ref{bre:res:fig-pll} aporta una descripción de la
frontera que el balance armónico directo no ofrecía. Los valores numéricos
del retorno y los sondeos no sustituyen encierros intervalares ni una
prueba global de separación de cuencas.

\subsection{Sistemas fraccionarios no Chua}
\label{bre:res:nonchua-fo}

\subsubsection{Lorenz generalizado: estabilidad local y absorción}
El campo de \eqref{bre:geo:lorenz} corresponde a
\((\sigma,r,a)=(3.4,6.8,-0.5)\), con \(\sigma=-ar\), y se considera a
\(q=0.995\) \cite{Danca2017}. Para calcular los equilibrios, la tercera
ecuación da \(z=xy\). La segunda queda
\(rx-y-x^2y=0\), y como \(1+x^2>0\),
\[
 y=\frac{rx}{1+x^2},\qquad z=\frac{rx^2}{1+x^2}.
\]
Si \(x=0\), se obtiene el origen. Si \(x\ne0\), sustituir estas
expresiones en la primera ecuación y usar \(\sigma=-ar\) da, con
\(u=x^2>0\),
\[
 0=arx\left[1-\frac r{1+u}-\frac{ru}{(1+u)^2}\right].
\]
Como \(arx\ne0\), se puede dividir y multiplicar por \((1+u)^2\):
\begin{equation}
 \begin{aligned}
 0&=(1+u)^2-r(1+u)-ru\\
  &=u^2+2(1-r)u+(1-r),\\
 u_\pm&=r-1\pm\sqrt{r(r-1)}.
 \end{aligned}
 \label{bre:res:lorenz-equilibria-algebra}
\end{equation}
Para \(r>1\), \(\sqrt{r(r-1)}>r-1\), de modo que \(u_-<0\) se excluye
y \(u_+>0\) genera dos signos de \(x\). Son todos los equilibrios: el
origen y
\[
 E_\pm=\begin{pmatrix}
 \pm3.47564776512854\\\pm1.80689408468029\\6.28012738724303
 \end{pmatrix}.
\]
El Jacobiano es
\begin{equation}
 J_L=\begin{pmatrix}
 -\sigma&\sigma-az&-ay\\r-z&-1&-x\\y&x&-1
 \end{pmatrix}.
 \label{bre:res:lorenz-jacobian}
\end{equation}
En el origen, la tercera coordenada se desacopla, por lo que
\[
 p_0(s)=(s+1)[(s+\sigma)(s+1)-\sigma r].
\]
La sustitución de los equilibrios exactos en el campo lo anula; el residuo
máximo al evaluar sus coordenadas numéricas fue \(1.07\times10^{-14}\).
El origen tiene autovalores
\[
 \begin{aligned}
 \lambda_1&=-7.15580467734555,\\
 \lambda_2&=2.75580467734555,\qquad \lambda_3=-1,
 \end{aligned}
\]
y es inestable. En \(E_\pm\), los autovalores son
\[
 \begin{aligned}
 \lambda_1&=-5.37029951773135,\\
 \lambda_{2,3}&=-0.0148502411343235
                    \pm3.83248917224189\ii.
 \end{aligned}
\]
El margen angular a \(q=0.995\) es \(0.0117287914887680>0\), de modo que
ambos equilibrios externos son localmente estables.
La eliminación de la sección~\ref{bre:geo:section} proporciona además dos
semillas FPP--A, relacionadas por \((x,y,z)\mapsto(-x,-y,z)\).
En las variables \(u=x^2\), \(v=y/x\), una raíz tiene
\((u,v,z)\simeq(0.5564665577,-16.9736823118,-7.4423053062)\), y se
reconstruye mediante \(x=\pm\sqrt u\), \(y=vx\). Su norma de campo es
\(23.42210895339\). Su integración y clasificación de destinos siguen
pendientes; resolver \(J_ff=0\) no determina su cuenca.

Para este campo sí puede demostrarse una cota global, válida para cualquier
reinicio y para \(0<q\leq1\). Considérese
\begin{equation}
 V(x,y,z)=x^2+y^2+\tfrac12(z-20.4)^2.
 \label{bre:res:lorenz-V}
\end{equation}
Su Hessiana \(\operatorname{diag}(2,2,1)\) es definida positiva. En cada
intervalo compacto de existencia, \(X-X_0=I^qF(X)\), con \(F(X)\)
acotado. La desigualdad convexa \eqref{bre:fund:convex}, aplicada a
\(V(X_0+\xi)-V(X_0)\), justifica
\({}^{\mathrm C}D^q V(X)\leq\nabla V(X)\cdot F(X)\); no se utiliza una
regla de cadena fraccionaria con igualdad
\cite{Gomoyunov2018Convex}. El cálculo específico es
\begin{align*}
 \nabla V\cdot F
 & =2x[-3.4(x-y)+0.5yz]\\
 &\quad+2y[6.8x-y-xz]+(z-20.4)(-z+xy)\\
 &=-6.8x^2-2y^2-z^2+20.4z\\
 &=-V+208.08-5.8x^2-y^2-\tfrac12z^2\\
 &\leq -V+208.08.
\end{align*}
Los términos \(xyz\) se cancelan porque sus coeficientes suman
\(1-2+1=0\); los términos \(xy\), porque
\(6.8+13.6-20.4=0\). La constante es
\(20.4^2/2=208.08\).

Escribiendo \(v(t)=V(X(t))\) y
\(r(t)=208.08-v(t)-{}^{\mathrm C}D^qv(t)\geq0\), la fórmula lineal de
variación \eqref{bre:fund:variation} da
\begin{align*}
 v(t)&=v(0)E_q(-t^q)+208.08[1-E_q(-t^q)]\\
 &\quad-\int_0^t(t-s)^{q-1}E_{q,q}(-(t-s)^q)r(s)\,\dd s.
\end{align*}
El núcleo es no negativo. Para \(0<q<1\), esto se sigue de la
representación positiva \cite{Podlubny1999}
\[
 E_q(-t^q)=\frac{\sin(\pi q)}{\pi}
 \int_0^\infty
 \frac{e^{-ut}u^{q-1}}{u^{2q}+2u^q\cos(\pi q)+1}\,\dd u,
\]
pues al derivar para \(t>0\) se obtiene
\(-\frac{\dd}{\dd t}E_q(-t^q)=t^{q-1}E_{q,q}(-t^q)\geq0\).
Integrar esta identidad da
\(\int_0^t u^{q-1}E_{q,q}(-u^q)\,\dd u=1-E_q(-t^q)\), usado para
la fuerza constante. La representación implica también
\(0<E_q(-t^q)\leq1\) y \(E_q(-t^q)\to0\). Para \(q=1\), las mismas
propiedades corresponden a \(e^{-t}\). Por consiguiente,
\begin{equation}
 \begin{aligned}
 V(X(t))&\leq208.08+[V(X_0)-208.08]E_q(-t^q),\\
 V(X(t))&\leq\max\{V(X_0),208.08\}.
 \end{aligned}
 \label{bre:res:lorenz-bound}
\end{equation}
Los subniveles de \(V\) son compactos. La cota impide que una solución
maximal abandone todos los compactos en tiempo finito; la continuación de
Volterra con historia conservada establece existencia global y unicidad.
Para una familia acotada \(B\) de reinicios, sea
\(C_B=\sup_{X_0\in B}V(X_0)<\infty\). Entonces
\[
 \sup_{X_0\in B}V(X(t;X_0))
 \leq208.08+(C_B-208.08)_+E_q(-t^q).
\]
Dado \(\epsilon>0\), elegir un tiempo desde el cual el último término sea
menor que \(\epsilon\) prueba absorción uniforme por
\(\{V\leq208.08+\epsilon\}\). Esta conclusión controla todas las
trayectorias de reinicio; no identifica un atractor oculto ni demuestra
atracción en un espacio funcional admisible aún por construir.

\subsubsection{Rabinovich--Fabrikant: candidatos algebraicos}
El campo \eqref{bre:geo:rf} se evalúa con
\((a,b,q)=(0.1,0.2876,0.998)\) \cite{Danca2017}. La ecuación estacionaria
tercera es \(z(b+xy)=0\). Si \(z=0\), las primeras dos se escriben,
con \(u=x^2\), como
\[
 \begin{pmatrix}a&u-1\\1-u&a\end{pmatrix}
 \begin{pmatrix}x\\y\end{pmatrix}=0.
\]
Su determinante \(a^2+(u-1)^2>0\) exige \(x=y=0\); esta rama sólo da el
origen. Si \(z\ne0\), entonces \(xy=-b\ne0\), por lo que \(x\ne0\) y
\(y=-b/x\). La primera ecuación, multiplicada por \(x\), da
\(-b(z-1+u)+au=0\), de donde
\[
 z=1+\frac{a-b}{b}u.
\]
La segunda, también multiplicada por \(x\), es
\(u(3z+1-u)-ab=0\). Sustituir \(z\) y multiplicar por \(b\) produce
\begin{equation}
 \begin{aligned}
 &(3a-4b)u^2+4bu-ab^2=0,\\
 &u_\pm=\frac{b[-2\pm\sqrt{4+a(3a-4b)}]}{3a-4b}.
 \end{aligned}
 \label{bre:res:rf-equilibria-algebra}
\end{equation}
Para los parámetros elegidos, \(3a-4b=-0.8504\); el discriminante es
positivo y las raíces tienen suma y producto positivos. Ambas son
positivas y distintas. Cada una permite reconstruir dos equilibrios:
\(x=\pm\sqrt u\), \(y=-b/x\),
\(z=1+(a-b)u/b\). La sustitución verifica \(z\ne0\) y las tres ecuaciones.
Se obtienen exactamente cinco equilibrios:
\[
 \begin{gathered}
 E_0=(0,0,0),\\
 E_{1,2}=\begin{pmatrix}
 \mp1.15997695583801\\\pm0.247935959893466\\0.122306917444673
 \end{pmatrix},\\
 E_{3,4}=\begin{pmatrix}
 \mp0.085021329989515\\\pm3.38268055834300\\0.995284804098130
 \end{pmatrix}.
 \end{gathered}
\]
El Jacobiano es
\begin{equation}
 J_R=\begin{pmatrix}
 a+2xy&z-1+x^2&y\\
 3z+1-3x^2&a&3x\\
 -2yz&-2xz&-2(b+xy)
 \end{pmatrix}.
 \label{bre:res:rf-jacobian}
\end{equation}
En el origen, sus bloques dan
\(p_0(s)=(s+2b)[(s-a)^2+1]\) y autovalores
\(-2b,a\pm\ii\). El par complejo tiene orden crítico
\(2\arctan(1/a)/\pi=0.936548965138893\), menor que 0.998.
Así, el origen es inestable. La evaluación en \(E_{3,4}\) también da
inestabilidad, mientras \(E_{1,2}\) son estables, con margen angular
\(0.0435197213462810\). Los equilibrios exactos anulan el campo; el
residuo máximo de sus coordenadas numéricas fue \(1.06\times10^{-14}\). El conteo algebraico exhaustivo de
la sección~\ref{bre:geo:section} produce ocho FPP--A: cuatro en \(z=0\) y
cuatro fuera de ese plano. La realización multicanal permite formular el
balance de \eqref{bre:bal:mimo-balance}; ni sus ecuaciones ni las raíces
FPP--A reemplazan la continuación e integración del campo original.
Quedan pendientes los destinos de los cuatro pares de semillas, sus
controles de acotación y paso, y el sondeo de los cinco equilibrios.

\subsubsection{Muñoz--Pacheco: integración desde FPP--A}
El campo \eqref{bre:geo:munoz} tiene \(a=0.35\) y se evalúa a
\(q=0.97\) \cite{munoz2018}. No tiene equilibrios: \(f_2=0\) exigiría
\(|x|=1\), \(f_3=0\) impondría \(z=-xy\), y la primera componente
quedaría
\[
 f_1=-x(y^2-y+a)=-x\bigl[(y-1/2)^2+0.10\bigr]\ne0.
\] Las cuatro raíces suaves de
\eqref{bre:geo:munoz-roots} son
\[
 \begin{gathered}
 p_{1,2}=\begin{pmatrix}
 \pm0.816784043380077\\0.591607978309962\\\mp0.333568086760154
 \end{pmatrix},\\
 p_{3,4}=\begin{pmatrix}
 \pm3.183215956619923\\-0.591607978309962\\\mp5.066431913239846
 \end{pmatrix}.
 \end{gathered}
\]
Son condiciones iniciales algebraicas, separadas de la superficie no suave
\(x=0\). PECE con historia completa, \(h=0.01\) y \(T=35\), produjo las
cuatro trayectorias de la Fig.~\ref{bre:res:fig-munoz}. Las normas máximas
observadas fueron 6.421 para el par interior y 10.178 para el exterior.
Estos máximos sólo establecen finitud durante la ventana calculada.
La ausencia de equilibrios hace vacua la parte de la definición de ocultedad
que compara sus vecindades; siguen siendo necesarias existencia de un
conjunto invariante, atracción y acotación de la respuesta.

\begin{figure*}[t]
 \centering
 \includegraphics[width=0.96\textwidth]{munoz_semillas.png}
 \caption{Muñoz--Pacheco, \(q=0.97\), desde cuatro semillas FPP--A exactas,
 con \(h=0.01\) y \(T=35\). Las cruces marcan las condiciones iniciales.
 Las trayectorias son resultados preliminares de integración causal.}
 \label{bre:res:fig-munoz}
\end{figure*}

En el diagnóstico FPP--H de \eqref{bre:geo:history}, la comparación de
\(h=0.01\) y \(h=0.005\) sobre \([2,17.5]\) no produjo un mínimo
convergente a cero. El mínimo normalizado fino fue aproximadamente 0.0871
cerca de \(t=11.965\). El valor grueso 0.0254 en \(t=32.33\) queda fuera
del horizonte fino y no constituye una comparación de convergencia.
No se ha identificado un evento FPP--H. La familia direccional de frontera
\((0,y_0,0)\) tampoco ofrece un candidato acotado: sus soluciones exactas
son \(x=z=0\), \(y=y_0+t^q/\Gamma(q+1)\), como se dedujo en
\eqref{bre:geo:munoz-boundary}. La comparación de paso y horizonte de las
cuatro semillas suaves, y la reproducción del caso publicado a \(q=0.996\),
permanecen pendientes.

\subsubsection{Comparaciones de exponentes con otros parámetros}
Los ensayos anteriores no se confunden con dos comparaciones dinámicas de
referencia \cite{fischer2020}. En el Lorenz estándar
\[
 \begin{aligned}
 \Caputo qx&=\sigma(y-x),\\
 \Caputo qy&=\rho x-y-xz,\\
 \Caputo qz&=xy-\beta z,
 \end{aligned}
\]
con \((\sigma,\rho,\beta,q)=(10,200,8/3,0.985)\), el dato \(X_0=(0.1,0.1,0.1)\), \(h=0.001\), \(T=500\) e intervalo de
ortogonalización \(T_R=5\) dieron
\((-0.0027791532,-0.0868071712,-1.6224574503)\). La diferencia máxima
respecto de \((-0.0026,-0.0870,-1.6225)\) fue menor que
\(1.93\times10^{-4}\). La concordancia cuantitativa verifica ese cálculo
finito de referencia; no aporta un candidato caótico ni resultados sobre
el Lorenz generalizado anterior.

Para RF con \((a,b,q)=(0.1,0.98,0.999)\), desde el mismo dato inicial,
\(h=0.005\), \(T=300\) y \(T_R=0.02\), el espectro calculado fue
\((0.0611165030,0.0038981399,-1.8324646836)\), frente a
\((0.0749,0.0018,-2.0850)\). La diferencia del tercer exponente,
\(0.2525353\), excedió la tolerancia 0.05 y fue reproducida por otro
solucionador. El primer exponente positivo motiva estudiar la dinámica,
pero el ensayo no incluyó localización de una semilla ni sondeos de
vecindades. Por tanto, estos parámetros y sus diagnósticos no completan
la evaluación del candidato RF con \(b=0.2876\).

\section{Cuencas muestreadas y transporte de las condiciones iniciales}
\label{bre:bas:section}

Los retratos de fase describen la trayectoria alcanzada desde una semilla;
los sondeos y cortes siguientes muestran cómo cambia el destino al cambiar
el dato inicial. En Caputo, cada punto del muestreo define un PVI nuevo con
terminal inferior cero y la misma convención de reinicio. Los colores
clasifican trayectorias en una ventana finita. La continuación, en cambio,
transporta una condición inicial entre sistemas: su gráfica debe leerse
junto con el parámetro que cambia y el tratamiento de la memoria.

\subsection{Chua no suave: sondeos dentro de bolas}
\label{bre:bas:pwl}

Para el caso de la Sec.~\ref{bre:res:chua-pwl}, con $q=0.9998$, se
muestrearon los tres equilibrios mediante
$X_0=e+rU^{1/3}d$, con $U$ uniforme en $[0,1]$ y $d$ uniforme en la
esfera unidad, independientes. Por tanto, $r$ es el radio de la bola y
la distancia de una sonda concreta puede ser menor que $r$. Se usaron
ABM--PECE con memoria completa, $h=0.01$, horizonte máximo $T=200$ y
descarte de 50. El criterio de contacto fue el percentil 90 de distancias
euclidianas a la nube objetivo, con umbral $0.5$ y hasta 2000 puntos de
cada cola. Los criterios de equilibrio y salida pudieron detener antes
una integración; una salida numérica no demuestra explosión matemática.

El protocolo local reunió 7200 trayectorias en los siete radios desde
$10^{-5}$ hasta $10^{-2}$ y no registró contactos. La extensión a
$r=0.03,0.1,0.3$ añadió 10\,200 trayectorias. Hasta $r=0.1$ el
acumulado fue $0/13\,800$; en las 3600 sondas de $r=0.3$ aparecieron
37 contactos, 19 desde la bola de $E_+$ y 18 desde la de $E_-$.
Sus distancias iniciales efectivas estuvieron entre $0.1941805298$ y
$0.2997254735$. El primer radio ensayado con acceso delimita una escala
de búsqueda; no estima la distancia mínima exacta a la cuenca.

\begin{figure*}[t]
 \centering
 \includegraphics[width=0.74\textwidth]{pwl_sondeos_volumetricos.pdf}
 \caption{Chua no suave de Caputo: nueve historias representativas de un
 banco de 17\,400 sondas volumétricas. Se distinguen el objetivo, los
 equilibrios y las trayectorias con destinos diferentes. Los triángulos
 señalan el primer cruce del encuadre por las dos historias de salida
 representadas; no son estados finales. La figura no dibuja la frontera
 de una cuenca ni representa todas las sondas.}
 \label{bre:bas:fig-pwl-volume}
\end{figure*}

La Fig.~\ref{bre:bas:fig-pwl-volume} permite comparar transitorios que
parten de vecindades semejantes y terminan en clases distintas. Sus ejes
conservan las coordenadas físicas, con límites ajustados por componente;
la caja tridimensional no impone una misma escala numérica a los tres
ejes. Los contactos a radio macroscópico se interpretan separadamente de
la ausencia observada dentro del protocolo local.

\begin{figure*}[t]
 \centering
 \includegraphics[width=0.98\textwidth]{pwl_nube_volumetrica_conteos.pdf}
 \caption{Condiciones iniciales y destinos del sondeo volumétrico PWL.
 Arriba: las 1200 condiciones efectivamente integradas en la bola de
 radio $0.3$ de cada equilibrio, trasladadas a coordenadas locales
 $X_0-e$, con la misma escala física en los tres ejes. El signo $+$
 marca el centro; no se normalizan ni se superponen bolas de radios
 diferentes. Abajo: conteos apilados de los diez radios, 5800
 trayectorias por equilibrio. Los colores indican la clasificación de
 horizonte finito obtenida mediante ABM--PECE con $h=0.01$, memoria
 completa, $T\leq200$ y descarte de 50. El amarillo muestra los
 37 contactos, todos en las bolas de radio $0.3$ de $E_\pm$.
 La densidad visual del panel superior no representa una densidad
 invariante del flujo ni una probabilidad certificada de atracción.}
 \label{bre:bas:fig-pwl-clouds}
\end{figure*}

La Fig.~\ref{bre:bas:fig-pwl-clouds} representa las condiciones iniciales
para el radio mayor ensayado. Los equilibrios de
\eqref{bre:res:chua-equilibrium-reduction} usados en la traslación son
\[
 E_0=0,\qquad E_\pm\simeq\pm\begin{pmatrix}
 6.5883078865\\0.0028364023\\-6.5854714843
 \end{pmatrix}.
\]
Todas las condiciones iniciales satisfacen $\norm{X_0-e}/r\leq1$.
Los conteos suman 5827 destinos de equilibrio, 5826 de otra clase,
5710 salidas y 37 contactos.

\subsection{Chua arctangente: superficies esféricas y cortes planos}
\label{bre:bas:arctan}

En c590 se mantienen los parámetros de
\eqref{bre:res:c590-parameters} y $q=0.9999$. Las 13\,500 sondas
pertenecen a superficies $X_0=e+rd$, mediante direcciones axiales y una
distribución de Fibonacci sobre la esfera. Se usaron nueve radios,
$10^{-5},10^{-4},10^{-3},10^{-2},0.03,0.1,0.3,1,2$, con
100, 200, $\ldots$, 900 direcciones por equilibrio, respectivamente.
No se trata de un muestreo uniforme del volumen de las bolas.
Las integraciones utilizaron ABM con memoria completa, $h=0.005$,
$T=180$ y descarte de 90, conservando las reglas de parada del protocolo.

Hasta $r=0.3$ hubo cero contactos entre 8400 trayectorias. En $r=1$
hubo 22 de 2400, todos desde $E_0$; en $r=2$, 588 de 2700, repartidos
en 289, 150 y 149 desde $E_0,E_+,E_-$. El total fue 610 contactos.
Los otros destinos fueron 4 equilibrios, 131 salidas y 12\,755
trayectorias de otra clase. Esta última categoría no determina por sí
sola un atractor diferente: puede contener transitorios todavía no
resueltos por la ventana y el clasificador.

\begin{figure*}[t]
 \centering
 \includegraphics[width=0.65\textwidth]{arctan_sondeos_superficiales.pdf}
 \caption{c590: nueve historias representativas seleccionadas de los
 sondeos sobre superficies esféricas hasta $r=1$. El conteo completo de
 13\,500 pruebas incorpora además el bloque $r=2$, que no se representa
 mediante historias adicionales en esta figura. Las circunferencias
 aparentes son proyecciones de trayectorias; no delimitan bolas de muestreo
 ni prueban periodicidad exacta del PVI de Caputo.}
 \label{bre:bas:fig-arctan-sphere}
\end{figure*}

Un análisis complementario fijó $z(0)$ y clasificó datos
$(x(0),y(0),z(0))$ en cinco planos. Los cortes se eligieron a partir de
las coordenadas de los equilibrios y los percentiles 10, 50 y 90 de la
nube postransitoria simétrica $A_+\cup A_-$:
\[
 z(0)\in\{-12.56295,-4.64761,0,4.64761,12.56295\}.
\]
El rectángulo común fue
$x(0)\in[-11.238014,11.238014]$,
$y(0)\in[-4.311997,4.311997]$. Se mantuvieron $h=0.005$, $T=180$
y descarte de 90. El contacto con una nube de referencia se evaluó
mediante el percentil 90 y umbral $0.5901016489$; la tolerancia de
clasificación a equilibrio fue $0.5$. Estas tolerancias describen
vecindad numérica en coordenadas físicas, no igualdad de conjuntos límite.

\begin{figure*}[t]
 \centering
 \includegraphics[width=\textwidth]{c590_cuencas_multicorte.pdf}
 \caption{Clasificación por destinos de c590 sobre cinco cortes $xy$.
 Las estrellas identifican equilibrios contenidos en el plano; los
 círculos sólo proyectan equilibrios de otros planos, con su separación
 $\Delta z$. La retícula visual es $241\times241$ por corte y contiene
 relleno categórico: los 290\,405 nodos mostrados corresponden a
 62\,416 trayectorias efectivamente integradas, no a una integración
 independiente por píxel. El gris incluye relleno y omisiones heredadas;
 la etiqueta gráfica de divergencia sólo se interpreta como resultado
 dinámico en los nodos cuya integración fue realizada.}
 \label{bre:bas:fig-c590-slices}
\end{figure*}

La retícula se refinó de forma adaptativa: se conservaron 35\,941
integraciones de la malla precedente y se integraron 26\,475 centros
de celdas cuyos cuatro vértices tenían clases distintas. Los nodos
restantes recibieron relleno categórico sin suavizado. Una máscara
separa los nodos calculados de ese relleno y de las omisiones grises
heredadas; por ello las áreas coloreadas no proporcionan directamente
fracciones de volumen de las cuencas. Entre los nodos realmente
integrados hubo 17\,369 contactos con $A_+$, 17\,514 con $A_-$,
843 destinos de equilibrio, 10\,331 de otra clase acotada y 16\,359
salidas, sin fallos numéricos registrados. Los cortes muestran regiones
de acceso y zonas donde conviene refinar; no establecen conectividad
global, dimensión de frontera ni exclusión de una vecindad tridimensional.

\subsection{Chua entero: destinos cerca de los equilibrios}
\label{bre:bas:integer}

Para el Chua entero de la Sec.~\ref{bre:res:chua-integer}, el control
comparado tomó las mismas 504 condiciones en las dos implementaciones:
tres equilibrios, siete radios entre $10^{-5}$ y $10^{-2}$ y 24 puntos
uniformes dentro de cada bola. Con $T=500$, descarte de 120,
tolerancia de equilibrio $10^{-3}$ y tolerancia normalizada a la nube
$0.08$, coincidieron las 504 clases: 152 trayectorias hacia $E_0$,
198 acotadas no triviales y 154 salidas. Ninguna hizo contacto con el
candidato. La concordancia reproduce el resultado del muestreo común;
no amplía sus radios ni convierte sus etiquetas finitas en una prueba
asintótica.

\begin{figure*}[t]
 \centering
 \includegraphics[width=0.72\textwidth]{chua_entero_sondeos.png}
 \caption{Trayectorias del control de vecindades del Chua entero, con la
 nube objetivo en verde. Los ejes $x0,x1,x2$ designan las tres componentes
 del estado $(x,y,z)$, no tiempos iniciales. La escala amplia conserva
 las trayectorias que abandonan la región del candidato; el retrato
 detallado del objetivo aparece en la Fig.~\ref{bre:res:fig-chua-integer}.}
 \label{bre:bas:fig-integer-probes}
\end{figure*}

\subsection{MAVPD: selección por continuación y sondeo del candidato}
\label{bre:bas:mavpd}

La continuación descrita en la Sec.~\ref{bre:res:mavpd} transportó la
rama de menor frecuencia desde $(\xi,\gamma)=(3.1,0.1)$ hasta
$\xi=2.85$, y después modificó $\gamma$. La
Fig.~\ref{bre:bas:fig-mavpd} compara el exponente máximo finito con
$-\max\operatorname{Re}\lambda(DF(E_\pm))$, positivo cuando los
equilibrios laterales son linealmente estables. El exponente de una
trayectoria y el margen local de un equilibrio miden objetos distintos.
Los siete desplazamientos respecto de
$\gamma_H=0.14380379839949112$ no definen una curva de bifurcación
continua certificada: los segmentos sólo enlazan valores ensayados.

En los desplazamientos $0.002,0.003,0.005$, el sondeo preliminar de
$E_0$ dio 6 contactos de 12. En $0.008$ y $0.010$ hubo cero contactos
y exponentes máximos mayores que $0.5$. La regla prefijada seleccionó
el mayor de estos desplazamientos, $\gamma=0.15380379839949113$;
un exponente positivo aislado no fue el único criterio.
En cada nivel del transporte se descartaron 60 unidades y se retuvieron
10, con DOP853, tolerancias $2\times10^{-10}$ y $2\times10^{-12}$,
y paso máximo $0.02$.

\begin{figure*}[t]
 \centering
 \includegraphics[width=0.49\textwidth]{mavpd_continuacion_cribado.pdf}\hfill
 \includegraphics[width=0.49\textwidth]{mavpd_sondeos_destinos.pdf}
 \caption{MAVPD. Izquierda: cribado de los niveles de continuación en
 $\gamma$; la tasa azul es un exponente de tiempo finito y la naranja
 es el margen lineal de $E_\pm$. Derecha: destinos de las 112 sondas
 del candidato seleccionado. Las barras 40, 36 y 36 cuentan condiciones
 lanzadas alrededor de $E_0,E_+,E_-$ que terminaron en algún equilibrio;
 no indican que todas hayan vuelto al equilibrio del que partieron.
 La barra de contactos con el objetivo es nula en los tres grupos.}
 \label{bre:bas:fig-mavpd}
\end{figure*}

El sondeo final utilizó 108 direcciones esféricas deterministas: tres
equilibrios, radios $10^{-5},10^{-3},10^{-2}$ y doce direcciones por
radio. Se añadieron cuatro condiciones en ambos sentidos del autovector
inestable de $E_0$, a distancias $10^{-7}$ y $10^{-4}$. Se descartaron
400 unidades y se conservaron 100; DOP853 usó tolerancias
$2\times10^{-10}$ y $2\times10^{-12}$ y paso máximo $0.03$.
Todas las sondas terminaron en un equilibrio, sin contactos con el
candidato, ambigüedades ni fallos. Es una exploración localizada de
destinos cerca de los equilibrios, no un mapa volumétrico completo.

\subsection{PLL: continuación de la ganancia y vecindades cilíndricas}
\label{bre:bas:pll}

Con $(\tau_1,\tau_2,\omega_\Delta)=(0.0448,0.0185,178.9)$, la
continuación aumenta la ganancia $L$ desde cero hasta 500 en pasos de 5.
La Fig.~\ref{bre:bas:fig-pll} representa la velocidad angular y la
coordenada del filtro en la sección $\theta=0$. Cada punto representa
un ciclo obtenido del problema de retorno del nivel correspondiente;
no es una muestra consecutiva de una trayectoria a parámetros fijos.
En $L=0$, la coordenada $x$ procede de la solución explícita de la
Sec.~\ref{bre:res:pll}, pues la inversión basada en dividir por $L$
sólo es válida cuando $L>0$.

\begin{figure*}[t]
 \centering
 \includegraphics[width=0.94\textwidth]{pll_continuacion_ganancia.pdf}\\[4pt]
 \includegraphics[width=0.66\textwidth]{pll_sondeos_cilindricos.pdf}
 \caption{PLL. Arriba: continuación de las coordenadas de retorno hasta
 $L=500$. Abajo: las 96 condiciones iniciales de los sondeos, expresadas
 mediante $u=x-x_e$ y la fase desplazada y envuelta $v$. Los dos grupos
 corresponden al foco y a la silla. La figura inferior representa datos
 iniciales, no una frontera de cuenca; sus 96 trayectorias terminaron en
 el foco bajo el protocolo indicado.}
 \label{bre:bas:fig-pll}
\end{figure*}

En el nivel final se usó la distancia cilíndrica escalada con escalas
$s_u=\tau_1/2=0.0224$ y $s_v=\pi$. Alrededor del foco y de la silla
se tomaron doce puntos uniformes dentro de cada bola escalada, para
$r=10^{-5},10^{-4},10^{-3},10^{-2}$: 96 condiciones en total.
El horizonte fue 5 y la cola evaluada tuvo longitud 1, con paso máximo
$0.002$, tolerancia al foco $10^{-5}$ y umbral de contacto con la nube
$0.01$. Las 48 sondas de cada equilibrio terminaron en el foco.
La coexistencia del ciclo corredor con el foco y la separatriz calculada
explica la utilidad de continuar un ciclo lejano; la exclusión matemática
de vecindades completas exigiría validar la separación de cuencas.

\subsection{Kalman--Fitts: del retorno por conmutación al ciclo suave}
\label{bre:bas:kalman}

La Fig.~\ref{bre:bas:fig-kalman} completa la ruta de la
Sec.~\ref{bre:res:kalman}: a partir del retorno del sistema con signo se
transporta el estado por los 51 niveles de
$(1-\lambda)\operatorname{sign}(c^{\T}X)+
\lambda\tanh(c^{\T}X/0.01)$. Los puntos superiores son estados al
finalizar cada etapa, de duración 120 y paso $0.01$; su variación refleja
tanto el cambio de sistema como la fase alcanzada. No son amplitudes
extremas, puntos fijos ni un diagrama de bifurcación.

\begin{figure*}[t]
 \centering
 \includegraphics[width=0.78\textwidth]{kalman_continuacion_signo_tanh.pdf}\\[4pt]
 \includegraphics[width=0.96\textwidth]{kalman_ciclo_proyecciones.pdf}
 \caption{Kalman--Fitts. Arriba: los cuatro estados terminales de la
 continuación desde signo hasta tangente hiperbólica. Abajo: proyecciones
 del régimen final del sistema suave, con $(x,y,z)=(x_1,x_2,x_3)$;
 la cuarta coordenada permanece en la integración. La geometría se
 contrasta con el periodo numérico $11.9169960234$ y los controles
 de retorno e integración de la Sec.~\ref{bre:res:kalman}.}
 \label{bre:bas:fig-kalman}
\end{figure*}

El control local alrededor del único equilibrio utilizó 48 datos
uniformes dentro de bolas de cuatro dimensiones: doce para cada radio
$10^{-5},10^{-4},10^{-3},10^{-2}$. Con horizonte 600, descarte 150 y
paso $0.02$, ninguno contactó con el ciclo. Cuarenta recibieron etiqueta
de equilibrio y ocho quedaron acotados sin contacto bajo la tolerancia
utilizada. Estos ocho resultados no se reinterpretan como convergencias
asintóticas ya resueltas. La trayectoria final, los retornos y el sondeo
local aportan controles complementarios a la continuación.

\FloatBarrier
\section{Aplicación de la metodología y alcance de las conclusiones}
\label{bre:closing}

\subsection{Condiciones que deben acompañar a cada resultado}

El procedimiento comienza con una ecuación bien definida. Para cada caso
se especifican el campo, sus parámetros, las coordenadas y escalas, el
dominio y las superficies donde cambia su regularidad. En orden
fraccionario se añaden el operador, el orden de cada componente, el
terminal inferior y la inicialización. El argumento de existencia local
se aplica en el dominio utilizado; la existencia global exige una cota
adicional o se limita la conclusión al intervalo anterior a la salida.
Un cambio de unidades temporales modifica el segundo miembro mediante
la potencia $q$ de la escala, como se deduce del cambio de variable en
la integral de Caputo.

La construcción algebraica de una semilla requiere sustituirla en las
ecuaciones originales. En las eliminaciones se examinan por separado
los factores divididos y los denominadores nulos; se descartan los
equilibrios cuando se buscan PP. En un balance armónico se conservan el
signo de realimentación, el término medio y la fase de reconstrucción.
El residuo se calcula en las coordenadas físicas, incluyendo el efecto
de los armónicos descartados. Una raíz aproximada de las ecuaciones
proyectadas no establece la existencia de una solución de la ecuación
funcional completa.

La continuación debe llegar al campo objetivo. En una EDO, el estado final
de una etapa basta para comenzar la siguiente. En Caputo se conserva la
contribución de todas las etapas anteriores y se tratan las discontinuidades
temporales de la familia. Si después se utiliza el estado final como dato
de un PVI autónomo con memoria nula, se declara ese reinicio y se integra
su nueva trayectoria. El balance estacionario, la trayectoria durante la
continuación y la trayectoria autónoma de reinicio corresponden a tres
problemas relacionados, con condiciones iniciales distintas.

La validación numérica compara aproximaciones del mismo problema.
Refinar paso, ampliar horizonte y variar una perturbación inicial
responden a preguntas distintas. La comparación entre códigos tampoco
sustituye una cota de error si ambos comparten una inicialización
inconsistente. El residuo integral y los ejemplos exactos de control
permiten examinar el operador efectivamente aproximado. Para soluciones
sensibles, la separación puntual de trayectorias a tiempos largos se
complementa con retornos y observables definidos sobre ventanas
equivalentes; no se exige una coincidencia punto por punto incompatible
con la sensibilidad dinámica.

La identificación de un destino atractivo exige algo más que acotación
observada. Una región positivamente invariante puede contener varios
destinos. Un exponente variacional estimado positivo, una señal de banda
ancha o un estadístico cercano al régimen caótico aportan diagnósticos
complementarios, sin constituir por separado una demostración de caos.
La ausencia de contactos en un conjunto finito de condiciones iniciales
tampoco implica la exclusión de una vecindad abierta completa.

\subsection{Cobertura de las cuatro familias}

La Tabla~\ref{bre:closing:coverage} organiza el alcance de los cálculos.
La distinción entre tipo Chua y no Chua corresponde a la estructura del
campo; la distinción entre orden entero y fraccionario corresponde al
operador temporal. Una representación de Lur'e en un sistema no Chua no
cambia esa clasificación.

\begin{table*}[t]
\centering
\caption{Resultados disponibles y operaciones adicionales para una conclusión más fuerte.}
\label{bre:closing:coverage}
\small
\begin{tabularx}{\textwidth}{@{}L{0.14\textwidth}L{0.20\textwidth}Y Y@{}}
\toprule
Familia & Sistemas & Resultado establecido o calculado & Desarrollo adicional\\
\midrule
Entero tipo Chua & Chua PWL & Equilibrios, estabilidad, balance y semilla;
trayectorias, comparación entre implementaciones e indicadores finitos. &
Control de atracción y separación de vecindades completas para una prueba de ocultedad.\\
Entero no Chua & MAVPD, Kalman--Fitts, PLL &
Balance o demostración de su obstrucción; rutas alternativas, retornos e
indicadores. Resolución de PP y ausencia de PP según el sistema. &
Validación de los operadores de retorno y de sus dominios; regiones
invariantes que separen cuencas. En MAVPD, prueba de caos adicional a los exponentes estimados.\\
Fraccionario tipo Chua & PWL y arctangente &
Formulación causal, estabilidad sectorial, semillas y trayectorias con
memoria. Sondeos de reinicio; controles de paso donde se indican. &
Completar el refinamiento del caso PWL; controlar el espacio de perfiles,
la atracción y la exclusión de vecindades para fortalecer la clasificación.\\
Fraccionario no Chua & Lorenz generalizado, RF, Muñoz--Pacheco &
Equilibrios y PP/FPP--A; acotación y absorción de Lorenz. Integraciones
piloto de semillas Muñoz y diagnósticos específicos RF. &
Completar localización autónoma con memoria y controles por caso;
establecer atracción y cuencas. En RF, no trasladar resultados entre
valores distintos de $b$.\\
\bottomrule
\end{tabularx}
\end{table*}

Para Lorenz generalizado, la cota absorbente permite continuar todas las
soluciones de reinicio y restringe la región de búsqueda. No identifica
el conjunto límite al que llega cada semilla. El siguiente cálculo
consiste en integrar las semillas algebraicas bajo un mismo PVI,
comparar los destinos al refinar la discretización y analizar su
atracción. Para RF, cada conjunto de parámetros requiere sus propios
equilibrios, márgenes sectoriales, condiciones iniciales y controles de
existencia durante el horizonte usado. Para Muñoz--Pacheco, la ausencia
de equilibrios elimina una condición de exclusión local, pero no
demuestra acotación ni atracción: la trayectoria exacta no acotada sobre
$x=z=0$ impide inferir acotación global a partir de los pilotos.

El análisis topológico puede reducir esta parte del trabajo cuando se
construye un objeto al que aplicar un teorema. Un grado no nulo permite
probar un cero del mapa correspondiente; un retorno contractivo en un
dominio invariante permite probar un punto fijo atractivo; una
separatriz demostrada puede excluir el acceso desde una región. El
objeto, su dominio, la regularidad y la condición de frontera deben
verificarse antes de aplicar el resultado. En orden fraccionario, la
proyección puntual de una trayectoria no proporciona automáticamente
un mapa de retorno: es necesario conservar la memoria o justificar una
reducción funcional compatible.

Los desarrollos anteriores permiten ejecutar la búsqueda y reconocer el
alcance preciso de cada resultado. Las cuestiones abiertas se expresan
como operaciones matemáticas o numéricas concretas: construir un
dominio invariante, validar un retorno, controlar el defecto integral,
identificar una familia admisible de historias o excluir vecindades
completas. De este modo, la metodología distingue los pasos resueltos
de las hipótesis que todavía requieren demostración o cálculo.

\clearpage
\bibliographystyle{ieeetr}
\bibliography{referencias}
\end{document}
